%=====================================================================
%  BUET -- MATH 243 / MATH 245 / MATH 241
%  Probability and Statistics: Topic-wise Question Bank WITH FULL SOLUTIONS
%  Sessions: 2011-2012 to 2023-2024 (12 Exam Sessions)
%  Compile with: pdflatex or lualatex (run twice for Table of Contents)
%
%  COVER CREDITS: name is set below. For the logos, place the official
%  image files next to this .tex file (Overleaf: upload to the project):
%      gemini-logo.png  (or .pdf / .jpg)
%      claude-logo.png  (or .pdf / .jpg)
%  If a file is missing, the cover simply shows the text without that logo.
%=====================================================================
\newcommand{\preparedby}{MAHDI HASAN}
\documentclass[11pt,a4paper]{article}
\usepackage[a4paper,top=24mm,bottom=22mm,left=19mm,right=19mm,headheight=15pt,headsep=8mm]{geometry}
\usepackage[T1]{fontenc}
\usepackage{mathpazo}
\usepackage[scaled=0.90]{helvet}
\usepackage{microtype}
\usepackage{amsmath,amssymb,mathtools}
\usepackage{xcolor,colortbl}
\usepackage{eso-pic}
\usepackage{l3draw}
\usepackage{booktabs,array,tabularx,longtable}
\usepackage{fancyhdr}
\usepackage{multicol}
\usepackage{enumerate}
\usepackage{graphicx}
\usepackage[hidelinks]{hyperref}

%------------------------------------------------------------ palette
\definecolor{ink}{HTML}{1B263B}
\definecolor{inksoft}{HTML}{415A77}
\definecolor{coral}{HTML}{E76F51}
\definecolor{teal}{HTML}{2A9D8F}
\definecolor{sand}{HTML}{FBF7F0}
\definecolor{mist}{HTML}{E8F4F2}
\definecolor{slate}{HTML}{6B7280}
\definecolor{gold}{HTML}{E9C46A}
\definecolor{goldbg}{HTML}{FDF6E3}
\definecolor{rose}{HTML}{FCE9E4}

\ExplSyntaxOn
% l3draw renamed \draw_linewidth:n -> \draw_set_linewidth:n; keep older TeX Live working
\cs_if_exist:NF \draw_set_linewidth:n { \cs_set_eq:NN \draw_set_linewidth:n \draw_linewidth:n }
\color_set:nnn {ink} {HTML} {1B263B}
\color_set:nnn {inksoft} {HTML} {415A77}
\color_set:nnn {coral} {HTML} {E76F51}
\color_set:nnn {teal} {HTML} {2A9D8F}
\color_set:nnn {sand} {HTML} {FBF7F0}
\color_set:nnn {mist} {HTML} {E8F4F2}
\color_set:nnn {slate} {HTML} {6B7280}
\color_set:nnn {gold} {HTML} {E9C46A}
\color_set:nnn {goldbg} {HTML} {FDF6E3}
\color_set:nnn {rose} {HTML} {FCE9E4}

\NewDocumentCommand{\drawlabel}{m m m m m}{
  \hcoffin_set:Nn \l_tmpa_coffin {#5}
  \draw_scope_begin:
  \draw_transform_shift:n {#1, #2}
  \draw_coffin_use:Nnn \l_tmpa_coffin {#3} {#4}
  \draw_scope_end:
}

\NewDocumentCommand{\coverplot}{}{
  \draw_begin:
  \draw_set_linewidth:n {0.6pt}
  \color_select:n {white!65}
  \draw_path_moveto:n {-0.3cm, 0cm}
  \draw_path_lineto:n {7.5cm, 0cm}
  \draw_path_use_clear:n {stroke}
  \draw_path_moveto:n {0cm, -0.4cm}
  \draw_path_lineto:n {0cm, 2.7cm}
  \draw_path_use_clear:n {stroke}
  \draw_set_linewidth:n {1.2pt}
  \color_select:n {white!80}
  \clist_map_inline:nn {
    {0.6, 0.22}, {1.2, 0.55}, {1.8, 1.12}, {2.4, 1.78}, {3.0, 2.25},
    {3.6, 2.45}, {4.2, 2.25}, {4.8, 1.78}, {5.4, 1.12}, {6.0, 0.55}, {6.6, 0.22}
  }{
    \draw_path_moveto:n {\clist_item:nn {##1}{1} cm, 0cm}
    \draw_path_lineto:n {\clist_item:nn {##1}{1} cm, \clist_item:nn {##1}{2} cm}
    \draw_path_use_clear:n {stroke}
    \color_select:n {coral}
    \draw_path_circle:nn {\clist_item:nn {##1}{1} cm, \clist_item:nn {##1}{2} cm} {1.8pt}
    \draw_path_use_clear:n {fill}
    \color_select:n {white!80}
  }
  \draw_set_linewidth:n {1.9pt}
  \color_select:n {gold}
  \draw_path_moveto:n {0.2cm, 0.06cm}
  \draw_path_curveto:nnn {1.5cm, 0.15cm} {2.4cm, 2.48cm} {3.6cm, 2.48cm}
  \draw_path_curveto:nnn {4.8cm, 2.48cm} {5.7cm, 0.15cm} {7.0cm, 0.06cm}
  \draw_path_use_clear:n {stroke}
  \draw_end:
}

% Circuit diagram for 2016-17 Q7(a)
\NewDocumentCommand{\circuitfig}{}{
  \draw_begin:
  \draw_set_linewidth:n {0.9pt}
  \color_select:n {ink}
  \draw_path_moveto:n {0cm, 0cm} \draw_path_lineto:n {0.9cm, 0cm}
  \draw_path_moveto:n {0.9cm, -1.6cm} \draw_path_lineto:n {0.9cm, 1.6cm}
  \draw_path_moveto:n {0.9cm, 1.6cm} \draw_path_lineto:n {7.1cm, 1.6cm}
  \draw_path_moveto:n {0.9cm, -1.6cm} \draw_path_lineto:n {7.1cm, -1.6cm}
  \draw_path_moveto:n {0.9cm, 0cm} \draw_path_lineto:n {1.7cm, 0cm}
  \draw_path_moveto:n {1.7cm, -0.6cm} \draw_path_lineto:n {1.7cm, 0.6cm}
  \draw_path_moveto:n {1.7cm, 0.6cm} \draw_path_lineto:n {6.3cm, 0.6cm}
  \draw_path_moveto:n {1.7cm, -0.6cm} \draw_path_lineto:n {6.3cm, -0.6cm}
  \draw_path_moveto:n {6.3cm, -0.6cm} \draw_path_lineto:n {6.3cm, 0.6cm}
  \draw_path_moveto:n {6.3cm, 0cm} \draw_path_lineto:n {7.1cm, 0cm}
  \draw_path_moveto:n {7.1cm, -1.6cm} \draw_path_lineto:n {7.1cm, 1.6cm}
  \draw_path_moveto:n {7.1cm, 0cm} \draw_path_lineto:n {8.0cm, 0cm}
  \draw_path_use_clear:n {stroke}
  \clist_map_inline:nn {
    {3.35cm, 1.25cm, 4.0cm, 1.6cm, 0.02},
    {3.35cm, -1.95cm, 4.0cm, -1.6cm, 0.02},
    {2.25cm, 0.25cm, 2.9cm, 0.6cm, 0.01},
    {4.45cm, 0.25cm, 5.1cm, 0.6cm, 0.01},
    {2.25cm, -0.95cm, 2.9cm, -0.6cm, 0.01},
    {4.45cm, -0.95cm, 5.1cm, -0.6cm, 0.01}
  }{
    \color_select:n {mist}
    \draw_path_rectangle:nn {\clist_item:nn {##1}{1}, \clist_item:nn {##1}{2}} {1.3cm, 0.7cm}
    \draw_path_use_clear:n {fill}
    \color_select:n {ink}
    \draw_path_rectangle:nn {\clist_item:nn {##1}{1}, \clist_item:nn {##1}{2}} {1.3cm, 0.7cm}
    \draw_path_use_clear:n {stroke}
    \drawlabel{\clist_item:nn {##1}{3}}{\clist_item:nn {##1}{4}}{hc}{vc}{\sffamily\bfseries\footnotesize\color{ink} \clist_item:nn {##1}{5}}
  }
  \draw_end:
}

% U-shaped component for 2016-17 Q7(b)
\NewDocumentCommand{\ushapefig}{}{
  \draw_begin:
  \draw_set_linewidth:n {0.9pt}
  \color_select:n {inksoft!25}
  \draw_path_moveto:n {0cm, 0cm}
  \draw_path_lineto:n {5.0cm, 0cm}
  \draw_path_lineto:n {5.0cm, 2.2cm}
  \draw_path_lineto:n {4.0cm, 2.2cm}
  \draw_path_lineto:n {4.0cm, 0.75cm}
  \draw_path_lineto:n {1.0cm, 0.75cm}
  \draw_path_lineto:n {1.0cm, 2.2cm}
  \draw_path_lineto:n {0cm, 2.2cm}
  \draw_path_close:
  \draw_path_use_clear:n {fill}
  \color_select:n {ink}
  \draw_path_moveto:n {0cm, 0cm}
  \draw_path_lineto:n {5.0cm, 0cm}
  \draw_path_lineto:n {5.0cm, 2.2cm}
  \draw_path_lineto:n {4.0cm, 2.2cm}
  \draw_path_lineto:n {4.0cm, 0.75cm}
  \draw_path_lineto:n {1.0cm, 0.75cm}
  \draw_path_lineto:n {1.0cm, 2.2cm}
  \draw_path_lineto:n {0cm, 2.2cm}
  \draw_path_close:
  \draw_path_use_clear:n {stroke}
  \draw_set_linewidth:n {0.6pt}
  \color_select:n {coral}
  \draw_path_moveto:n {0cm, -0.35cm} \draw_path_lineto:n {5.0cm, -0.35cm}
  \draw_path_moveto:n {0cm, -0.5cm} \draw_path_lineto:n {0cm, -0.15cm}
  \draw_path_moveto:n {5.0cm, -0.5cm} \draw_path_lineto:n {5.0cm, -0.15cm}
  \draw_path_moveto:n {0cm, 2.55cm} \draw_path_lineto:n {1.0cm, 2.55cm}
  \draw_path_moveto:n {4.0cm, 2.55cm} \draw_path_lineto:n {5.0cm, 2.55cm}
  \draw_path_moveto:n {1.0cm, 1.75cm} \draw_path_lineto:n {4.0cm, 1.75cm}
  \draw_path_use_clear:n {stroke}
  \drawlabel{2.5cm}{-0.35cm}{hc}{vc}{\colorbox{sand}{\sffamily\bfseries\footnotesize\color{ink} $A$}}
  \drawlabel{0.5cm}{2.55cm}{hc}{vc}{\colorbox{sand}{\sffamily\bfseries\footnotesize\color{ink} $B$}}
  \drawlabel{4.5cm}{2.55cm}{hc}{vc}{\colorbox{sand}{\sffamily\bfseries\footnotesize\color{ink} $C$}}
  \drawlabel{2.5cm}{1.75cm}{hc}{vc}{\colorbox{white}{\sffamily\bfseries\footnotesize\color{teal} $D$}}
  \draw_end:
}

% Joint PMF stem plot for 2023-24 Q2(b)(i)
\NewDocumentCommand{\jointpmffig}{}{
  \draw_begin:
  \draw_set_linewidth:n {0.7pt}
  \color_select:n {ink}
  \draw_path_moveto:n {0cm, 0cm} \draw_path_lineto:n {5.9cm, 0cm}
  \draw_path_moveto:n {0cm, 0cm} \draw_path_lineto:n {0cm, 3.6cm}
  \draw_path_moveto:n {0cm, 0cm} \draw_path_lineto:n {1.68cm, 1.42cm}
  \draw_path_use_clear:n {stroke}
  \drawlabel{6.1cm}{0cm}{l}{vc}{\small $x$}
  \drawlabel{1.82cm}{1.56cm}{l}{vc}{\small $y$}
  \drawlabel{0cm}{3.8cm}{hc}{b}{\small $f(x,y)$}
  \clist_map_inline:nn {
    {0cm, 0cm, 0.96cm},
    {2.5cm, 0cm, 2.89cm},
    {5.0cm, 0cm, 0.96cm},
    {0.65cm, 0.55cm, 2.48cm},
    {3.15cm, 0.55cm, 2.48cm},
    {1.30cm, 1.10cm, 1.48cm}
  }{
    \color_select:n {slate}
    \draw_path_circle:nn {\clist_item:nn {##1}{1}, \clist_item:nn {##1}{2}} {1.4pt}
    \draw_path_use_clear:n {fill}
    \draw_set_linewidth:n {1.6pt}
    \color_select:n {coral}
    \draw_path_moveto:n {\clist_item:nn {##1}{1}, \clist_item:nn {##1}{2}}
    \draw_path_lineto:n {\clist_item:nn {##1}{1}, \clist_item:nn {##1}{3}}
    \draw_path_use_clear:n {stroke}
    \color_select:n {teal}
    \draw_path_circle:nn {\clist_item:nn {##1}{1}, \clist_item:nn {##1}{3}} {2.2pt}
    \draw_path_use_clear:n {fill}
  }
  \drawlabel{-0.35cm}{0.98cm}{hc}{b}{\scriptsize $3/28$}
  \drawlabel{2.5cm}{3.08cm}{hc}{b}{\scriptsize $9/28$}
  \drawlabel{5.0cm}{1.15cm}{hc}{b}{\scriptsize $3/28$}
  \drawlabel{0.65cm}{2.67cm}{hc}{b}{\scriptsize $3/14$}
  \drawlabel{3.15cm}{2.67cm}{hc}{b}{\scriptsize $3/14$}
  \drawlabel{1.30cm}{1.67cm}{hc}{b}{\scriptsize $1/28$}
  \drawlabel{0cm}{-0.25cm}{hc}{t}{\tiny\sffamily\color{slate} (0,0)}
  \drawlabel{2.5cm}{-0.25cm}{hc}{t}{\tiny\sffamily\color{slate} (1,0)}
  \drawlabel{5.0cm}{-0.25cm}{hc}{t}{\tiny\sffamily\color{slate} (2,0)}
  \drawlabel{0.65cm}{0.32cm}{hc}{t}{\tiny\sffamily\color{slate} (0,1)}
  \drawlabel{3.15cm}{0.32cm}{hc}{t}{\tiny\sffamily\color{slate} (1,1)}
  \drawlabel{1.30cm}{0.86cm}{hc}{t}{\tiny\sffamily\color{slate} (0,2)}
  \draw_end:
}

% Scatter plot for 2023-24 Q5(b)
\NewDocumentCommand{\scatterfig}{}{
  \draw_begin:
  \draw_set_linewidth:n {0.7pt}
  \color_select:n {ink}
  \draw_path_moveto:n {0cm, 0cm} \draw_path_lineto:n {6.8cm, 0cm}
  \draw_path_moveto:n {0cm, 0cm} \draw_path_lineto:n {0cm, 3.4cm}
  \draw_path_use_clear:n {stroke}
  \clist_map_inline:nn {{15, 0.96cm}, {20, 2.56cm}, {25, 4.16cm}, {30, 5.76cm}}{
    \color_select:n {ink}
    \draw_path_moveto:n {\clist_item:nn {##1}{2}, -0.08cm}
    \draw_path_lineto:n {\clist_item:nn {##1}{2}, 0.08cm}
    \draw_path_use_clear:n {stroke}
    \drawlabel{\clist_item:nn {##1}{2}}{-0.22cm}{hc}{t}{\scriptsize\sffamily \clist_item:nn {##1}{1}}
  }
  \clist_map_inline:nn {{3.0, 0.75cm}, {3.5, 1.50cm}, {4.0, 2.25cm}, {4.5, 3.00cm}}{
    \color_select:n {ink}
    \draw_path_moveto:n {-0.08cm, \clist_item:nn {##1}{2}}
    \draw_path_lineto:n {0.08cm, \clist_item:nn {##1}{2}}
    \draw_path_use_clear:n {stroke}
    \drawlabel{-0.18cm}{\clist_item:nn {##1}{2}}{r}{vc}{\scriptsize\sffamily \clist_item:nn {##1}{1}}
  }
  \draw_set_linewidth:n {1.1pt}
  \color_select:n {teal}
  \draw_path_moveto:n {0.64cm, 1.130cm} \draw_path_lineto:n {5.76cm, 2.255cm}
  \draw_path_use_clear:n {stroke}
  \clist_map_inline:nn {
    {0.96cm, 0.75cm}, {1.92cm, 2.25cm}, {3.20cm, 1.50cm}, {3.52cm, 1.50cm}, {5.44cm, 2.25cm}
  }{
    \color_select:n {coral}
    \draw_path_circle:nn {\clist_item:nn {##1}{1}, \clist_item:nn {##1}{2}} {2.6pt}
    \draw_path_use_clear:n {fill}
  }
  \drawlabel{7.0cm}{0cm}{l}{vc}{\small\sffamily Price~($x$,~\$)}
  \drawlabel{0cm}{3.65cm}{hc}{b}{\small\sffamily Rating~($y$)}
  \drawlabel{0.96cm}{0.45cm}{hc}{t}{\scriptsize\sffamily C(15, 3.0)}
  \drawlabel{1.92cm}{2.48cm}{hc}{b}{\scriptsize\sffamily D(18, 4.0)}
  \drawlabel{2.75cm}{1.22cm}{hc}{t}{\scriptsize\sffamily A(22, 3.5)}
  \drawlabel{4.15cm}{1.22cm}{hc}{t}{\scriptsize\sffamily E(23, 3.5)}
  \drawlabel{5.44cm}{2.48cm}{hc}{b}{\scriptsize\sffamily B(29, 4.0)}
  \draw_end:
}

% Histogram \& Ogive for 2022-23 Q5(b)
\NewDocumentCommand{\histogivefig}{}{
  \draw_begin:
  \draw_set_linewidth:n {0.7pt}
  \color_select:n {ink}
  % Left plot: Histogram (x from 0 to 5.6cm, 7 bins of width 0.75cm starting at 0.2cm)
  \draw_path_moveto:n {0cm, 0cm} \draw_path_lineto:n {5.8cm, 0cm}
  \draw_path_moveto:n {0cm, 0cm} \draw_path_lineto:n {0cm, 3.4cm}
  \draw_path_use_clear:n {stroke}
  % Frequencies: 4, 5, 10, 25, 9, 3, 4 (scale: 1 unit = 0.11cm)
  \clist_map_inline:nn {
    {0.20cm, 0.44cm, 4},
    {0.95cm, 0.55cm, 5},
    {1.70cm, 1.10cm, 10},
    {2.45cm, 2.75cm, 25},
    {3.20cm, 0.99cm, 9},
    {3.95cm, 0.33cm, 3},
    {4.70cm, 0.44cm, 4}
  }{
    \color_select:n {mist}
    \draw_path_rectangle:nn {\clist_item:nn {##1}{1}, 0cm} {0.75cm, \clist_item:nn {##1}{2}}
    \draw_path_use_clear:n {fill}
    \color_select:n {teal}
    \draw_path_rectangle:nn {\clist_item:nn {##1}{1}, 0cm} {0.75cm, \clist_item:nn {##1}{2}}
    \draw_path_use_clear:n {stroke}
    \drawlabel{\clist_item:nn {##1}{1} + 0.375cm}{\clist_item:nn {##1}{2} + 0.12cm}{hc}{b}{\scriptsize\sffamily \clist_item:nn {##1}{3}}
  }
  \clist_map_inline:nn {{0.20cm, 17.5}, {0.95cm, 23.5}, {1.70cm, 29.5}, {2.45cm, 35.5}, {3.20cm, 41.5}, {3.95cm, 47.5}, {4.70cm, 53.5}, {5.45cm, 59.5}}{
    \drawlabel{\clist_item:nn {##1}{1}}{-0.18cm}{hc}{t}{\tiny\sffamily \clist_item:nn {##1}{2}}
  }
  \drawlabel{2.8cm}{-0.55cm}{hc}{t}{\scriptsize\sffamily\bfseries (i)~Frequency~Histogram~(mm)}
  \drawlabel{0cm}{3.55cm}{hc}{b}{\scriptsize\sffamily Freq~($f$)}
  % Right plot: Ogive (shifted by 7.2cm)
  \color_select:n {ink}
  \draw_path_moveto:n {7.2cm, 0cm} \draw_path_lineto:n {13.0cm, 0cm}
  \draw_path_moveto:n {7.2cm, 0cm} \draw_path_lineto:n {7.2cm, 3.4cm}
  \draw_path_use_clear:n {stroke}
  % Cum freqs: 0, 4, 9, 19, 44, 53, 56, 60 (scale: 1 unit = 0.05cm)
  \draw_set_linewidth:n {1.4pt}
  \color_select:n {coral}
  \draw_path_moveto:n {7.40cm, 0cm}
  \draw_path_lineto:n {8.15cm, 0.20cm}
  \draw_path_lineto:n {8.90cm, 0.45cm}
  \draw_path_lineto:n {9.65cm, 0.95cm}
  \draw_path_lineto:n {10.40cm, 2.20cm}
  \draw_path_lineto:n {11.15cm, 2.65cm}
  \draw_path_lineto:n {11.90cm, 2.80cm}
  \draw_path_lineto:n {12.65cm, 3.00cm}
  \draw_path_use_clear:n {stroke}
  \clist_map_inline:nn {
    {7.40cm, 0cm, 17.5}, {8.15cm, 0.20cm, 23.5}, {8.90cm, 0.45cm, 29.5}, {9.65cm, 0.95cm, 35.5},
    {10.40cm, 2.20cm, 41.5}, {11.15cm, 2.65cm, 47.5}, {11.90cm, 2.80cm, 53.5}, {12.65cm, 3.00cm, 59.5}
  }{
    \color_select:n {ink}
    \draw_path_circle:nn {\clist_item:nn {##1}{1}, \clist_item:nn {##1}{2}} {1.8pt}
    \draw_path_use_clear:n {fill}
    \drawlabel{\clist_item:nn {##1}{1}}{-0.18cm}{hc}{t}{\tiny\sffamily \clist_item:nn {##1}{3}}
  }
  \drawlabel{10.0cm}{-0.55cm}{hc}{t}{\scriptsize\sffamily\bfseries (ii)~Cumulative~Frequency~Ogive}
  \drawlabel{7.2cm}{3.55cm}{hc}{b}{\scriptsize\sffamily Cum.\ Freq}
  \draw_end:
}

% Ogive for 2021-22 Q5(a) (Monthly wages)
\NewDocumentCommand{\wageogivefig}{}{
  \draw_begin:
  \draw_set_linewidth:n {0.7pt}
  \color_select:n {ink}
  \draw_path_moveto:n {0cm, 0cm} \draw_path_lineto:n {8.2cm, 0cm}
  \draw_path_moveto:n {0cm, 0cm} \draw_path_lineto:n {0cm, 3.6cm}
  \draw_path_use_clear:n {stroke}
  % Wage boundaries: 1900(0.3cm), 1950(1.0cm), 2000(1.7cm), 2050(2.4cm), 2100(3.1cm), 2150(3.8cm), 2200(4.5cm), 2400(7.3cm)
  % Cum freqs: 0, 6, 16, 38, 68, 84, 96, 111 (scale: 1 worker = 0.028cm)
  \draw_set_linewidth:n {1.5pt}
  \color_select:n {teal}
  \draw_path_moveto:n {0.3cm, 0cm}
  \draw_path_lineto:n {1.0cm, 0.168cm}
  \draw_path_lineto:n {1.7cm, 0.448cm}
  \draw_path_lineto:n {2.4cm, 1.064cm}
  \draw_path_lineto:n {3.1cm, 1.904cm}
  \draw_path_lineto:n {3.8cm, 2.352cm}
  \draw_path_lineto:n {4.5cm, 2.688cm}
  \draw_path_lineto:n {7.3cm, 3.108cm}
  \draw_path_use_clear:n {stroke}
  % Median line at N/2 = 55.5 -> y = 1.554cm, x = 2.4 + 0.583*0.7 = 2.808cm (2079.17 Tk)
  \draw_set_linewidth:n {0.8pt}
  \color_select:n {coral}
  \draw_path_moveto:n {0cm, 1.554cm} \draw_path_lineto:n {2.808cm, 1.554cm}
  \draw_path_lineto:n {2.808cm, 0cm}
  \draw_path_use_clear:n {stroke}
  \clist_map_inline:nn {
    {0.3cm, 0cm, 1900}, {1.0cm, 0.168cm, 1950}, {1.7cm, 0.448cm, 2000},
    {2.4cm, 1.064cm, 2050}, {3.1cm, 1.904cm, 2100}, {3.8cm, 2.352cm, 2150},
    {4.5cm, 2.688cm, 2200}, {7.3cm, 3.108cm, 2400}
  }{
    \color_select:n {ink}
    \draw_path_circle:nn {\clist_item:nn {##1}{1}, \clist_item:nn {##1}{2}} {1.8pt}
    \draw_path_use_clear:n {fill}
    \drawlabel{\clist_item:nn {##1}{1}}{-0.20cm}{hc}{t}{\tiny\sffamily \clist_item:nn {##1}{3}}
  }
  \drawlabel{-0.15cm}{1.554cm}{r}{vc}{\tiny\sffamily\color{coral} $N/2=55.5$}
  \drawlabel{2.81cm}{-0.42cm}{hc}{t}{\tiny\sffamily\bfseries\color{coral} Med~$\approx 2079.2$}
  \drawlabel{8.3cm}{0cm}{l}{vc}{\scriptsize\sffamily Wages~(Tk)}
  \drawlabel{0cm}{3.75cm}{hc}{b}{\scriptsize\sffamily Cum.\ Workers}
  \draw_end:
}
\ExplSyntaxOff

%------------------------------------------------------------ cover logo slot
% \logoslot{basename}: inline logo from basename.pdf/.png/.jpg if present
\newcommand{\logoslot}[1]{%
  \IfFileExists{#1.pdf}{\raisebox{-0.22\height}{\includegraphics[height=5mm]{#1.pdf}}\hspace{2.2mm}}{%
  \IfFileExists{#1.png}{\raisebox{-0.22\height}{\includegraphics[height=5mm]{#1.png}}\hspace{2.2mm}}{%
  \IfFileExists{#1.jpg}{\raisebox{-0.22\height}{\includegraphics[height=5mm]{#1.jpg}}\hspace{2.2mm}}{%
  \typeout{*** Cover: #1.(pdf|png|jpg) not found -- logo omitted ***}}}}}

%------------------------------------------------------------ text setup
\setlength{\parindent}{0pt}
\setlength{\parskip}{4pt}
\linespread{1.06}

%------------------------------------------------------------ handy macros
\newcommand{\ok}{\textcolor{teal}{\ensuremath{\checkmark}}}
\newcommand{\no}{\textcolor{coral}{\ensuremath{\times}}}
\newcommand{\key}[1]{\textcolor{teal}{\textbf{#1}}}
\newcommand{\ans}[1]{\par\smallskip\noindent\colorbox{mist}{\parbox{\dimexpr\linewidth-2\fboxsep\relax}{\vspace{2pt}\sffamily\small\textbf{\color{teal}Answer.}\ \color{ink}#1\vspace{2pt}}}\par}
\newcommand{\sub}[1]{\par\smallskip{\sffamily\bfseries\small\color{inksoft}#1}\par}
\newcommand{\fig}[1]{\par\smallskip\begin{center}#1\end{center}\vspace{-2pt}}

%------------------------------------------------------------ headers
\pagestyle{fancy}
\fancyhf{}
\renewcommand{\headrulewidth}{0pt}
\fancyhead[L]{\sffamily\scriptsize\color{slate}MATH 243 / 245 / 241 \textbullet\ PROBABILITY AND STATISTICS \textbullet\ SOLVED}
\fancyhead[R]{\sffamily\scriptsize\color{coral}\leftmark}
\fancyfoot[C]{\sffamily\small\color{ink}\thepage}
\fancypagestyle{plain}{\fancyhf{}\fancyfoot[C]{\sffamily\small\color{ink}\thepage}}

%------------------------------------------------------------ topic opener
\NewDocumentCommand{\topic}{o m m}{%
  \clearpage\phantomsection
  \addcontentsline{toc}{section}{\protect\numberline{#2}#3}%
  \IfValueTF{#1}{\markboth{\MakeUppercase{#1}}{}}{\markboth{\MakeUppercase{#3}}{}}% short running header if given
  \AddToShipoutPictureBG*{%
    \AtPageUpperLeft{%
      \put(0,-9mm){\textcolor{ink}{\rule{\paperwidth}{9mm}}}%
      \put(0,-10.2mm){\textcolor{coral}{\rule{\paperwidth}{1.2mm}}}%
    }%
  }%
  \vspace*{-2mm}
  {\noindent\sffamily\bfseries\fontsize{46}{46}\selectfont\color{coral}#2\par}
  \vspace{1mm}
  {\noindent\sffamily\bfseries\LARGE\color{ink}#3\par}
  \vspace{2mm}
  {\noindent\color{teal}\rule{0.22\linewidth}{2pt}\par}
  \vspace{4mm}}

%------------------------------------------------------------ question card
\newcommand{\marksbadge}[1]{{\setlength{\fboxsep}{2.4pt}\colorbox{coral}{\textcolor{white}{\sffamily\bfseries\scriptsize\,MARKS #1\,}}}}

\newsavebox{\qbox}
\newbox\qfullbox
\makeatletter
\newenvironment{question}[3]{%
  \def\q@title{\textcolor{ink}{#1}\enspace{\mdseries\textcolor{slate}{#2}}\hfill\marksbadge{#3}}%
  \par\vspace{10pt}\penalty10000%
  \begin{lrbox}{\qbox}%
  \begin{minipage}{\dimexpr\linewidth-3.5pt-20pt\relax}%
    \vspace{4pt}%
    {\sffamily\bfseries\small\q@title\par}\vspace{4pt}%
    \setlength{\parindent}{0pt}%
    \setlength{\parskip}{3pt}%
}{%
    \vspace{4pt}%
  \end{minipage}%
  \end{lrbox}%
  \setbox\qfullbox=\hbox{{\setlength{\fboxsep}{0pt}\colorbox{coral}{\hspace{3.5pt}{\setlength{\fboxsep}{10pt}\colorbox{sand}{\usebox{\qbox}}}}}}%
  \setbox\qfullbox=\hbox{\raise\dp\qfullbox\box\qfullbox}%
  \par\kern0pt\penalty10000%
  \ifdim\pagetotal>0pt
    \ifdim\pagegoal<\maxdimen
      \ifdim\dimexpr\pagegoal-\pagetotal\relax<\dimexpr\ht\qfullbox+80pt\relax
        \newpage
      \fi
    \fi
  \fi
  \noindent\box\qfullbox\par\vspace{3pt}%
}

%------------------------------------------------------------ repeated-question banner
\newcommand{\repeated}[1]{%
  \par\vspace{10pt}\kern0pt\penalty10000%
  \ifdim\pagetotal>0pt
    \ifdim\pagegoal<\maxdimen
      \ifdim\dimexpr\pagegoal-\pagetotal\relax<150pt
        \newpage
      \fi
    \fi
  \fi
  \noindent{\setlength{\fboxsep}{3.5pt}\colorbox{teal}{\textcolor{white}{\sffamily\bfseries\scriptsize REPEATED $\times$#1}}}\ %
  {\sffamily\scriptsize\color{teal}\bfseries Same question in #1 papers --- stacked below, solved once.}\par\vspace{-4pt}\nopagebreak%
}

%------------------------------------------------------------ source note
\newsavebox{\notebox}
\newbox\nfullbox
\newenvironment{note}{%
  \par\vspace{4pt}%
  \begin{lrbox}{\notebox}%
  \begin{minipage}{\dimexpr\linewidth-2.5pt-16pt\relax}%
    \vspace{2pt}\small{\sffamily\bfseries\scriptsize\color{ink}NOTE\quad}%
}{%
    \vspace{2pt}%
  \end{minipage}%
  \end{lrbox}%
  \setbox\nfullbox=\hbox{{\setlength{\fboxsep}{0pt}\colorbox{gold}{\hspace{2.5pt}{\setlength{\fboxsep}{8pt}\colorbox{goldbg}{\usebox{\notebox}}}}}}%
  \setbox\nfullbox=\hbox{\raise\dp\nfullbox\box\nfullbox}%
  \noindent\box\nfullbox\par\vspace{4pt}%
}

%------------------------------------------------------------ breakable solution panel
\newbox\sol@box
\newdimen\sol@avail
\newenvironment{solution}{%
  \par\vspace{4pt}\penalty10000%
  \setbox\sol@box\vbox\bgroup
    \hsize=\dimexpr\linewidth-14pt\relax
    \textwidth=\hsize
    \linewidth=\hsize
    \setlength{\parindent}{0pt}%
    \setlength{\parskip}{4pt}%
    \emergencystretch=2.5em
    {\sffamily\bfseries\scriptsize\color{teal}SOLUTION}\par\vspace{2pt}%
}{%
  \egroup
  \sol@split
  \par\vspace{10pt}%
}
\newcommand{\sol@split}{%
  \par\kern0pt\penalty10000%
  \ifvoid\sol@box\else
    \ifdim\pagetotal=0pt
      \sol@avail=\dimexpr\textheight-18pt\relax
    \else
      \ifdim\pagegoal=\maxdimen
        \sol@avail=\dimexpr\textheight-18pt\relax
      \else
        \sol@avail=\dimexpr\pagegoal-\pagetotal-22pt\relax
      \fi
    \fi
    % FIX: only jump to a new page when there is truly no room (< ~3 lines)
    \ifdim\dimexpr\ht\sol@box+\dp\sol@box\relax>\sol@avail
      \ifdim\sol@avail<45pt
        \newpage
        \sol@avail=\dimexpr\textheight-18pt\relax
      \fi
    \fi
    \ifdim\dimexpr\ht\sol@box+\dp\sol@box\relax>\sol@avail
      {\vbadness=10000 \vfuzz=\maxdimen
       \splittopskip=0pt
       \splitmaxdepth=\maxdepth
       \global\setbox\@ne=\vsplit\sol@box to \sol@avail
       \global\setbox\tw@=\vbox{\unvcopy\@ne}%
       % FIX: \vsplit may legitimately shrink display/list spacing to fit;
       % keep the piece at its fitted height instead of re-measuring it at
       % natural height (that made whole solutions jump to the next page).
       \ifdim\dimexpr\ht\tw@+\dp\tw@\relax>\sol@avail
         \global\setbox\@ne=\vbox to \sol@avail{\unvbox\@ne}%
       \else
         \global\setbox\@ne=\box\tw@
       \fi}%
      \ifdim\ht\@ne>\dimexpr\sol@avail+8pt\relax
        % a single unbreakable block taller than the space left: move it on
        % (unless we are already at the top of a page -- then just print it)
        \ifdim\sol@avail<\dimexpr\textheight-40pt\relax
          \global\setbox\sol@box=\vbox{\unvbox\@ne\unvbox\sol@box}%
          \newpage
        \else
          \noindent\textcolor{teal}{\vrule width 1.6pt}\hspace{11pt}\box\@ne\par\penalty\z@ % FIX: legal page break after each piece
          \ifvoid\sol@box\else\newpage\fi
        \fi
      \else
        \noindent\textcolor{teal}{\vrule width 1.6pt}\hspace{11pt}\box\@ne\par\penalty\z@ % FIX: legal page break after each piece
        \ifvoid\sol@box\else\newpage\fi
      \fi
      \expandafter\sol@split
    \else
      \noindent\textcolor{teal}{\vrule width 1.6pt}\hspace{11pt}\vbox{\unvbox\sol@box}\par
    \fi
  \fi
}
\makeatother

\hypersetup{pdftitle={BUET MATH 243 / 245 / 241 Probability and Statistics -- Topic-wise Solved Term-Final Questions},pdfauthor={MAHDI HASAN}}

\begin{document}

%===================================================================== COVER
\thispagestyle{empty}
\AddToShipoutPictureBG*{%
  \AtPageUpperLeft{%
    \put(0,-150mm){\textcolor{ink}{\rule{\paperwidth}{150mm}}}%
    \put(0,-153mm){\textcolor{coral}{\rule{\paperwidth}{3mm}}}%
    \put(20mm,-25mm){\textcolor{coral}{\sffamily\bfseries\large BUET \textbullet\ MATH 243 / 245 / 241}}%
    \put(19mm,-42mm){\parbox[t]{120mm}{\sffamily\bfseries\fontsize{40}{44}\selectfont\color{white}Probability and\\Statistics}}%
    \put(20mm,-78mm){\parbox[t]{110mm}{\sffamily\Large\color{gold}Topic-wise questions with\\complete worked solutions}}%
    \put(114mm,-122mm){\coverplot}%
    \put(20mm,-172mm){\parbox[t]{170mm}{\sffamily
      \begin{tabular}{@{}l@{\hspace{12mm}}l@{\hspace{12mm}}l@{\hspace{12mm}}l@{}}
      {\color{coral}\fontsize{26}{26}\selectfont\bfseries 18} &
      {\color{coral}\fontsize{26}{26}\selectfont\bfseries 125} &
      {\color{coral}\fontsize{26}{26}\selectfont\bfseries 12} &
      {\color{coral}\fontsize{26}{26}\selectfont\bfseries 3}\\[2mm]
      \small\color{ink} topic sections & \small\color{ink} exam questions & \small\color{ink} exam sessions & \small\color{ink} course codes
      \end{tabular}}}%
    \put(20mm,-202mm){\parbox[t]{170mm}{\sffamily\small\color{slate}%
      MATH 243, MATH 245, and MATH 241 term finals (CSE), sessions 2011--2012 to 2023--2024.\\[2mm]
      Every Probability, Statistics, and Out-of-Syllabus question from the exam papers is organized
      topic-wise (most recent session first, lowest course number first) and followed by a complete,
      step-by-step descriptive and mathematical solution. Diagrams, histograms, ogives, scatter plots,
      and circuits are redrawn in vector graphics.}}%
    \put(90mm,-238mm){\textcolor{teal}{\rule{30mm}{1.4mm}}}%
    \put(0,-248mm){\parbox[t]{\paperwidth}{\centering\sffamily
      {\fontsize{12}{15}\selectfont\bfseries\color{ink}\preparedby}\\[3.5mm]
      {\small\color{ink}\logoslot{gemini-logo}Solutions created using \textbf{Gemini}}\\[2.6mm]
      {\small\color{ink}\logoslot{claude-logo}Layout fixed, solutions checked \& corrected by \textbf{Claude}}}}%
  }%
}
\mbox{}
\clearpage

%===================================================================== CONTENTS
\pagenumbering{roman}
\thispagestyle{plain}
{\sffamily\bfseries\Huge\color{ink}Contents\par}\vspace{2mm}
{\color{teal}\rule{0.2\linewidth}{2pt}}\vspace{6mm}
\makeatletter
\renewcommand*\l@section[2]{%
  \par\noindent\vspace{3pt}{\sffamily\large #1}\nobreak{\color{slate!60}\leaders\hbox{\,.\,}\hfill}\nobreak{\sffamily\color{coral}\bfseries#2}\par}
\renewcommand*\numberline[1]{\makebox[13mm][l]{\color{coral}\bfseries#1}}
\@starttoc{toc}
\makeatother
\clearpage

%===================================================================== HOW TO READ
\thispagestyle{plain}
{\sffamily\bfseries\Huge\color{ink}How to read this book\par}\vspace{2mm}
{\color{teal}\rule{0.2\linewidth}{2pt}}\vspace{6mm}

\begin{question}{2099--2100 \textbullet\ Q0(a)}{Course \textbullet\ Date \textbullet\ ID}{10}
A question card. The coral badge on the right shows the printed marks. Tables and figures from the original
exam paper are re-typeset or redrawn in vector graphics inside the card.
\end{question}
\begin{solution}
The solution panel always follows its question with complete step-by-step mathematical derivations, intermediate tables, and descriptive explanations. Key results are highlighted:
\ans{like this.}
\end{solution}

\repeated{2}
\begin{question}{Paper A}{Q1}{12}\small A question that was asked in a recent session\dots\end{question}
\begin{question}{Paper B}{Q4}{15}\small \dots and asked again in an earlier session\dots\end{question}
\begin{solution}One comprehensive shared solution covers every stacked card.\end{solution}

\begin{note}Gold boxes flag a misprint or discrepancy in the source scan (such as a swapped label or typo in a piecewise function) and state the exact reading used in the solution.\end{note}

\sub{Conventions \& Ordering Rules}
\begin{itemize}
\item \textbf{Sorting Order:} Within each topic, questions are ordered from the \textbf{most recent session first} (\texttt{2023--2024} down to \texttt{2011--2012}); when sessions are identical, the lowest course number (\texttt{MATH 241} $<$ \texttt{MATH 243} $<$ \texttt{MATH 245}) appears first.
\item \textbf{Standard Notation:} $\Phi(z) = P(Z \le z)$ denotes the cumulative distribution function (CDF) of a standard normal variable $Z \sim N(0,1)$; $t_{\alpha,\nu}$, $\chi^2_{\alpha,\nu}$, and $F_{\alpha(\nu_1,\nu_2)}$ denote upper-tail critical values with $\nu$ degrees of freedom.
\item \textbf{Moments \& Shape Coefficients:} Central moments are $\mu_r = \frac{1}{N}\sum f_i(x_i - \bar{x})^r$; Karl Pearson's shape coefficients are $\beta_1 = \mu_3^2/\mu_2^3$ and $\beta_2 = \mu_4/\mu_2^2$, with moment coefficient of skewness $\gamma_1 = \mu_3/\mu_2^{3/2}$ and excess kurtosis $\gamma_2 = \beta_2 - 3$.
\item \textbf{Cross-Listed Topics:} Questions that span multiple syllabus bullets (e.g.\ a joint distribution problem that also asks for marginal distributions and covariance) are solved in full under their primary topic and cross-referenced in secondary topics.
\end{itemize}
\clearpage
\pagenumbering{arabic}

%===================================================================== TOPIC 01
\topic{01}{Introduction to Probability}

%---------------------------------------------------------------- P001
\begin{question}{2023--2024 \textbullet\ Q1(a)}{MATH 243 \textbullet\ 21/05/2026 \textbullet\ P001}{6}
Can two events be mutually exclusive and independent simultaneously? Justify your answer.
\end{question}
\begin{solution}
Let $A$ and $B$ be two events defined on the same sample space $S$.
\begin{enumerate}[(i)]
\item \textbf{Condition for Mutually Exclusive (Disjoint) Events:} Two events $A$ and $B$ are mutually exclusive if they cannot occur together, meaning $A \cap B = \emptyset$. By the axioms of probability,
\[ P(A \cap B) = 0. \]
\item \textbf{Condition for Independent Events:} Two events $A$ and $B$ are statistically independent if the occurrence of one does not affect the probability of the other, which requires
\[ P(A \cap B) = P(A)\,P(B). \]
\end{enumerate}
\textbf{Simultaneous Satisfaction:} For $A$ and $B$ to be both mutually exclusive and independent simultaneously, both conditions must hold at the same time:
\[ P(A)\,P(B) = P(A \cap B) = 0 \implies P(A) = 0 \quad \text{or} \quad P(B) = 0. \]
Therefore:
\begin{itemize}
\item If $A$ and $B$ are \textbf{non-trivial events} with positive probabilities ($P(A) > 0$ and $P(B) > 0$), then $P(A)P(B) > 0 \neq 0$. In this case, knowing that $A$ has occurred immediately implies $B$ cannot occur ($P(B \mid A) = 0 \neq P(B)$), so they \textbf{cannot} be mutually exclusive and independent at the same time.
\item Only in the \textbf{trivial degenerate case} where at least one event has zero probability ($P(A) = 0$ or $P(B) = 0$, such as the impossible event $\emptyset$) can two events satisfy both definitions simultaneously.
\end{itemize}
\ans{No, two events with non-zero probabilities ($P(A)>0$ and $P(B)>0$) cannot be mutually exclusive and independent simultaneously, because $P(A\cap B)=0$ and $P(A\cap B)=P(A)P(B)>0$ are mutually contradictory.}
\end{solution}

%---------------------------------------------------------------- P002
\begin{question}{2023--2024 \textbullet\ Q1(b)}{MATH 243 \textbullet\ 21/05/2026 \textbullet\ P002}{12}
Define multiplication law of probability. Suppose we have two boxes $A$ and $B$: the first box contains 5 white balls and 7 red balls, while the second box contains 3 white balls and 8 red balls. One ball is transferred unseen from the first box to the second box, and then a ball is drawn from the second box. What is the probability that this ball is red?
\end{question}
\begin{solution}
\textbf{Part 1: Multiplication Law of Probability.}
For any two events $E_1$ and $E_2$ in a sample space with $P(E_1) > 0$, the probability that both $E_1$ and $E_2$ occur simultaneously is equal to the probability of $E_1$ multiplied by the conditional probability of $E_2$ given that $E_1$ has already occurred:
\[ P(E_1 \cap E_2) = P(E_1)\,P(E_2 \mid E_1) = P(E_2)\,P(E_1 \mid E_2). \]
If $E_1$ and $E_2$ are independent, this reduces to $P(E_1 \cap E_2) = P(E_1)\,P(E_2)$.

\medskip
\textbf{Part 2: Two-Box Transfer Problem.}
\begin{itemize}
\item \textbf{Initial composition of Box $A$:} $5\text{ white} + 7\text{ red} = 12\text{ balls}$.
\item \textbf{Initial composition of Box $B$:} $3\text{ white} + 8\text{ red} = 11\text{ balls}$.
\end{itemize}
Let $W_1$ and $R_1$ denote the mutually exclusive and exhaustive events of transferring a white ball or a red ball from Box $A$ to Box $B$, respectively, and let $R_2$ be the event of drawing a red ball from Box $B$ after the transfer.
\begin{enumerate}[(i)]
\item \textbf{Case 1 (White ball transferred, $W_1$):}
\[ P(W_1) = \frac{5}{12}. \]
Box $B$ now contains $4\text{ white}$ and $8\text{ red}$ balls ($12\text{ balls}$ total), so
\[ P(R_2 \mid W_1) = \frac{8}{12}. \]
\item \textbf{Case 2 (Red ball transferred, $R_1$):}
\[ P(R_1) = \frac{7}{12}. \]
Box $B$ now contains $3\text{ white}$ and $9\text{ red}$ balls ($12\text{ balls}$ total), so
\[ P(R_2 \mid R_1) = \frac{9}{12}. \]
\end{enumerate}
By the law of total probability and the multiplication rule:
\[ \begin{aligned} P(R_2) &= P(W_1)\,P(R_2 \mid W_1) + P(R_1)\,P(R_2 \mid R_1) \\ &= \left(\frac{5}{12}\times\frac{8}{12}\right) + \left(\frac{7}{12}\times\frac{9}{12}\right) = \frac{40 + 63}{144} = \frac{103}{144} \approx 0.7153. \end{aligned} \]
\ans{$P(\text{ball drawn from Box } B \text{ is red}) = \dfrac{103}{144} \approx 0.7153$.}
\end{solution}

%---------------------------------------------------------------- P003
\begin{question}{2022--2023 \textbullet\ Q1(a)}{MATH 243 \textbullet\ 01/02/2025 \textbullet\ P003}{8}
One bag contains 4 white balls and 3 black balls, and a second bag contains 3 white balls and 5 black balls. One ball is drawn from the first bag and placed unseen in the second bag. What is the probability that a ball now drawn from the second bag is black?
\end{question}
\begin{solution}
\begin{itemize}
\item \textbf{Initial composition of Bag 1:} $4\text{ white} + 3\text{ black} = 7\text{ balls}$.
\item \textbf{Initial composition of Bag 2:} $3\text{ white} + 5\text{ black} = 8\text{ balls}$.
\end{itemize}
Let $W_1$ and $B_1$ denote the events that the ball transferred from Bag 1 to Bag 2 is white or black, respectively, and let $B_2$ be the event that the ball subsequently drawn from Bag 2 is black:
\begin{enumerate}[(i)]
\item \textbf{If a white ball is transferred ($W_1$):} $P(W_1) = \dfrac{4}{7}$. Bag 2 then has $4\text{ white}$ and $5\text{ black}$ balls ($9\text{ total}$), giving $P(B_2 \mid W_1) = \dfrac{5}{9}$.
\item \textbf{If a black ball is transferred ($B_1$):} $P(B_1) = \dfrac{3}{7}$. Bag 2 then has $3\text{ white}$ and $6\text{ black}$ balls ($9\text{ total}$), giving $P(B_2 \mid B_1) = \dfrac{6}{9}$.
\end{enumerate}
Applying the law of total probability:
\[ P(B_2) = P(W_1)\,P(B_2 \mid W_1) + P(B_1)\,P(B_2 \mid B_1) = \left(\frac{4}{7}\times\frac{5}{9}\right) + \left(\frac{3}{7}\times\frac{6}{9}\right) = \frac{20 + 18}{63} = \frac{38}{63} \approx 0.6032. \]
\ans{$P(\text{ball drawn from second bag is black}) = \dfrac{38}{63} \approx 0.6032$.}
\end{solution}

%---------------------------------------------------------------- P004
\begin{question}{2021--2022 \textbullet\ Q6(a)}{MATH 245 \textbullet\ 30/04/2023 \textbullet\ P004}{10}
Suppose 500 lottery tickets are sold, and among them, 200 tickets pay off at least the cost of the ticket. Now suppose you buy 5 tickets. Find the probability that you will win back at least the cost of 3 tickets.
\end{question}
\begin{solution}
The problem interprets ``winning back at least the cost of 3 tickets'' as obtaining at least $3$ winning tickets among the $n = 5$ tickets purchased from the finite population of $N = 500$ tickets ($K = 200$ winning tickets, $N - K = 300$ non-winning tickets).

\medskip
\textbf{Method 1: Exact Hypergeometric Distribution (Sampling without Replacement).}
Let $X$ be the number of winning tickets among the 5 purchased. Then
\[ P(X = x) = \frac{\binom{200}{x}\binom{300}{5-x}}{\binom{500}{5}}, \qquad x = 0, 1, 2, 3, 4, 5. \]
Here $\binom{500}{5} = 255{,}244{,}687{,}600$. Summing over $x = 3, 4, 5$:
\begin{align*}
P(X = 3) &= \frac{\binom{200}{3}\binom{300}{2}}{\binom{500}{5}} = \frac{1{,}313{,}400 \times 44{,}850}{255{,}244{,}687{,}600} \approx 0.230782, \\[3pt]
P(X = 4) &= \frac{\binom{200}{4}\binom{300}{1}}{\binom{500}{5}} = \frac{64{,}684{,}950 \times 300}{255{,}244{,}687{,}600} \approx 0.076027, \\[3pt]
P(X = 5) &= \frac{\binom{200}{5}\binom{300}{0}}{\binom{500}{5}} = \frac{2{,}535{,}650{,}040 \times 1}{255{,}244{,}687{,}600} \approx 0.009934.
\end{align*}
Adding the three terms gives:
\[ P(X \ge 3) = 0.230782 + 0.076027 + 0.009934 = 0.316744. \]

\medskip
\textbf{Method 2: Binomial Approximation (since $n/N = 5/500 = 0.01 \ll 0.05$).}
With $p = \dfrac{200}{500} = 0.4$ and $q = 1 - p = 0.6$:
\[ \begin{aligned} P(X \ge 3) &= \sum_{x=3}^{5} \binom{5}{x}(0.4)^x(0.6)^{5-x} = 10(0.4)^3(0.6)^2 + 5(0.4)^4(0.6) + (0.4)^5 \\ &= 0.2304 + 0.0768 + 0.01024 = 0.31744. \end{aligned} \]
\ans{$P(X \ge 3) = 0.3167$ (exact hypergeometric) or $0.3174$ (binomial approximation).}
\end{solution}

%---------------------------------------------------------------- P005
\begin{question}{2016--2017 \textbullet\ Q7(a)}{MATH 245 \textbullet\ 05/08/2017 \textbullet\ P005}{10}
The following circuit operates if and only if there is a path of functional devices from left to right. Assume that devices fail independently and that the probability of failure of each device is as shown below. What is the probability that the circuit does not operate?
\fig{\circuitfig}
\end{question}
\begin{solution}
The circuit consists of three main parallel branches between the left and right terminals:
\begin{enumerate}[(i)]
\item \textbf{Top branch ($B_1$):} Contains a single device with failure probability $P(B_1\text{ fails}) = 0.02$.
\item \textbf{Bottom branch ($B_3$):} Contains a single device with failure probability $P(B_3\text{ fails}) = 0.02$.
\item \textbf{Middle branch ($B_2$):} Consists of two identical parallel sub-branches ($M_1$ and $M_2$). Each sub-branch contains two devices in series, each with failure probability $0.01$ (and thus working probability $1 - 0.01 = 0.99$).
\begin{itemize}
\item A series sub-branch operates if and only if both of its devices operate:
\[ P(M_i\text{ works}) = 0.99 \times 0.99 = 0.9801 \implies P(M_i\text{ fails}) = 1 - 0.9801 = 0.0199 \quad (i=1,2). \]
\item Because $M_1$ and $M_2$ are in parallel, the middle branch $B_2$ fails if and only if both $M_1$ and $M_2$ fail:
\[ P(B_2\text{ fails}) = P(M_1\text{ fails})\,P(M_2\text{ fails}) = (0.0199)^2 = 0.00039601. \]
\end{itemize}
\end{enumerate}
\textbf{Overall Circuit Failure Probability:} Since the three main branches $B_1, B_2, B_3$ are connected in parallel, the entire circuit fails to operate if and only if all three branches fail simultaneously:
\[ \begin{aligned} P(\text{Circuit fails}) &= P(B_1\text{ fails})\,P(B_2\text{ fails})\,P(B_3\text{ fails}) \\ &= (0.02)\times(0.00039601)\times(0.02) = 1.58404 \times 10^{-7}. \end{aligned} \]
\ans{$P(\text{circuit does not operate}) = 1.58404 \times 10^{-7}$ (so $P(\text{circuit operates}) \approx 0.99999984$).}
\end{solution}

%---------------------------------------------------------------- P006
\begin{question}{2014--2015 \textbullet\ Q8(a)}{MATH 241 \textbullet\ 06/02/2016 \textbullet\ P006}{12}
A box contains 8 red, 3 white, and 9 blue balls. If three balls are drawn at random, determine the probability that: (i) all 3 are red, (ii) 2 are red and 1 is white, (iii) at least 1 is white.
\end{question}
\begin{solution}
Total number of balls in the box is $N = 8 + 3 + 9 = 20$. The total number of equally likely ways to draw $3$ balls out of $20$ is:
\[ n(S) = \binom{20}{3} = \frac{20 \times 19 \times 18}{3 \times 2 \times 1} = 1140. \]
\begin{enumerate}[(i)]
\item \textbf{All 3 balls are red:} The number of ways to choose 3 red balls from 8 red balls is $\binom{8}{3} = 56$.
\[ P(\text{all 3 red}) = \frac{\binom{8}{3}}{\binom{20}{3}} = \frac{56}{1140} = \frac{14}{285} \approx 0.04912. \]
\item \textbf{2 are red and 1 is white:} The number of ways to choose 2 red balls from 8 and 1 white ball from 3 is $\binom{8}{2}\binom{3}{1} = 28 \times 3 = 84$.
\[ P(\text{2 red, 1 white}) = \frac{\binom{8}{2}\binom{3}{1}}{\binom{20}{3}} = \frac{84}{1140} = \frac{7}{95} \approx 0.07368. \]
\item \textbf{At least 1 ball is white:} Using the complement event (none of the 3 balls is white, meaning all 3 are chosen from the $8 + 9 = 17$ non-white balls):
\[ P(\text{at least 1 white}) = 1 - P(\text{no white}) = 1 - \frac{\binom{17}{3}}{\binom{20}{3}} = 1 - \frac{680}{1140} = \frac{460}{1140} = \frac{23}{57} \approx 0.40351. \]
\end{enumerate}
\ans{(i) $\dfrac{14}{285} \approx 0.0491$, \quad (ii) $\dfrac{7}{95} \approx 0.0737$, \quad (iii) $\dfrac{23}{57} \approx 0.4035$.}
\end{solution}

%---------------------------------------------------------------- P007
\begin{question}{2013--2014 \textbullet\ Q7(b)(i)}{MATH 241 \textbullet\ 13/12/2014 \textbullet\ P007}{7}
3 cards are drawn from a deck of 52 cards one after another without replacement. Find the probability that they will come from different suits.
\end{question}
\begin{solution}
A standard deck of 52 cards consists of 4 suits (Spades, Hearts, Diamonds, Clubs) with 13 cards in each suit. Let $E_1, E_2, E_3$ be the events that the 1st, 2nd, and 3rd cards belong to distinct suits:
\begin{itemize}
\item \textbf{First card ($E_1$):} Can be of any suit, so $P(E_1) = \dfrac{52}{52} = 1$.
\item \textbf{Second card ($E_2 \mid E_1$):} Out of the 51 remaining cards, $52 - 13 = 39$ cards belong to the 3 other suits, so $P(E_2 \mid E_1) = \dfrac{39}{51}$.
\item \textbf{Third card ($E_3 \mid E_1 \cap E_2$):} Out of the 50 remaining cards, $52 - 26 = 26$ cards belong to the 2 remaining unused suits, so $P(E_3 \mid E_1 \cap E_2) = \dfrac{26}{50}$.
\end{itemize}
By the multiplication rule of probability:
\[ P(\text{3 different suits}) = 1 \times \frac{39}{51} \times \frac{26}{50} = \frac{13}{17} \times \frac{13}{25} = \frac{169}{425} \approx 0.39765. \]
\textit{(Combinatorial check: $\dfrac{\binom{4}{3}\times 13^3}{\binom{52}{3}} = \dfrac{4 \times 2197}{22100} = \dfrac{8788}{22100} = \dfrac{169}{425}$.)}
\ans{$P(\text{3 cards from different suits}) = \dfrac{169}{425} \approx 0.3976$.}
\end{solution}

%---------------------------------------------------------------- P008
\begin{question}{2011--2012 \textbullet\ Q7(b)(ii)}{MATH 241 \textbullet\ 07/01/2013 \textbullet\ P008}{8}
The probability that $A$ can solve a problem is $\frac{3}{7}$, that $B$ can solve it is $\frac{2}{3}$, and that $C$ can solve it is $\frac{4}{5}$. If all of them try independently, find the probability that the problem will be solved.
\end{question}
\begin{solution}
Given the individual probabilities of solving the problem:
\[ P(A) = \frac{3}{7}, \qquad P(B) = \frac{2}{3}, \qquad P(C) = \frac{4}{5}. \]
The probabilities that each person \textbf{fails} to solve the problem are:
\[ P(A^c) = 1 - \frac{3}{7} = \frac{4}{7}, \qquad P(B^c) = 1 - \frac{2}{3} = \frac{1}{3}, \qquad P(C^c) = 1 - \frac{4}{5} = \frac{1}{5}. \]
Since $A, B,$ and $C$ attempt the problem independently, the probability that \textbf{none} of them solves it is:
\[ P(A^c \cap B^c \cap C^c) = P(A^c)\,P(B^c)\,P(C^c) = \frac{4}{7} \times \frac{1}{3} \times \frac{1}{5} = \frac{4}{105}. \]
The problem is solved if at least one of $A, B,$ or $C$ solves it:
\[ P(A \cup B \cup C) = 1 - P(A^c \cap B^c \cap C^c) = 1 - \frac{4}{105} = \frac{101}{105} \approx 0.9619. \]
\ans{$P(\text{problem will be solved}) = \dfrac{101}{105} \approx 0.9619$.}
\end{solution}

%===================================================================== TOPIC 02
\topic[Odds \& Conditional Probability]{02}{Odds, Joint Probability \& Conditional Probability (Bayes' Theorem)}

\begin{note}
\textbf{Syllabus Coverage Note:} No standalone question on \emph{Odds} appeared in the 2011--2024 papers (for reference, the odds in favor of event $E$ are $P(E):P(E^c)$, and odds against $E$ are $P(E^c):P(E)$). \emph{Joint Probability Tables} for discrete random variables are solved in full under \textbf{Topic 04 (P019)}, while bivariate event intersections are solved below in \textbf{P012} and \textbf{P013}.
\end{note}

%---------------------------------------------------------------- P009
\begin{question}{2023--2024 \textbullet\ Q1(c)}{MATH 243 \textbullet\ 21/05/2026 \textbullet\ P009}{17}
An opinion poll in a city shows that 45\% of city dwellers support party $A$, 40\% support party $B$, and the remaining 15\% support other parties $C$. Previous records reveal that in city elections, 65\% of $A$, 80\% of $B$, and 50\% of other party supporters turned up to cast their votes. A person in the city is chosen at random and it is learned that he did not cast his vote in the last election. What is the probability that he is a supporter of $A$? $B$? $C$?
\end{question}
\begin{solution}
Let $A, B,$ and $C$ denote the events that a randomly chosen citizen supports party $A$, party $B$, or other parties $C$, respectively:
\[ P(A) = 0.45, \qquad P(B) = 0.40, \qquad P(C) = 0.15. \]
Let $V$ denote the event that a citizen cast their vote in the last election, and $V^c$ the event that they \textbf{did not} cast their vote. From the given voter turnout rates:
\begin{align*}
P(V \mid A) &= 0.65 \implies P(V^c \mid A) = 1 - 0.65 = 0.35, \\
P(V \mid B) &= 0.80 \implies P(V^c \mid B) = 1 - 0.80 = 0.20, \\
P(V \mid C) &= 0.50 \implies P(V^c \mid C) = 1 - 0.50 = 0.50.
\end{align*}
By the law of total probability, the overall probability that a randomly selected citizen did not vote is:
\begin{align*}
P(V^c) &= P(A)P(V^c \mid A) + P(B)P(V^c \mid B) + P(C)P(V^c \mid C) \\
       &= (0.45)(0.35) + (0.40)(0.20) + (0.15)(0.50) \\
       &= 0.1575 + 0.0800 + 0.0750 = 0.3125.
\end{align*}
Applying Bayes' theorem for each party:
\begin{enumerate}[(i)]
\item \textbf{Supporter of Party $A$:}
\[ P(A \mid V^c) = \frac{P(A)P(V^c \mid A)}{P(V^c)} = \frac{0.1575}{0.3125} = \frac{63}{125} = 0.504. \]
\item \textbf{Supporter of Party $B$:}
\[ P(B \mid V^c) = \frac{P(B)P(V^c \mid B)}{P(V^c)} = \frac{0.0800}{0.3125} = \frac{32}{125} = 0.256. \]
\item \textbf{Supporter of Other Parties $C$:}
\[ P(C \mid V^c) = \frac{P(C)P(V^c \mid C)}{P(V^c)} = \frac{0.0750}{0.3125} = \frac{30}{125} = \frac{6}{25} = 0.240. \]
\end{enumerate}
\textit{(Check: $0.504 + 0.256 + 0.240 = 1.000$.)}
\ans{$P(A \mid V^c) = 0.504 \ \left(\dfrac{63}{125}\right)$, \quad $P(B \mid V^c) = 0.256 \ \left(\dfrac{32}{125}\right)$, \quad $P(C \mid V^c) = 0.240 \ \left(\dfrac{6}{25}\right)$.}
\end{solution}

%---------------------------------------------------------------- P010
\begin{question}{2022--2023 \textbullet\ Q1(b)}{MATH 243 \textbullet\ 01/02/2025 \textbullet\ P010}{15}
In an assembly plant, three machines $B_1, B_2$, and $B_3$ produce 30\%, 45\%, and 25\% of the products, respectively. Past experience shows that 2\%, 3\%, and 2\% of the products made by each machine, respectively, are defective. If a finished product is randomly selected and found to be defective, what is the probability that this product was made by machine $B_i$ ($i = 1, 2, 3$)?
\end{question}
\begin{solution}
Let $B_1, B_2, B_3$ denote the events that a randomly selected product is manufactured by machine $B_1, B_2, B_3$, respectively, and let $D$ denote the event that the product is defective:
\[ \begin{aligned} P(B_1) &= 0.30, & P(B_2) &= 0.45, & P(B_3) &= 0.25, \\ P(D \mid B_1) &= 0.02, & P(D \mid B_2) &= 0.03, & P(D \mid B_3) &= 0.02. \end{aligned} \]
By the law of total probability, the overall probability of selecting a defective product is:
\begin{align*}
P(D) &= P(B_1)P(D \mid B_1) + P(B_2)P(D \mid B_2) + P(B_3)P(D \mid B_3) \\
     &= (0.30)(0.02) + (0.45)(0.03) + (0.25)(0.02) = 0.0060 + 0.0135 + 0.0050 = 0.0245.
\end{align*}
By Bayes' theorem, the posterior probabilities for each machine given $D$ are:
\begin{align*}
P(B_1 \mid D) &= \frac{P(B_1)P(D \mid B_1)}{P(D)} = \frac{0.0060}{0.0245} = \frac{12}{49} \approx 0.2449, \\[4pt]
P(B_2 \mid D) &= \frac{P(B_2)P(D \mid B_2)}{P(D)} = \frac{0.0135}{0.0245} = \frac{27}{49} \approx 0.5510, \\[4pt]
P(B_3 \mid D) &= \frac{P(B_3)P(D \mid B_3)}{P(D)} = \frac{0.0050}{0.0245} = \frac{10}{49} \approx 0.2041.
\end{align*}
\ans{$P(B_1 \mid D) = \dfrac{12}{49} \approx 0.2449$, \quad $P(B_2 \mid D) = \dfrac{27}{49} \approx 0.5510$, \quad $P(B_3 \mid D) = \dfrac{10}{49} \approx 0.2041$.}
\end{solution}

%---------------------------------------------------------------- P011
\begin{question}{2019--2020 \textbullet\ Q7(a)}{MATH 245 \textbullet\ 25/09/2021 \textbullet\ P011}{17}
A truth serum properly judges 90\% of guilty suspects, while 10\% of guilty suspects are improperly found innocent; innocent suspects are misjudged 1\% of the time. If a suspect is selected from a group of suspects of which only 5\% have ever committed a crime, and the serum indicates that he is guilty, what is the probability that he is innocent?
\end{question}
\begin{solution}
Define the events:
\begin{itemize}
\item $G$: the suspect has actually committed a crime (guilty), with prior $P(G) = 0.05$.
\item $I = G^c$: the suspect is actually innocent, with prior $P(I) = 1 - 0.05 = 0.95$.
\item $T_G$: the truth serum indicates that the suspect is guilty.
\end{itemize}
From the reliability characteristics of the serum:
\[ P(T_G \mid G) = 0.90, \qquad P(T_G \mid I) = 0.01. \]
By the law of total probability, the probability that the serum indicates guilt is:
\[ P(T_G) = P(G)P(T_G \mid G) + P(I)P(T_G \mid I) = (0.05)(0.90) + (0.95)(0.01) = 0.0450 + 0.0095 = 0.0545. \]
Using Bayes' theorem, the probability that the suspect is actually innocent given that the serum indicates guilt is:
\[ P(I \mid T_G) = \frac{P(I)P(T_G \mid I)}{P(T_G)} = \frac{0.0095}{0.0545} = \frac{95}{545} = \frac{19}{109} \approx 0.1743. \]
\ans{$P(\text{Innocent} \mid \text{Serum indicates guilty}) = \dfrac{19}{109} \approx 0.1743$ ($17.43\%$).}
\end{solution}

%---------------------------------------------------------------- P012
\begin{question}{2018--2019 \textbullet\ Q6(b)}{MATH 245 \textbullet\ 03/11/2019 \textbullet\ P012}{15}
A group of 20 students went on a picnic. Among them, 5 got sunburnt, 8 got bitten by mosquitoes, and 10 returned home safely without any mishap. What is the probability that: (i) a boy who got sunburnt was ignored by the mosquitoes? and (ii) a mosquito-bitten boy also got sunburnt?
\end{question}
\begin{solution}
Let $S$ be the event that a student got sunburnt, and $M$ be the event that a student was bitten by mosquitoes, in a group of $n(\Omega) = 20$ students:
\[ n(S) = 5, \qquad n(M) = 8, \qquad n((S \cup M)^c) = 10. \]
Therefore, the number of students who suffered at least one mishap is:
\[ n(S \cup M) = 20 - 10 = 10. \]
By the principle of inclusion--exclusion (De Morgan's and addition laws):
\[ n(S \cap M) = n(S) + n(M) - n(S \cup M) = 5 + 8 - 10 = 3. \]
Thus, $n(S \cap M^c) = n(S) - n(S \cap M) = 5 - 3 = 2$.
\begin{enumerate}[(i)]
\item \textbf{Probability that a sunburnt student was ignored by mosquitoes ($P(M^c \mid S)$):}
\[ P(M^c \mid S) = \frac{P(S \cap M^c)}{P(S)} = \frac{n(S \cap M^c)}{n(S)} = \frac{2}{5} = 0.40. \]
\item \textbf{Probability that a mosquito-bitten student also got sunburnt ($P(S \mid M)$):}
\[ P(S \mid M) = \frac{P(S \cap M)}{P(M)} = \frac{n(S \cap M)}{n(M)} = \frac{3}{8} = 0.375. \]
\end{enumerate}
\ans{(i) $P(M^c \mid S) = \dfrac{2}{5} = 0.40$, \qquad (ii) $P(S \mid M) = \dfrac{3}{8} = 0.375$.}
\end{solution}

%---------------------------------------------------------------- P013
\begin{question}{2015--2016 \textbullet\ Q8(a)}{MATH 241 \textbullet\ 07/08/2016 \textbullet\ P013}{13}
The probability that a married man watches a certain television show is 0.25 and that a married woman watches the show is 0.35. The probability that a man watches the show, given that his wife does, is 0.8. Find the probability that:
(i) a married couple watches the show.
(ii) a husband watches the show given that his wife does not.
\end{question}
\begin{solution}
Let $H$ be the event that the husband watches the show and $W$ be the event that the wife watches the show.
\begin{note}
The printed data are internally inconsistent. If $P(H)=0.25$, $P(W)=0.35$, and $P(H \mid W)=0.8$ are used literally, then $P(H \cap W) = P(W)P(H \mid W) = 0.35 \times 0.8 = 0.28 > P(H) = 0.25$, which makes $P(H \cap W^c) = P(H) - P(H \cap W) = -0.03 < 0$. (A conditional probability $P(H\mid W)=0.8$ requires $P(H)\ge P(W)P(H\mid W)$.) Below, we show the literal substitution and the two most plausible consistent readings: the marginals swapped ($P(H)=0.35$, $P(W)=0.25$), or the $0.8$ meant as $P(W\mid H)$.
\end{note}
\textbf{Case A: Literal Formula Evaluation as Printed ($P(H)=0.25,\ P(W)=0.35,\ P(H \mid W)=0.8$).}
\begin{enumerate}[(i)]
\item \textbf{Probability that both watch ($P(H \cap W)$):}
\[ P(H \cap W) = P(W)\,P(H \mid W) = 0.35 \times 0.80 = 0.28. \]
\item \textbf{Probability that the husband watches given his wife does not ($P(H \mid W^c)$):}
\[ \begin{aligned} P(H \mid W^c) &= \frac{P(H \cap W^c)}{P(W^c)} = \frac{P(H) - P(H \cap W)}{1 - P(W)} \\ &= \frac{0.25 - 0.28}{1 - 0.35} = \frac{-0.03}{0.65} = -0.0462 \quad (\text{invalid probability due to typo}). \end{aligned} \]
\end{enumerate}

\textbf{Case B: Consistent Reading 1 (Swapped marginal labels: $P(H)=0.35,\ P(W)=0.25,\ P(H \mid W)=0.8$).}
\begin{enumerate}[(i)]
\item $P(H \cap W) = P(W)\,P(H \mid W) = 0.25 \times 0.80 = 0.20$.
\item $P(H \mid W^c) = \dfrac{P(H) - P(H \cap W)}{1 - P(W)} = \dfrac{0.35 - 0.20}{1 - 0.25} = \dfrac{0.15}{0.75} = 0.20$.
\end{enumerate}

\textbf{Case C: Consistent Reading 2 ($P(H)=0.25,\ P(W)=0.35$, with $0.8$ being $P(W \mid H)=0.8$).}
\begin{enumerate}[(i)]
\item $P(H \cap W) = P(H)\,P(W \mid H) = 0.25 \times 0.80 = 0.20$.
\item $P(H \mid W^c) = \dfrac{P(H) - P(H \cap W)}{1 - P(W)} = \dfrac{0.25 - 0.20}{1 - 0.35} = \dfrac{0.05}{0.65} = \dfrac{1}{13} \approx 0.0769$.
\end{enumerate}
\ans{Literal formula gives (i) $0.28$; under the corrected consistent parameters, (i) $P(H\cap W)=0.20$ and (ii) $P(H\mid W^c) = 0.20$ (or $1/13 \approx 0.0769$).}
\end{solution}

%===================================================================== TOPIC 03
\topic[Random Variables \& Distributions]{03}{Random Variables, PMF, PDF \& Cumulative Distribution Function (CDF)}

\begin{note}
This section covers the syllabus topics \emph{Random variable and types of Random Variable}, \emph{Discrete and Continuous Probability Distribution}, and \emph{Discrete Distribution Function (CDF) and Continuous Distribution Function (pdf)}.
\end{note}

%---------------------------------------------------------------- P014
\begin{question}{2023--2024 \textbullet\ Q2(a)}{MATH 243 \textbullet\ 21/05/2026 \textbullet\ P014}{15}
The waiting time, in hours, between successive speeders spotted by a radar unit is a continuous random variable $X$ with cumulative distribution function:
\[ F(x) = \begin{cases} 1 - e^{-8x}, & x \ge 0 \\ 0, & \text{otherwise.} \end{cases} \]
Find the probability of waiting less than 12 minutes between successive speeders using: (i) the cumulative distribution function of $X$, and (ii) the probability density function of $X$.
\end{question}
\begin{solution}
First, convert the waiting time threshold of $12\text{ minutes}$ into hours so its units match $X$:
\[ x = \frac{12\text{ minutes}}{60\text{ minutes/hour}} = 0.2\text{ hours}. \]
\begin{enumerate}[(i)]
\item \textbf{Using the Cumulative Distribution Function $F(x)$:}
Since $X$ is a continuous random variable, $P(X < 0.2) = P(X \le 0.2) = F(0.2)$:
\[ P(X < 0.2) = F(0.2) = 1 - e^{-8(0.2)} = 1 - e^{-1.6} \approx 1 - 0.201897 = 0.798103. \]
\item \textbf{Using the Probability Density Function $f(x)$:}
Differentiating $F(x)$ with respect to $x$ yields the probability density function:
\[ f(x) = \frac{d}{dx}F(x) = \begin{cases} 8e^{-8x}, & x \ge 0 \\ 0, & \text{otherwise.} \end{cases} \]
Integrating $f(x)$ over the interval $[0, 0.2]$:
\[ P(X < 0.2) = \int_{0}^{0.2} 8e^{-8x}\,dx = \Big[-e^{-8x}\Big]_{0}^{0.2} = 1 - e^{-1.6} \approx 0.798103. \]
\end{enumerate}
\ans{$P(\text{waiting less than 12 minutes}) = 1 - e^{-1.6} \approx 0.7981$ (by both CDF and PDF methods).}
\end{solution}

%---------------------------------------------------------------- P015
\begin{question}{2022--2023 \textbullet\ Q1(c)}{MATH 243 \textbullet\ 01/02/2025 \textbullet\ P015}{12}
The time to failure in hours of an electronic component in a manufactured DVD player has the density function:
\[ f(x) = \begin{cases} \dfrac{1}{2000}\exp\left(-\dfrac{x}{2000}\right), & \text{when } x \ge 0 \\[6pt] 0, & \text{when } x < 0 \end{cases} \]
(i) Find the cumulative distribution function for $f(x)$.
(ii) Determine the probability that the component (and thus the DVD player) lasts more than 1000 hours before the component needs to be replaced.
(iii) Determine the probability that the component fails before 2000 hours.
\end{question}
\begin{solution}
\begin{enumerate}[(i)]
\item \textbf{Cumulative Distribution Function $F(x)$:}
For $x < 0$, $F(x) = \int_{-\infty}^{x} 0\,dt = 0$. For $x \ge 0$:
\[ F(x) = \int_{-\infty}^{x} f(t)\,dt = \int_{0}^{x} \frac{1}{2000}e^{-t/2000}\,dt = \Big[-e^{-t/2000}\Big]_{0}^{x} = 1 - e^{-x/2000}. \]
Thus, the complete piecewise CDF is:
\[ F(x) = \begin{cases} 1 - e^{-x/2000}, & x \ge 0 \\ 0, & x < 0. \end{cases} \]
\item \textbf{Probability that the component lasts more than 1000 hours ($P(X > 1000)$):}
\[ P(X > 1000) = 1 - F(1000) = 1 - \left(1 - e^{-1000/2000}\right) = e^{-0.5} \approx 0.60653. \]
\item \textbf{Probability that the component fails before 2000 hours ($P(X < 2000)$):}
\[ P(X < 2000) = F(2000) = 1 - e^{-2000/2000} = 1 - e^{-1} \approx 1 - 0.36788 = 0.63212. \]
\end{enumerate}
\ans{(i) $F(x) = 1 - e^{-x/2000}$ for $x \ge 0$ ($0$ for $x < 0$); \quad (ii) $e^{-0.5} \approx 0.6065$; \quad (iii) $1 - e^{-1} \approx 0.6321$.}
\end{solution}

%---------------------------------------------------------------- P016
\begin{question}{2021--2022 \textbullet\ Q6(c)}{MATH 245 \textbullet\ 30/04/2023 \textbullet\ P016}{15}
An overseas shipment of 5 automobiles contains 2 that have slight paint blemishes. If an agency receives 3 automobiles at random, find the probability distribution and $F(x)$ of the random variable $X$ representing the number of automobiles purchased by the agency with paint blemishes.
\end{question}
\begin{solution}
Out of $N = 5$ automobiles, $K = 2$ have paint blemishes and $N - K = 3$ are blemish-free. The agency receives $n = 3$ automobiles. The random variable $X$ (number of blemished automobiles received) can take values $x \in \{0, 1, 2\}$, with hypergeometric probability mass function:
\[ P(X = x) = f(x) = \frac{\binom{2}{x}\binom{3}{3-x}}{\binom{5}{3}}, \qquad \text{where } \binom{5}{3} = 10. \]
Evaluating $f(x)$ for each possible value of $x$:
\begin{itemize}
\item For $x = 0$: \quad $f(0) = \dfrac{\binom{2}{0}\binom{3}{3}}{10} = \dfrac{1 \times 1}{10} = \dfrac{1}{10} = 0.10$.
\item For $x = 1$: \quad $f(1) = \dfrac{\binom{2}{1}\binom{3}{2}}{10} = \dfrac{2 \times 3}{10} = \dfrac{6}{10} = 0.60$.
\item For $x = 2$: \quad $f(2) = \dfrac{\binom{2}{2}\binom{3}{1}}{10} = \dfrac{1 \times 3}{10} = \dfrac{3}{10} = 0.30$.
\end{itemize}
\begin{center}
\small
\begin{tabular}{@{}cccc@{}}
\toprule
$x$ & $0$ & $1$ & $2$ \\
\midrule
$f(x) = P(X = x)$ & $0.10$ & $0.60$ & $0.30$ \\
$F(x) = P(X \le x)$ & $0.10$ & $0.70$ & $1.00$ \\
\bottomrule
\end{tabular}
\end{center}
Summing $f(t)$ for $t \le x$ gives the step-function cumulative distribution function $F(x)$:
\[ F(x) = P(X \le x) = \begin{cases} 0, & x < 0 \\ 0.10, & 0 \le x < 1 \\ 0.70, & 1 \le x < 2 \\ 1.00, & x \ge 2. \end{cases} \]
\ans{$f(0)=0.1,\ f(1)=0.6,\ f(2)=0.3$; \quad $F(x) = 0$ ($x<0$), $0.1$ ($0\le x<1$), $0.7$ ($1\le x<2$), $1.0$ ($x\ge 2$).}
\end{solution}

%---------------------------------------------------------------- P017
\begin{question}{2020--2021 \textbullet\ Q6(b)}{MATH 245 \textbullet\ 2020--2021 \textbullet\ P017}{10}
Suppose the measurement error $X$ of a certain physical quantity is decided by the density function:
\[ f(x) = \begin{cases} k(3-x^2), & -1 \le x \le 1, \\ 0, & \text{elsewhere.} \end{cases} \]
(i) Determine $k$ that renders $f(x)$ a valid density function.
(ii) Find the probability that a random error in measurement is less than $\frac{1}{2}$.
(iii) For this particular measurement, it is undesirable if the magnitude of the error exceeds 0.8. What is the probability that this occurs?
\end{question}
\begin{solution}
\begin{enumerate}[(i)]
\item \textbf{Finding $k$:} For $f(x)$ to be a valid probability density function, $f(x) \ge 0$ and $\int_{-\infty}^{\infty} f(x)\,dx = 1$. Since $3 - x^2$ is even on $[-1, 1]$:
\[ \int_{-1}^{1} k(3 - x^2)\,dx = 2k\int_{0}^{1} (3 - x^2)\,dx = 2k\left[3x - \frac{x^3}{3}\right]_{0}^{1} = 2k\left(3 - \frac{1}{3}\right) = \frac{16k}{3} = 1 \implies k = \frac{3}{16}. \]
\item \textbf{Probability that error is less than $\frac{1}{2}$ ($P(X < \frac{1}{2})$):}
\[ \begin{aligned} P\left(X < \frac{1}{2}\right) &= \frac{3}{16}\int_{-1}^{1/2} (3 - x^2)\,dx = \frac{3}{16}\left[3x - \frac{x^3}{3}\right]_{-1}^{1/2} = \frac{3}{16}\left[\left(\frac{3}{2} - \frac{1}{24}\right) - \left(-3 + \frac{1}{3}\right)\right] \\ &= \frac{3}{16}\left(\frac{35}{24} + \frac{8}{3}\right) = \frac{3}{16}\times\frac{99}{24} = \frac{99}{128} \approx 0.77344. \end{aligned} \]
\item \textbf{Probability that magnitude of error exceeds $0.8$ ($P(|X| > 0.8)$):}
By symmetry of $f(x)$ about $x = 0$:
\[ \begin{aligned} P(|X| > 0.8) &= 2\int_{0.8}^{1} \frac{3}{16}(3 - x^2)\,dx = \frac{3}{8}\left[3x - \frac{x^3}{3}\right]_{0.8}^{1} \\ &= \frac{3}{8}\left[\frac{8}{3} - \left(2.4 - \frac{0.512}{3}\right)\right] = \frac{3}{8}\times\frac{1.312}{3} = 0.164 = \frac{41}{250}. \end{aligned} \]
\end{enumerate}
\ans{(i) $k = \dfrac{3}{16}$, \quad (ii) $P(X < \tfrac{1}{2}) = \dfrac{99}{128} \approx 0.7734$, \quad (iii) $P(|X| > 0.8) = \dfrac{41}{250} = 0.164$.}
\end{solution}

%---------------------------------------------------------------- P018
\begin{question}{2018--2019 \textbullet\ Q7(a)}{MATH 245 \textbullet\ 03/11/2019 \textbullet\ P018}{20}
What is density function? What condition must a density function satisfy? Suppose that the distribution function of a random variable $X$ is as follows:
\[ F(x) = \begin{cases} 0, & x \le 0 \\ \dfrac{x}{8}, & 0 < x < 2 \\[4pt] \dfrac{y^2}{16} \ \text{[sic: }\dfrac{x^2}{16}\text{]}, & 2 \le x < 4 \\[4pt] 1, & x \ge 4 \end{cases} \]
(i) Find the probability density function of $X$.
(ii) Find $p(x \ge 1.5)$ and $p(x \ge 1 \mid x \le 3)$.
\end{question}
\begin{solution}
\begin{note}
In the printed exam paper, the third branch of $F(x)$ is misprinted as $\frac{y^2}{16}$ instead of $\frac{x^2}{16}$. With $\frac{x^2}{16}$, $F(x)$ is continuous everywhere ($F(2^-)=\frac{2}{8}=\frac{1}{4}=F(2^+)$ and $F(4^-)=\frac{16}{16}=1=F(4^+)$).
\end{note}
\textbf{Definition \& Conditions of a Probability Density Function (pdf):}
For a continuous random variable $X$, the probability density function $f(x)$ is the function whose integral over any interval $[a,b]$ gives the probability $P(a \le X \le b) = \int_a^b f(x)\,dx$. A valid pdf must satisfy two conditions:
\begin{enumerate}[1.]
\item \textbf{Non-negativity:} $f(x) \ge 0$ for all $x \in (-\infty, \infty)$.
\item \textbf{Normalization (Unit Total Area):} $\displaystyle\int_{-\infty}^{\infty} f(x)\,dx = 1$.
\end{enumerate}

\medskip
\textbf{(i) Probability Density Function of $X$:}
Differentiating $F(x)$ piecewise with respect to $x$ ($f(x) = F'(x)$):
\[ f(x) = \begin{cases} \dfrac{1}{8}, & 0 < x < 2 \\[6pt] \dfrac{x}{8}, & 2 \le x < 4 \\[6pt] 0, & \text{elsewhere.} \end{cases} \]
\textit{(Check: $\int_0^2 \frac{1}{8}\,dx + \int_2^4 \frac{x}{8}\,dx = \frac{2}{8} + \frac{16-4}{16} = \frac{1}{4} + \frac{3}{4} = 1$. \ok)}

\medskip
\textbf{(ii) Evaluating $P(X \ge 1.5)$ and $P(X \ge 1 \mid X \le 3)$:}
\begin{itemize}
\item Since $X$ is continuous, $P(X \ge 1.5) = 1 - F(1.5)$. Because $1.5 \in (0, 2)$, $F(1.5) = \dfrac{1.5}{8} = \dfrac{3}{16}$:
\[ P(X \ge 1.5) = 1 - \frac{3}{16} = \frac{13}{16} = 0.8125. \]
\item By the definition of conditional probability:
\[ P(X \ge 1 \mid X \le 3) = \frac{P(1 \le X \le 3)}{P(X \le 3)} = \frac{F(3) - F(1)}{F(3)}. \]
Since $1 \in (0, 2)$ and $3 \in [2, 4)$, we have $F(1) = \dfrac{1}{8} = \dfrac{2}{16}$ and $F(3) = \dfrac{3^2}{16} = \dfrac{9}{16}$. Hence:
\[ P(X \ge 1 \mid X \le 3) = \frac{\frac{9}{16} - \frac{2}{16}}{\frac{9}{16}} = \frac{7/16}{9/16} = \frac{7}{9} \approx 0.7778. \]
\end{itemize}
\ans{(i) $f(x) = \frac{1}{8}$ ($0<x<2$), $\frac{x}{8}$ ($2\le x<4$), $0$ elsewhere; \quad (ii) $P(X\ge 1.5) = \dfrac{13}{16} = 0.8125$ and $P(X\ge 1\mid X\le 3) = \dfrac{7}{9} \approx 0.7778$.}
\end{solution}

%===================================================================== TOPIC 04
\topic[Joint, Marginal \& Conditional]{04}{Joint, Marginal \& Conditional Distributions}

\begin{note}
This section covers the syllabus topics \emph{Joint Probability Distribution (Discrete, Continuous)}, \emph{Marginal Distribution}, and \emph{Conditional Distribution and Independence of Random Variable}.
\end{note}

%---------------------------------------------------------------- P019
\begin{question}{2023--2024 \textbullet\ Q2(b)}{MATH 243 \textbullet\ 21/05/2026 \textbullet\ P019}{20}
From a bag of pens containing three red pens, two blue pens and three black pens, a random sample of two pens is selected. If $X$ is the number of red pens, and $Y$ is the number of blue pens in the sample, then work on the following:
(i) Plot $f(x,y)$ versus $(x,y)$, where $f(x,y)$ denotes joint distribution.
(ii) $P[(X,Y) \in A]$, where $A = \{(x,y) \text{ such that } x + y \le 1\}$.
(iii) Marginal distributions of $X$ and $Y$.
(iv) Variance of $X$.
(v) Covariance between $X$ and $Y$.
\end{question}
\begin{solution}
The bag contains $3\text{ red}$, $2\text{ blue}$, and $3\text{ black}$ pens ($8\text{ pens}$ total). A sample of $2$ pens is drawn without replacement, so the total number of outcomes is $\binom{8}{2} = 28$. The joint probability mass function is:
\[ f(x,y) = P(X=x, Y=y) = \frac{\binom{3}{x}\binom{2}{y}\binom{3}{2-x-y}}{\binom{8}{2}}, \qquad x \ge 0,\ y \ge 0,\ x + y \le 2. \]
Evaluating $f(x,y)$ for all six valid pairs $(x,y)$:
\begin{itemize}
\item $f(0,0) = \dfrac{\binom{3}{0}\binom{2}{0}\binom{3}{2}}{28} = \dfrac{3}{28}$, \qquad $f(0,1) = \dfrac{\binom{3}{0}\binom{2}{1}\binom{3}{1}}{28} = \dfrac{6}{28} = \dfrac{3}{14}$, \qquad $f(0,2) = \dfrac{\binom{3}{0}\binom{2}{2}\binom{3}{0}}{28} = \dfrac{1}{28}$,
\item $f(1,0) = \dfrac{\binom{3}{1}\binom{2}{0}\binom{3}{1}}{28} = \dfrac{9}{28}$, \qquad $f(1,1) = \dfrac{\binom{3}{1}\binom{2}{1}\binom{3}{0}}{28} = \dfrac{6}{28} = \dfrac{3}{14}$, \qquad $f(2,0) = \dfrac{\binom{3}{2}\binom{2}{0}\binom{3}{0}}{28} = \dfrac{3}{28}$.
\end{itemize}
\begin{center}
\small
\begin{tabular}{@{}ccccc@{}}
\toprule
$f(x,y)$ & $y = 0$ & $y = 1$ & $y = 2$ & \textbf{Marginal $g(x) = P(X=x)$} \\
\midrule
$x = 0$ & $3/28$ & $6/28$ & $1/28$ & $10/28 = 5/14$ \\
$x = 1$ & $9/28$ & $6/28$ & $0$ & $15/28$ \\
$x = 2$ & $3/28$ & $0$ & $0$ & $3/28$ \\
\midrule
\textbf{Marginal $h(y) = P(Y=y)$} & $15/28$ & $12/28 = 3/7$ & $1/28$ & $\mathbf{1}$ \\
\bottomrule
\end{tabular}
\end{center}

\textbf{(i) Plot of $f(x,y)$ versus $(x,y)$:}
\fig{\jointpmffig}

\textbf{(ii) Probability $P[(X,Y) \in A]$ where $A = \{(x,y) : x + y \le 1\}$:}
The points in $A$ are $(0,0), (0,1),$ and $(1,0)$:
\[ P[(X,Y) \in A] = f(0,0) + f(0,1) + f(1,0) = \frac{3}{28} + \frac{6}{28} + \frac{9}{28} = \frac{18}{28} = \frac{9}{14} \approx 0.6429. \]

\textbf{(iii) Marginal Distributions of $X$ and $Y$:}
From the row and column totals of the joint probability table above:
\begin{itemize}
\item Marginal PMF of $X$: \quad $g(0) = \dfrac{5}{14}$, \quad $g(1) = \dfrac{15}{28}$, \quad $g(2) = \dfrac{3}{28}$.
\item Marginal PMF of $Y$: \quad $h(0) = \dfrac{15}{28}$, \quad $h(1) = \dfrac{3}{7}$, \quad $h(2) = \dfrac{1}{28}$.
\end{itemize}

\textbf{(iv) Variance of $X$:}
\[ E(X) = \sum_{x=0}^{2} x\,g(x) = 0\left(\frac{10}{28}\right) + 1\left(\frac{15}{28}\right) + 2\left(\frac{3}{28}\right) = \frac{21}{28} = \frac{3}{4} = 0.75, \]
\[ E(X^2) = \sum_{x=0}^{2} x^2\,g(x) = 0 + 1^2\left(\frac{15}{28}\right) + 2^2\left(\frac{3}{28}\right) = \frac{27}{28}, \]
\[ \operatorname{Var}(X) = E(X^2) - [E(X)]^2 = \frac{27}{28} - \left(\frac{3}{4}\right)^2 = \frac{27}{28} - \frac{9}{16} = \frac{108 - 63}{112} = \frac{45}{112} \approx 0.4018. \]

\textbf{(v) Covariance between $X$ and $Y$:}
First compute $E(Y)$ and $E(XY)$:
\[ E(Y) = \sum_{y=0}^{2} y\,h(y) = 0 + 1\left(\frac{12}{28}\right) + 2\left(\frac{1}{28}\right) = \frac{14}{28} = \frac{1}{2} = 0.50, \]
\[ E(XY) = \sum_{x}\sum_{y} xy\,f(x,y) = (1)(1)f(1,1) = \frac{6}{28} = \frac{3}{14} \quad (\text{other terms vanish}). \]
Therefore,
\[ \operatorname{Cov}(X,Y) = E(XY) - E(X)E(Y) = \frac{3}{14} - \left(\frac{3}{4}\right)\left(\frac{1}{2}\right) = \frac{12 - 21}{56} = -\frac{9}{56} \approx -0.1607. \]
\ans{(ii) $P(A) = \dfrac{9}{14}$; \quad (iii) $g(x) = \left\{\frac{5}{14}, \frac{15}{28}, \frac{3}{28}\right\}$, $h(y) = \left\{\frac{15}{28}, \frac{3}{7}, \frac{1}{28}\right\}$; \quad (iv) $\operatorname{Var}(X) = \dfrac{45}{112}$; \quad (v) $\operatorname{Cov}(X,Y) = -\dfrac{9}{56}$.}
\end{solution}

%---------------------------------------------------------------- P020
\begin{question}{2022--2023 \textbullet\ Q2(a)}{MATH 243 \textbullet\ 01/02/2025 \textbullet\ P020}{15}
The joint density for the random variables $(X, Y)$, where $X$ is the unit temperature change and $Y$ is the proportion of spectrum shift that a certain atomic particle produces, is:
\[ f(x, y) = \begin{cases} 10xy^2, & 0 < x < y < 1 \\ 0, & \text{elsewhere.} \end{cases} \]
(i) Find marginal densities $g(x), h(y)$, and conditional density $f(y|x)$.
(ii) Find the probability that the spectrum shifts more than half of the total observations, given that the temperature is increased by 0.25 unit.
\end{question}
\begin{solution}
\begin{enumerate}[(i)]
\item \textbf{Marginal and Conditional Densities:}
The support is the triangular region $0 < x < y < 1$.
\begin{itemize}
\item \textbf{Marginal density of $X$, $g(x)$:} For a fixed $x \in (0, 1)$, $y$ ranges from $x$ to $1$:
\[ g(x) = \int_{x}^{1} 10xy^2\,dy = 10x\left[\frac{y^3}{3}\right]_{x}^{1} = \frac{10}{3}x(1 - x^3), \quad 0 < x < 1 \ (0\text{ elsewhere}). \]
\item \textbf{Marginal density of $Y$, $h(y)$:} For a fixed $y \in (0, 1)$, $x$ ranges from $0$ to $y$:
\[ h(y) = \int_{0}^{y} 10xy^2\,dx = 10y^2\left[\frac{x^2}{2}\right]_{0}^{y} = 5y^4, \quad 0 < y < 1 \ (0\text{ elsewhere}). \]
\item \textbf{Conditional density of $Y$ given $X = x$, $f(y \mid x)$:} For $0 < x < y < 1$:
\[ f(y \mid x) = \frac{f(x,y)}{g(x)} = \frac{10xy^2}{\frac{10}{3}x(1 - x^3)} = \frac{3y^2}{1 - x^3}, \qquad x < y < 1 \quad (0\text{ elsewhere}). \]
\end{itemize}
\item \textbf{Conditional Probability $P(Y > 0.5 \mid X = 0.25)$:}
Substituting $x = 0.25 = \frac{1}{4}$ into $f(y \mid x)$ and integrating over $y \in (0.5, 1)$:
\[ \begin{aligned} P(Y > 0.5 \mid X = 0.25) &= \int_{1/2}^{1} \frac{3y^2}{1 - (1/4)^3}\,dy = \frac{1}{63/64}\Big[y^3\Big]_{1/2}^{1} \\ &= \frac{7/8}{63/64} = \frac{8}{9} \approx 0.8889. \end{aligned} \]
\end{enumerate}
\ans{(i) $g(x)=\frac{10}{3}x(1-x^3)$ ($0<x<1$), $h(y)=5y^4$ ($0<y<1$), $f(y\mid x)=\dfrac{3y^2}{1-x^3}$ ($x<y<1$); \quad (ii) $P(Y>0.5\mid X=0.25) = \dfrac{8}{9} \approx 0.8889$.}
\end{solution}

%---------------------------------------------------------------- P021
\begin{question}{2020--2021 \textbullet\ Q5(c)}{MATH 245 \textbullet\ 2020--2021 \textbullet\ P021}{10}
Each rear tire on an experimental airplane is supposed to be filled to a pressure of 40 pounds per square inch (psi). Let $X$ denote the actual air pressure for the right tire and $Y$ denote the actual air pressure for the left tire. Suppose that $X$ and $Y$ are random variables with the joint density function:
\[ f(x,y) = \begin{cases} k(x^2+y^2), & 30 \le x < 50,\ 30 \le y < 50, \\ 0, & \text{elsewhere.} \end{cases} \]
(i) Find the value of $k$.
(ii) Find $P(30 \le X \le 40 \text{ and } 40 \le Y < 50)$.
(iii) Find the probability that both tires are underfilled.
\end{question}
\begin{solution}
\begin{enumerate}[(i)]
\item \textbf{Finding $k$:} Setting the total volume under $f(x,y)$ over $[30,50]\times[30,50]$ equal to $1$:
\[ \begin{aligned}
\int_{30}^{50}\!\int_{30}^{50} k(x^2 + y^2)\,dx\,dy &= k\int_{30}^{50} \left(\frac{98{,}000}{3} + 20y^2\right) dy \\
&= k\left[\frac{98{,}000}{3}(20) + 20\left(\frac{98{,}000}{3}\right)\right] = \frac{3{,}920{,}000}{3}k = 1 \\
\implies k &= \frac{3}{3{,}920{,}000}.
\end{aligned} \]
\item \textbf{Probability $P(30 \le X \le 40 \text{ and } 40 \le Y < 50)$:}
\[ \begin{aligned}
P(30 \le X \le 40,\ 40 \le Y < 50) &= k\int_{40}^{50}\!\left(\frac{37{,}000}{3} + 10y^2\right) dy \\
&= \frac{3}{3{,}920{,}000}\times\frac{370{,}000 + 610{,}000}{3} = \frac{1}{4} = 0.25.
\end{aligned} \]
\item \textbf{Probability that both tires are underfilled ($X < 40$ and $Y < 40$):}
Since the target pressure is $40\text{ psi}$, both tires are underfilled when $30 \le X < 40$ and $30 \le Y < 40$:
\[ \begin{aligned} P(30 \le X < 40,\ 30 \le Y < 40) &= k\int_{30}^{40}\!\int_{30}^{40} (x^2 + y^2)\,dx\,dy = k\left[20 \times \frac{40^3 - 30^3}{3}\right] \\ &= \frac{3}{3{,}920{,}000}\times\frac{740{,}000}{3} = \frac{37}{196} \approx 0.18878. \end{aligned} \]
\end{enumerate}
\ans{(i) $k = \dfrac{3}{3{,}920{,}000}$, \quad (ii) $P(30\le X\le 40,\ 40\le Y<50) = \dfrac{1}{4} = 0.25$, \quad (iii) $P(\text{both underfilled}) = \dfrac{37}{196} \approx 0.1888$.}
\end{solution}

%---------------------------------------------------------------- P022
\begin{question}{2012--2013 \textbullet\ Q8(a)}{MATH 241 \textbullet\ 02/06/2014 \textbullet\ P022}{18}
For the joint density function
\[ f(x,y) = \begin{cases} \dfrac{6-x-y}{3}, & 0 < x < 2,\ 2 < y < 4 \\ 0, & \text{elsewhere,} \end{cases} \]
find $P(1 < Y < 3 \mid x = 2)$.
\end{question}
\begin{solution}
\begin{note}
\textbf{Correction (checked by Claude):} with the printed constant $\tfrac{1}{3}$, $\int_0^2\!\int_2^4 \frac{6-x-y}{3}\,dy\,dx = \frac{8}{3} \neq 1$, so it is not a valid density. The correct constant is $\tfrac{1}{8}$, i.e.\ $f(x,y) = \frac{6-x-y}{8}$ (this is the standard textbook version). The constant cancels in $f(y\mid x)$, so the final answer is unaffected; only $g(x)$ changes. Note also that $x=2$ lies on the boundary of $0<x<2$, so $f(y\mid x=2)$ is taken as the limit $x\to 2^-$.
\end{note}
\textbf{Step 1: Find the marginal density $g(x)$ of $X$.}
For $0 < x < 2$, integrate $f(x,y) = \frac{6-x-y}{8}$ with respect to $y$ from $2$ to $4$:
\[ g(x) = \int_{2}^{4} \frac{6 - x - y}{8}\,dy = \frac{1}{8}\left[(6 - x)y - \frac{y^2}{2}\right]_{2}^{4} = \frac{2(6 - x) - 6}{8} = \frac{6 - 2x}{8} = \frac{3 - x}{4}. \]
\textit{(Check: $\int_0^2 \frac{3-x}{4}\,dx = \frac{6-2}{4} = 1$. \ok)}
\textbf{Step 2: Find the conditional density $f(y \mid x)$.}
For $2 < y < 4$:
\[ f(y \mid x) = \frac{f(x,y)}{g(x)} = \frac{(6 - x - y)/8}{(6 - 2x)/8} = \frac{6 - x - y}{6 - 2x}, \qquad 2 < y < 4. \]
Evaluating at the boundary $x = 2$ (by continuity of the formula as $x \to 2^-$):
\[ f(y \mid x = 2) = \frac{6 - 2 - y}{6 - 4} = \frac{4 - y}{2}, \qquad 2 < y < 4 \quad (\text{and } 0 \text{ elsewhere}). \]
\textbf{Step 3: Compute $P(1 < Y < 3 \mid x = 2)$.}
Since $f(y \mid x=2) = 0$ for $1 < y \le 2$, the integral over $(1, 3)$ reduces to $(2, 3)$:
\[ P(1 < Y < 3 \mid x = 2) = \int_{2}^{3} \frac{4 - y}{2}\,dy = \frac{1}{2}\left[4y - \frac{y^2}{2}\right]_{2}^{3} = \frac{1}{2}\left(4 - \frac{5}{2}\right) = \frac{3}{4} = 0.75. \]
\ans{$P(1 < Y < 3 \mid x = 2) = \dfrac{3}{4} = 0.75$.}
\end{solution}

%===================================================================== TOPIC 05
\topic[Expectation, Variance \& Covariance]{05}{Mathematical Expectation, Variance \& Covariance}

\begin{note}
For the variance and covariance of discrete bivariate random variables ($\operatorname{Var}(X)$ and $\operatorname{Cov}(X,Y)$), see also \textbf{2023--2024 MATH 243 Q2(b)(iv)--(v) [P019]} under Topic 04.
\end{note}

%---------------------------------------------------------------- P023
\begin{question}{2022--2023 \textbullet\ Q2(b)}{MATH 243 \textbullet\ 01/02/2025 \textbullet\ P023}{10}
The length of time, in minutes, for an airplane to obtain clearance for takeoff at a certain airport is a random variable $Y = 3X - 2$, where $X$ has the density function:
\[ f(x) = \begin{cases} \dfrac{1}{4}\exp\left(-\dfrac{x}{4}\right), & x > 0 \\ 0, & \text{elsewhere.} \end{cases} \]
Find the mean and variance of the random variable $Y$.
\end{question}
\begin{solution}
First, compute $E(X)$ and $E(X^2)$ using integration by parts (or the gamma integral $\int_0^\infty x^n e^{-x/\beta}\,dx = n!\,\beta^{n+1}$ with $\beta = 4$):
\begin{align*}
E(X) &= \int_{0}^{\infty} x\cdot\frac{1}{4}e^{-x/4}\,dx = \frac{1}{4}\left(1!\times 4^2\right) = 4, \\[4pt]
E(X^2) &= \int_{0}^{\infty} x^2\cdot\frac{1}{4}e^{-x/4}\,dx = \frac{1}{4}\left(2!\times 4^3\right) = 32.
\end{align*}
Hence, the variance of $X$ is:
\[ \operatorname{Var}(X) = E(X^2) - [E(X)]^2 = 32 - 4^2 = 16. \]
Using the linear properties of expectation and variance for $Y = 3X - 2$:
\begin{align*}
E(Y) &= E(3X - 2) = 3E(X) - 2 = 3(4) - 2 = 10\text{ minutes}, \\[4pt]
\operatorname{Var}(Y) &= \operatorname{Var}(3X - 2) = 3^2\operatorname{Var}(X) = 9 \times 16 = 144\text{ minutes}^2.
\end{align*}
\ans{Mean $E(Y) = 10\text{ minutes}$, \qquad Variance $\operatorname{Var}(Y) = 144\text{ minutes}^2$.}
\end{solution}

%---------------------------------------------------------------- P024
\begin{question}{2020--2021 \textbullet\ Q8(a)}{MATH 245 \textbullet\ 2020--2021 \textbullet\ P024}{10}
The hospitalization period, in days, for patients following treatment for a certain type of kidney disorder is a random variable $Y = X + 4$, where $X$ has the density function:
\[ f(x) = \begin{cases} \dfrac{32}{(x+4)^3}, & x > 0, \\ 0, & \text{elsewhere.} \end{cases} \]
Find the average number of days that a person is hospitalized following treatment for this disorder.
\end{question}
\begin{solution}
The average number of days a patient is hospitalized is the mathematical expectation $E(Y) = E(X + 4)$:
\[ E(Y) = \int_{-\infty}^{\infty} (x + 4)f(x)\,dx = \int_{0}^{\infty} (x + 4)\cdot\frac{32}{(x + 4)^3}\,dx = \int_{0}^{\infty} \frac{32}{(x + 4)^2}\,dx. \]
Evaluating the improper integral:
\[ E(Y) = \left[-\frac{32}{x + 4}\right]_{0}^{\infty} = \lim_{b\to\infty}\left(-\frac{32}{b + 4}\right) - \left(-\frac{32}{0 + 4}\right) = 0 + \frac{32}{4} = 8\text{ days}. \]
\ans{The average hospitalization period is $E(Y) = 8\text{ days}$.}
\end{solution}

%---------------------------------------------------------------- P025
\begin{question}{2019--2020 \textbullet\ Q7(b)}{MATH 245 \textbullet\ 25/09/2021 \textbullet\ P025}{18}
A continuous random variable $X$ has the density function:
\[ f(x) = \begin{cases} e^{-x}, & x > 0 \\ 0, & \text{Otherwise.} \end{cases} \]
Find the expected value of $g(X) = e^{2X/3}$.
\end{question}
\begin{solution}
By the law of the unconscious statistician (expected value of a function of a continuous random variable):
\[ E[g(X)] = \int_{-\infty}^{\infty} g(x)f(x)\,dx = \int_{0}^{\infty} e^{2x/3}\cdot e^{-x}\,dx = \int_{0}^{\infty} e^{-x/3}\,dx. \]
Evaluating the improper integral:
\[ E\left[e^{2X/3}\right] = \Big[-3e^{-x/3}\Big]_{0}^{\infty} = \lim_{b\to\infty}\left(-3e^{-b/3}\right) - \left(-3e^{0}\right) = 0 + 3 = 3. \]
\ans{$E[g(X)] = E\left[e^{2X/3}\right] = 3$.}
\end{solution}

%---------------------------------------------------------------- P026
\begin{question}{2018--2019 \textbullet\ Q7(b)}{MATH 245 \textbullet\ 03/11/2019 \textbullet\ P026}{15}
A lot containing 7 components is sampled by a quality inspector; the lot contains 4 good components and 3 defective components. A sample of 3 is taken by the inspector. Find the expected value of the number of good components in this sample taken by the inspector.
\end{question}
\begin{solution}
Let $X$ denote the number of good components in the sample of $n = 3$ components drawn without replacement from $N = 7$ components ($K = 4$ good, $N - K = 3$ defective). Then $X$ has the hypergeometric PMF:
\[ P(X = x) = \frac{\binom{4}{x}\binom{3}{3-x}}{\binom{7}{3}}, \qquad x = 0, 1, 2, 3, \quad \text{where } \binom{7}{3} = 35. \]
Computing the probability for each $x$:
\begin{itemize}
\item $P(X = 0) = \dfrac{\binom{4}{0}\binom{3}{3}}{35} = \dfrac{1 \times 1}{35} = \dfrac{1}{35}$,
\item $P(X = 1) = \dfrac{\binom{4}{1}\binom{3}{2}}{35} = \dfrac{4 \times 3}{35} = \dfrac{12}{35}$,
\item $P(X = 2) = \dfrac{\binom{4}{2}\binom{3}{1}}{35} = \dfrac{6 \times 3}{35} = \dfrac{18}{35}$,
\item $P(X = 3) = \dfrac{\binom{4}{3}\binom{3}{0}}{35} = \dfrac{4 \times 1}{35} = \dfrac{4}{35}$.
\end{itemize}
By definition of mathematical expectation:
\[ E(X) = \sum_{x=0}^{3} x\,P(X = x) = 0 + 1\left(\frac{12}{35}\right) + 2\left(\frac{18}{35}\right) + 3\left(\frac{4}{35}\right) = \frac{60}{35} = \frac{12}{7} \approx 1.7143. \]
\textit{(Direct hypergeometric mean check: $E(X) = n\cdot\dfrac{K}{N} = 3 \times \dfrac{4}{7} = \dfrac{12}{7}$. \ok)}
\ans{The expected number of good components is $E(X) = \dfrac{12}{7} \approx 1.7143$.}
\end{solution}

%===================================================================== TOPIC 06
\topic{06}{Binomial Distribution}

%---------------------------------------------------------------- P027
\begin{question}{2023--2024 \textbullet\ Q3(c)}{MATH 243 \textbullet\ 21/05/2026 \textbullet\ P027}{12}
The probability that a person recovers from a rare blood disease is 0.4. If 15 people are known to have contracted this disease, what is the probability that:
(i) at least 10 people survive?
(ii) from 3 to 8 survive?
(iii) exactly 5 people die?
\end{question}
\begin{solution}
Let $X$ be the number of patients who survive out of $n = 15$ independent patients, where the probability of survival (recovery) is $p = 0.4$ and the probability of death is $q = 1 - p = 0.6$. Then $X \sim \operatorname{Binomial}(n = 15, p = 0.4)$, with PMF:
\[ P(X = x) = b(x; 15, 0.4) = \binom{15}{x}(0.4)^x(0.6)^{15-x}, \qquad x = 0, 1, \dots, 15. \]
\begin{enumerate}[(i)]
\item \textbf{At least 10 people survive ($P(X \ge 10)$):}
\begin{align*}
P(X \ge 10) &= \sum_{x=10}^{15} \binom{15}{x}(0.4)^x(0.6)^{15-x} \\
&= 0.024486 + 0.007420 + 0.001649 + 0.000254 + 0.000024 + 0.000001 = 0.033833.
\end{align*}
\item \textbf{From 3 to 8 survive ($P(3 \le X \le 8)$):}
\begin{align*}
P(3 \le X \le 8) &= \sum_{x=3}^{8} \binom{15}{x}(0.4)^x(0.6)^{15-x} \\
&= b(3) + b(4) + b(5) + b(6) + b(7) + b(8) \\
&= 0.063388 + 0.126776 + 0.185938 + 0.206598 + 0.177084 + 0.118056 \\
&= 0.877839.
\end{align*}
\item \textbf{Exactly 5 people die:}
If 5 out of 15 people die, then exactly $15 - 5 = 10$ people survive ($X = 10$):
\[ P(\text{5 die}) = P(X = 10) = \binom{15}{10}(0.4)^{10}(0.6)^{5} = 3003(0.4)^{10}(0.6)^5 \approx 0.024486. \]
\end{enumerate}
\ans{(i) $P(X \ge 10) = 0.0338$, \quad (ii) $P(3 \le X \le 8) = 0.8778$, \quad (iii) $P(\text{5 die}) = P(X=10) = 0.0245$.}
\end{solution}

%---------------------------------------------------------------- P028
\begin{question}{2020--2021 \textbullet\ Q7(a)}{MATH 245 \textbullet\ 2020--2021 \textbullet\ P028}{17}
For a particular type of circuit board, the probability of a diode failure is 0.03 and the board contains 200 diodes.
(i) What is the mean number of failures among the diodes?
(ii) What is the variance?
(iii) The board will work if there are no defective diodes. What is the probability that a board will work?
\end{question}
\begin{solution}
Let $X$ denote the number of failed diodes on a board of $n = 200$ independent diodes, each failing with probability $p = 0.03$ ($q = 1 - p = 0.97$). Thus $X \sim \operatorname{Binomial}(n = 200, p = 0.03)$.
\begin{enumerate}[(i)]
\item \textbf{Mean number of failures:}
\[ \mu = E(X) = np = 200 \times 0.03 = 6. \]
\item \textbf{Variance of the number of failures:}
\[ \sigma^2 = \operatorname{Var}(X) = npq = 200 \times 0.03 \times 0.97 = 5.82 \quad (\text{standard deviation } \sigma = \sqrt{5.82} \approx 2.4125). \]
\item \textbf{Probability that a board will work ($X = 0$):}
\begin{itemize}
\item \textbf{Exact Binomial Calculation:}
\[ P(X = 0) = \binom{200}{0}(0.03)^0(0.97)^{200} = (0.97)^{200} \approx 0.002261. \]
\item \textbf{Poisson Approximation (with $\lambda = np = 6$):}
\[ P(X = 0) \approx \frac{e^{-6}6^0}{0!} = e^{-6} \approx 0.002479. \]
\end{itemize}
\end{enumerate}
\ans{(i) $\mu = 6$, \quad (ii) $\sigma^2 = 5.82$, \quad (iii) $P(X=0) = (0.97)^{200} \approx 0.00226$ (or $e^{-6} \approx 0.00248$ via Poisson approximation).}
\end{solution}

%---------------------------------------------------------------- P029
\begin{question}{2014--2015 \textbullet\ Q8(b)}{MATH 241 \textbullet\ 06/02/2016 \textbullet\ P029}{12}
The probability that an entering college student will graduate is 0.4. Determine the probability that out of 5 students: (i) none, (ii) at least 1, (iii) all will graduate.
\end{question}
\begin{solution}
Let $X \sim \operatorname{Binomial}(n = 5, p = 0.4)$, where $q = 1 - 0.4 = 0.6$:
\begin{enumerate}[(i)]
\item \textbf{None will graduate ($X = 0$):}
\[ P(X = 0) = \binom{5}{0}(0.4)^0(0.6)^5 = (0.6)^5 = 0.07776. \]
\item \textbf{At least 1 will graduate ($X \ge 1$):}
\[ P(X \ge 1) = 1 - P(X = 0) = 1 - 0.07776 = 0.92224. \]
\item \textbf{All 5 will graduate ($X = 5$):}
\[ P(X = 5) = \binom{5}{5}(0.4)^5(0.6)^0 = (0.4)^5 = 0.01024. \]
\end{enumerate}
\ans{(i) $0.07776$, \quad (ii) $0.92224$, \quad (iii) $0.01024$.}
\end{solution}

%---------------------------------------------------------------- P030
\begin{question}{2013--2014 \textbullet\ Q7(b)(ii)}{MATH 241 \textbullet\ 13/12/2014 \textbullet\ P030}{8}
If 10\% of the instruments produced by a company are defective, find the probability that out of 12 instruments, at least 4 instruments will be defective.
\end{question}
\begin{solution}
Let $X$ be the number of defective instruments in a random sample of $n = 12$, where $p = 0.10$ and $q = 0.90$. Using the complement rule:
\[ P(X \ge 4) = 1 - P(X \le 3) = 1 - \sum_{x=0}^{3} \binom{12}{x}(0.10)^x(0.90)^{12-x}. \]
Computing the individual terms for $x = 0, 1, 2, 3$:
\begin{align*}
P(X = 0) &= \binom{12}{0}(0.1)^0(0.9)^{12} = 1 \times (0.9)^{12} \approx 0.282430, \\
P(X = 1) &= \binom{12}{1}(0.1)^1(0.9)^{11} = 12 \times 0.1 \times 0.313811 \approx 0.376573, \\
P(X = 2) &= \binom{12}{2}(0.1)^2(0.9)^{10} = 66 \times 0.01 \times 0.348678 \approx 0.230128, \\
P(X = 3) &= \binom{12}{3}(0.1)^3(0.9)^{9} = 220 \times 0.001 \times 0.387420 \approx 0.085233.
\end{align*}
Summing these gives $P(X \le 3) = 0.974363$, and thus:
\[ P(X \ge 4) = 1 - 0.974363 = 0.025637. \]
\ans{$P(X \ge 4) = 0.02564$.}
\end{solution}

%---------------------------------------------------------------- P031
\begin{question}{2012--2013 \textbullet\ Q6(b)}{MATH 241 \textbullet\ 02/06/2014 \textbullet\ P031}{10}
In testing a certain kind of truck tire over a rugged terrain, it is found that 20\% of the trucks fail to complete the test run without a blowout. Of the next 18 trucks tested, find the probability using (both binomial and Poisson distribution) that:
(i) from 4 to 7 have blowouts;
(ii) fewer than 3 have blowouts;
(iii) more than 5 have blowouts.
\end{question}
\begin{solution}
Here $n = 18$, $p = 0.20$, $q = 0.80$, and the Poisson parameter is $\lambda = np = 18 \times 0.20 = 3.6$.
\begin{itemize}
\item \textbf{Binomial PMF:} $P_B(X = x) = \binom{18}{x}(0.20)^x(0.80)^{18-x}$.
\item \textbf{Poisson PMF:} $P_P(X = x) = \dfrac{e^{-3.6}(3.6)^x}{x!}$.
\end{itemize}
\begin{center}
\small
\begin{tabular}{@{}ccccccccc@{}}
\toprule
$x$ & $0$ & $1$ & $2$ & $3$ & $4$ & $5$ & $6$ & $7$ \\
\midrule
Binomial $P_B(X=x)$ & $0.01801$ & $0.08106$ & $0.17226$ & $0.22968$ & $0.21533$ & $0.15073$ & $0.08165$ & $0.03499$ \\
Poisson $P_P(X=x)$ & $0.02732$ & $0.09837$ & $0.17706$ & $0.21247$ & $0.19122$ & $0.13768$ & $0.08261$ & $0.04248$ \\
\bottomrule
\end{tabular}
\end{center}
\begin{enumerate}[(i)]
\item \textbf{From 4 to 7 have blowouts ($P(4 \le X \le 7) = \sum_{x=4}^{7} P(X=x)$):}
\begin{itemize}
\item \textbf{Binomial:} $0.21533 + 0.15073 + 0.08165 + 0.03499 \approx \mathbf{0.48269}$.
\item \textbf{Poisson:} $0.19122 + 0.13768 + 0.08261 + 0.04248 = \mathbf{0.45399}$.
\end{itemize}
\item \textbf{Fewer than 3 have blowouts ($P(X < 3) = \sum_{x=0}^{2} P(X=x)$):}
\begin{itemize}
\item \textbf{Binomial:} $0.01801 + 0.08106 + 0.17226 = \mathbf{0.27134}$.
\item \textbf{Poisson:} $0.02732 + 0.09837 + 0.17706 = \mathbf{0.30275}$.
\end{itemize}
\item \textbf{More than 5 have blowouts ($P(X > 5) = 1 - \sum_{x=0}^{5} P(X=x)$):}
\begin{itemize}
\item \textbf{Binomial:} $1 - (0.27134 + 0.22968 + 0.21533 + 0.15073) = 1 - 0.86708 = \mathbf{0.13292}$.
\item \textbf{Poisson:} $1 - (0.30275 + 0.21247 + 0.19122 + 0.13768) = 1 - 0.84412 = \mathbf{0.15588}$.
\end{itemize}
\end{enumerate}
\ans{Binomial: (i) $0.4827$, (ii) $0.2713$, (iii) $0.1329$; \quad Poisson: (i) $0.4540$, (ii) $0.3028$, (iii) $0.1559$.}
\end{solution}

%---------------------------------------------------------------- P032
\begin{question}{2012--2013 \textbullet\ Q6(c)}{MATH 241 \textbullet\ 02/06/2014 \textbullet\ P032}{10}
Derive the mean and variance of the frequency function for the binomial distribution.
\end{question}
\begin{solution}
Let $X \sim \operatorname{Binomial}(n, p)$ with probability mass function $f(x) = \binom{n}{x}p^x q^{n-x}$ for $x = 0, 1, \dots, n$, where $p + q = 1$.

\medskip
\textbf{1. Derivation of the Mean ($\mu = E(X)$):}
\[ \begin{aligned} E(X) &= \sum_{x=0}^{n} x\,\binom{n}{x}p^x q^{n-x} = \sum_{x=1}^{n} x\cdot\frac{n!}{x!(n-x)!}p^x q^{n-x} \\ &= np\sum_{x=1}^{n} \frac{(n-1)!}{(x-1)!\,((n-1)-(x-1))!}p^{x-1}q^{(n-1)-(x-1)}. \end{aligned} \]
Letting $k = x - 1$ (so $k$ runs from $0$ to $n - 1$) and using the binomial expansion $(p + q)^{n-1} = 1$:
\[ \mu = E(X) = np\sum_{k=0}^{n-1} \binom{n-1}{k}p^k q^{n-1-k} = np(p + q)^{n-1} = np(1)^{n-1} = np. \]

\medskip
\textbf{2. Derivation of the Variance ($\sigma^2 = \operatorname{Var}(X)$):}
Write $x^2 = x(x - 1) + x$, so $E(X^2) = E[X(X - 1)] + E(X)$. Computing the second factorial moment $E[X(X - 1)]$:
\begin{align*}
E[X(X - 1)] &= \sum_{x=0}^{n} x(x - 1)\frac{n!}{x!(n-x)!}p^x q^{n-x} = n(n - 1)p^2\sum_{x=2}^{n} \frac{(n-2)!}{(x-2)!(n-x)!}p^{x-2}q^{n-x} \\
            &= n(n - 1)p^2(p + q)^{n-2} = n(n - 1)p^2.
\end{align*}
Therefore, $E(X^2) = n(n - 1)p^2 + np$, and the variance is:
\[ \sigma^2 = E(X^2) - [E(X)]^2 = n(n - 1)p^2 + np - (np)^2 = n^2p^2 - np^2 + np - n^2p^2 = np(1 - p) = npq. \]
\ans{Mean $\mu = np$, \qquad Variance $\sigma^2 = npq$.}
\end{solution}

%===================================================================== TOPIC 07
\topic{07}{Poisson Distribution}

%---------------------------------------------------------------- P033
\begin{question}{2023--2024 \textbullet\ Q3(b)}{MATH 243 \textbullet\ 21/05/2026 \textbullet\ P033}{12}
Establish the relation between binomial and Poisson distribution. Suppose telephone calls arrive at a switchboard at a mean rate of 0.5 calls per minute. Calculate the probability that two calls will arrive in a particular five-minute period.
\end{question}
\begin{solution}
\textbf{Part 1: Limiting Relation between Binomial and Poisson Distributions.}
Let $X \sim \operatorname{Binomial}(n, p)$ with PMF $b(x; n, p) = \binom{n}{x}p^x(1 - p)^{n-x}$. Suppose $n \to \infty$ and $p \to 0$ such that the mean $\mu = np$ remains constant (so $p = \mu/n$):
\begin{align*}
b(x; n, p) &= \frac{n(n - 1)(n - 2)\cdots(n - x + 1)}{x!}\left(\frac{\mu}{n}\right)^x\left(1 - \frac{\mu}{n}\right)^{n-x} \\
           &= \frac{\mu^x}{x!}\cdot 1\left(1 - \frac{1}{n}\right)\left(1 - \frac{2}{n}\right)\cdots\left(1 - \frac{x - 1}{n}\right)\left(1 - \frac{\mu}{n}\right)^n\left(1 - \frac{\mu}{n}\right)^{-x}.
\end{align*}
Taking the limit as $n \to \infty$ for any fixed integer $x \ge 0$:
\[ \lim_{n\to\infty}\left(1 - \frac{k}{n}\right) = 1, \qquad \lim_{n\to\infty}\left(1 - \frac{\mu}{n}\right)^{-x} = 1, \qquad \lim_{n\to\infty}\left(1 - \frac{\mu}{n}\right)^n = e^{-\mu}. \]
Hence,
\[ \lim_{\substack{n\to\infty,\ p\to 0 \\ np=\mu}} \binom{n}{x}p^x(1 - p)^{n-x} = \frac{e^{-\mu}\mu^x}{x!} = p(x; \mu). \]

\medskip
\textbf{Part 2: Switchboard Problem.}
Given arrival rate $\lambda_1 = 0.5\text{ calls/minute}$, the mean number of calls in a $t = 5\text{-minute}$ period is:
\[ \mu = \lambda_1 t = 0.5 \times 5 = 2.5\text{ calls}. \]
The probability that exactly $x = 2$ calls arrive in 5 minutes is:
\[ P(X = 2) = \frac{e^{-2.5}(2.5)^2}{2!} = \frac{0.082085 \times 6.25}{2} \approx 0.256516. \]
\ans{$P(X = 2) = \dfrac{e^{-2.5}(2.5)^2}{2!} \approx 0.2565$.}
\end{solution}

%---------------------------------------------------------------- P034
\begin{question}{2022--2023 \textbullet\ Q3(a)}{MATH 243 \textbullet\ 01/02/2025 \textbullet\ P034}{10}
Briefly discuss about Bernoulli and Poisson process. Derive Poisson distribution for Poisson random variable.
\end{question}
\begin{solution}
\textbf{1. Bernoulli Process:}
A Bernoulli process is a discrete-time sequence of independent and identically distributed (i.i.d.) binary trials $X_1, X_2, \dots$, where each trial results in either ``success'' (probability $p$) or ``failure'' (probability $q = 1 - p$), and $p$ remains constant across all trials. The number of successes in $n$ trials follows the Binomial distribution.

\medskip
\textbf{2. Poisson Process:}
A Poisson process is a continuous-time counting process $\{N(t), t \ge 0\}$ with mean rate $\lambda > 0$ satisfying:
\begin{enumerate}[(a)]
\item \emph{Stationarity:} The probability of $k$ occurrences in any interval of length $t$ depends only on $t$ (mean $\mu = \lambda t$).
\item \emph{Independent Increments:} Numbers of occurrences in disjoint time intervals are statistically independent.
\item \emph{Orderliness (No simultaneous events):} In an infinitesimal interval $\Delta t$, $P(\text{1 event}) = \lambda\Delta t + o(\Delta t)$ and $P(\ge 2\text{ events}) = o(\Delta t)$.
\end{enumerate}

\medskip
\textbf{3. Derivation of the Poisson Distribution:}
Divide a continuous interval of length $t$ into $n$ very small subintervals of length $\Delta t = t/n$. For large $n$, each subinterval acts as an independent Bernoulli trial with success probability $p = \lambda\Delta t = \dfrac{\lambda t}{n} = \dfrac{\mu}{n}$, where $\mu = \lambda t$.
The probability of observing $x$ events in the interval of length $t$ is the limit of the binomial probability as $n \to \infty$:
\[ \begin{aligned} P(X = x) &= \lim_{n\to\infty} \binom{n}{x}\left(\frac{\mu}{n}\right)^x\left(1 - \frac{\mu}{n}\right)^{n-x} \\ &= \frac{\mu^x}{x!}\lim_{n\to\infty}\frac{n!}{n^x(n-x)!}\left(1 - \frac{\mu}{n}\right)^n\left(1 - \frac{\mu}{n}\right)^{-x} = \frac{e^{-\mu}\mu^x}{x!}, \quad x = 0, 1, 2, \dots \end{aligned} \]
\ans{$p(x; \mu) = \dfrac{e^{-\mu}\mu^x}{x!}$ for $x = 0, 1, 2, \dots$, where $\mu = \lambda t$.}
\end{solution}

%---------------------------------------------------------------- P035
\begin{question}{2022--2023 \textbullet\ Q3(b)}{MATH 243 \textbullet\ 01/02/2025 \textbullet\ P035}{12}
Calculate the variance for a Poisson random variable. Now, suppose that, on average, 1 person in 1000 makes a numerical error in preparing his or her income tax return. If 10,000 returns are selected at random and examined, find the probability that 6, 7, or 8 of them contain an error.
\end{question}
\begin{solution}
\textbf{Part 1: Variance of a Poisson Random Variable.}
For $X \sim \operatorname{Poisson}(\mu)$ with PMF $p(x) = \dfrac{e^{-\mu}\mu^x}{x!}$, first find the mean $E(X)$:
\[ E(X) = \sum_{x=0}^{\infty} x\frac{e^{-\mu}\mu^x}{x!} = \mu e^{-\mu}\sum_{x=1}^{\infty} \frac{\mu^{x-1}}{(x-1)!} = \mu e^{-\mu}e^{\mu} = \mu. \]
Next, compute the second factorial moment $E[X(X - 1)]$:
\[ E[X(X - 1)] = \sum_{x=0}^{\infty} x(x - 1)\frac{e^{-\mu}\mu^x}{x!} = \mu^2 e^{-\mu}\sum_{x=2}^{\infty} \frac{\mu^{x-2}}{(x-2)!} = \mu^2 e^{-\mu}e^{\mu} = \mu^2. \]
Since $E(X^2) = E[X(X - 1)] + E(X) = \mu^2 + \mu$, the variance is:
\[ \operatorname{Var}(X) = E(X^2) - [E(X)]^2 = (\mu^2 + \mu) - \mu^2 = \mu. \]

\medskip
\textbf{Part 2: Income Tax Return Problem.}
Here $n = 10{,}000$ and $p = \dfrac{1}{1000} = 0.001$, so the Poisson mean is $\mu = np = 10{,}000 \times 0.001 = 10$.
The probability that $X \in \{6, 7, 8\}$ is:
\[ \begin{aligned} P(6 \le X \le 8) &= e^{-10}\left(\frac{10^6}{6!} + \frac{10^7}{7!} + \frac{10^8}{8!}\right) \\ &= 0.063055 + 0.090079 + 0.112599 = 0.265734. \end{aligned} \]
\ans{$\operatorname{Var}(X) = \mu$; \qquad $P(X = 6, 7, \text{ or } 8) \approx 0.2657$.}
\end{solution}

%---------------------------------------------------------------- P036
\begin{question}{2019--2020 \textbullet\ Q8(a)}{MATH 245 \textbullet\ 25/09/2021 \textbullet\ P036}{10}
In a manufacturing process where glass products are made, on average, 1 in every 1000 of these items produced has one or more bubbles. What is the probability that a random sample of 8000 will yield fewer than 7 items possessing bubbles?
\end{question}
\begin{solution}
Given $n = 8000$ and defect probability $p = \dfrac{1}{1000} = 0.001$. Since $n$ is large and $p$ is very small, the number of items with bubbles $X$ follows a Poisson distribution with mean:
\[ \mu = np = 8000 \times 0.001 = 8. \]
The probability of finding fewer than 7 items with bubbles ($X < 7$, i.e.\ $X \le 6$) is:
\[ P(X < 7) = \sum_{x=0}^{6} \frac{e^{-8}8^x}{x!} = e^{-8}\left(1 + 8 + \frac{64}{2} + \frac{512}{6} + \frac{4096}{24} + \frac{32768}{120} + \frac{262144}{720}\right). \]
Evaluating term by term ($e^{-8} \approx 0.00033546$):
\begin{center}
\small
\begin{tabular}{@{}cccccccc@{}}
\toprule
$x$ & $0$ & $1$ & $2$ & $3$ & $4$ & $5$ & $6$ \\
\midrule
$P(X=x)$ & $0.00034$ & $0.00268$ & $0.01073$ & $0.02863$ & $0.05725$ & $0.09160$ & $0.12214$ \\
\bottomrule
\end{tabular}
\end{center}
Summing from $x = 0$ to $6$:
\[ P(X < 7) = e^{-8}\times\frac{42037}{45} = 934.1556 \times 0.00033546 \approx 0.313374. \]
\ans{$P(X < 7) \approx 0.3134$.}
\end{solution}

%---------------------------------------------------------------- P037
\begin{question}{2016--2017 \textbullet\ Q8(a)}{MATH 245 \textbullet\ 05/08/2017 \textbullet\ P037}{20}
If $X$ is a binomial random variable with probability distribution $b(x; n, p)$, when $n \to \infty, p \to 0$ and $np \to \mu$ remains constant, prove that $b(x; n, p) \to p(x, \mu)$. In a certain industrial facility accidents occur infrequently. It is known that the probability of an accident on any given day is 0.005 and accidents are independent of each other. What is the probability that:
(i) in any given period of 400 days there will be an accident on one day, and
(ii) there are at most three days with an accident.
\end{question}
\begin{solution}
\textbf{Part 1: Proof of $b(x; n, p) \to p(x, \mu)$.}
\textit{(See the full step-by-step limiting derivation in \textbf{P033 [2023--2024 Q3(b)]} above.)}

\medskip
\textbf{Part 2: Industrial Facility Accident Probability.}
Here $n = 400\text{ days}$ and $p = 0.005$, so the Poisson mean is $\mu = np = 400 \times 0.005 = 2$.
\begin{enumerate}[(i)]
\item \textbf{An accident on exactly one day ($X = 1$):}
\[ P(X = 1) = \frac{e^{-2}(2)^1}{1!} = 2e^{-2} \approx 2(0.135335) = 0.270671. \]
\item \textbf{At most three days with an accident ($X \le 3$):}
\[ P(X \le 3) = \sum_{x=0}^{3} \frac{e^{-2}2^x}{x!} = e^{-2}\left(1 + 2 + \frac{4}{2} + \frac{8}{6}\right) = \frac{19}{3}e^{-2} \approx 6.3333 \times 0.135335 = 0.857123. \]
\end{enumerate}
\ans{(i) $P(X = 1) = 2e^{-2} \approx 0.2707$, \qquad (ii) $P(X \le 3) = \dfrac{19}{3}e^{-2} \approx 0.8571$.}
\end{solution}

%---------------------------------------------------------------- P038
\begin{question}{2011--2012 \textbullet\ Q7(b)(i)}{MATH 241 \textbullet\ 07/01/2013 \textbullet\ P038}{7}
An insurance company finds that the probability of dying of an individual belonging to a certain age group within a year is 0.0005. If 1000 persons of this age group buy a policy of this company, find the probability that the company will receive 5 or more claims during a given year.
\end{question}
\begin{solution}
Given $n = 1000$ policyholders and individual mortality probability $p = 0.0005$. Using the Poisson distribution with mean $\mu = np = 1000 \times 0.0005 = 0.5$:
\[ P(X \ge 5) = 1 - P(X \le 4) = 1 - \sum_{x=0}^{4} \frac{e^{-0.5}(0.5)^x}{x!}. \]
Evaluating the sum for $x = 0, 1, 2, 3, 4$ ($e^{-0.5} \approx 0.60653066$):
\[ \sum_{x=0}^{4} \frac{e^{-0.5}(0.5)^x}{x!} = e^{-0.5}\left(1 + 0.5 + \frac{0.25}{2} + \frac{0.125}{6} + \frac{0.0625}{24}\right) = e^{-0.5}(1.6484375) \approx 0.9998279. \]
Therefore,
\[ P(X \ge 5) = 1 - 0.9998279 = 0.0001721. \]
\ans{$P(X \ge 5) \approx 0.000172$.}
\end{solution}

%===================================================================== TOPIC 08
\topic[Normal Distribution]{08}{Normal Distribution \& Standard Normal Approximation}

%---------------------------------------------------------------- P039
\begin{question}{2023--2024 \textbullet\ Q4(a)}{MATH 243 \textbullet\ 21/05/2026 \textbullet\ P039}{11}
Show that the mode of normal distribution is equal to $\mu$, where $\mu$ is the mean of that distribution.
Phillips Bangladesh manufactures electric bulbs that have a length of life that is normally distributed with mean equal to 800 hours and standard deviation of 40 hours. Find the probability that a bulb burns between 778 and 834 hours. (Necessary table is attached)
\end{question}
\begin{solution}
\textbf{Part 1: Proof that the Mode of $N(\mu, \sigma^2)$ Equals $\mu$.}
The probability density function of a normal random variable $X \sim N(\mu, \sigma^2)$ is:
\[ f(x) = \frac{1}{\sigma\sqrt{2\pi}}\exp\left(-\frac{(x-\mu)^2}{2\sigma^2}\right), \qquad -\infty < x < \infty. \]
The mode is the value of $x$ that maximizes $f(x)$. Differentiating $f(x)$ with respect to $x$:
\[ f'(x) = \frac{1}{\sigma\sqrt{2\pi}}\exp\left(-\frac{(x-\mu)^2}{2\sigma^2}\right)\cdot\left(-\frac{x-\mu}{\sigma^2}\right) = -\frac{x-\mu}{\sigma^2}f(x). \]
Setting $f'(x) = 0$: since $f(x) > 0$ for all $x$, we obtain $x - \mu = 0 \implies x = \mu$.
Differentiating a second time using the product rule:
\[ f''(x) = -\frac{1}{\sigma^2}f(x) - \frac{x-\mu}{\sigma^2}f'(x) = \frac{f(x)}{\sigma^2}\left[\frac{(x-\mu)^2}{\sigma^2} - 1\right]. \]
Evaluating at $x = \mu$:
\[ f''(\mu) = -\frac{f(\mu)}{\sigma^2} = -\frac{1}{\sigma^3\sqrt{2\pi}} < 0. \]
Since $f'(\mu) = 0$ and $f''(\mu) < 0$, $f(x)$ attains its unique global maximum at $x = \mu$. Thus, $\text{Mode} = \mu$.

\medskip
\textbf{Part 2: Electric Bulb Lifetime Probability.}
Given $X \sim N(\mu = 800, \sigma^2 = 40^2)$. Standardizing to $Z = \dfrac{X - \mu}{\sigma}$ at $x_1 = 778$ and $x_2 = 834$:
\[ z_1 = \frac{778 - 800}{40} = \frac{-22}{40} = -0.55, \qquad z_2 = \frac{834 - 800}{40} = \frac{34}{40} = 0.85. \]
From the attached cumulative standard normal table ($\Phi(z) = P(Z \le z)$):
\[ \Phi(0.85) = 0.802337, \qquad \Phi(-0.55) = 1 - \Phi(0.55) = 1 - 0.708840 = 0.291160. \]
Therefore,
\[ P(778 < X < 834) = \Phi(0.85) - \Phi(-0.55) = 0.802337 - 0.291160 = 0.511177. \]
\ans{$\text{Mode} = \mu$; \qquad $P(778 < X < 834) \approx 0.5112$.}
\end{solution}

%---------------------------------------------------------------- P040
\begin{question}{2022--2023 \textbullet\ Q3(c)}{MATH 243 \textbullet\ 01/02/2025 \textbullet\ P040}{13}
A lawyer commutes daily from his suburban home to his midtown office. The average time for a one-way trip is 24 minutes, with a standard deviation of 3.8 minutes. Assuming the distribution of trip times is normally distributed, what is the probability that a trip will take at least half an hour?
(i) If the office opens at 9:00 am and the lawyer leaves his house at 8:45 am daily, what percentage of the time is he late for work?
(ii) Find the length of time above which we find the slowest 15\% of the trips.
(iii) Find the probability that 2 of the next 3 trips will take at least half an hour. (Necessary table is attached)
\end{question}
\begin{solution}
Let $X \sim N(\mu = 24, \sigma = 3.8)$ be the one-way trip time in minutes.

\medskip
\textbf{Preliminary: Probability that a trip takes at least half an hour ($X \ge 30$):}
\[ z = \frac{30 - 24}{3.8} = \frac{6}{3.8} \approx 1.58 \implies P(X \ge 30) = 1 - \Phi(1.58) = 1 - 0.9429 = 0.0571. \]
\begin{enumerate}[(i)]
\item \textbf{Percentage of the time the lawyer is late:}
Leaving at 8:45 am for a 9:00 am opening gives $15\text{ minutes}$ of travel time, so he is late whenever $X > 15$:
\[ z = \frac{15 - 24}{3.8} = \frac{-9}{3.8} \approx -2.37 \implies P(X > 15) = P(Z > -2.37) = \Phi(2.37) = 0.9911 \ (99.11\%). \]
\item \textbf{Length of time $x_0$ above which the slowest 15\% of trips fall:}
We require $P(X > x_0) = 0.15 \iff \Phi(z_0) = 1 - 0.15 = 0.85$. From the standard normal table, $z_0 \approx 1.04$ (or $1.0364$ interpolated):
\[ x_0 = \mu + z_0\sigma = 24 + (1.04)(3.8) = 27.952\text{ min} \quad (\text{or } 27.94\text{ min using } z_0 = 1.0364). \]
\item \textbf{Probability that 2 of the next 3 trips take at least half an hour:}
Each trip is an independent Bernoulli trial with $p = P(X \ge 30) = 0.0571$ and $q = 1 - p = 0.9429$. For $n = 3$ trips:
\[ P(Y = 2) = \binom{3}{2}p^2 q = 3(0.0571)^2(0.9429) \approx 0.00922. \]
\end{enumerate}
\ans{$P(X\ge 30) = 0.0571$; \quad (i) $99.11\%$; \quad (ii) $27.95\text{ minutes}$; \quad (iii) $0.00922$.}
\end{solution}

%---------------------------------------------------------------- P041
\begin{question}{2022--2023 \textbullet\ Q4(a)}{MATH 243 \textbullet\ 01/02/2025 \textbullet\ P041}{10}
Write down some properties of normal curve. Find the mean, mode, and standard deviation of normal random variable. Hence comment on the mean and variance of standard normal random variable.
\end{question}
\begin{solution}
\textbf{1. Key Properties of the Normal Curve $f(x) = \frac{1}{\sigma\sqrt{2\pi}}e^{-(x-\mu)^2/(2\sigma^2)}$:}
\begin{itemize}
\item \textbf{Symmetry \& Bell Shape:} The curve is bell-shaped and perfectly symmetric about the vertical line $x = \mu$.
\item \textbf{Coincidence of Central Measures:} The mean, median, and mode coincide at $x = \mu$, where $f(x)$ reaches its maximum height $\frac{1}{\sigma\sqrt{2\pi}}$.
\item \textbf{Points of Inflection:} The curve has points of inflection at $x = \mu \pm \sigma$ (where $f''(x) = 0$).
\item \textbf{Asymptotic Behaviour \& Total Area:} The tails approach the horizontal axis asymptotically as $x \to \pm\infty$, and the total area under the curve over $(-\infty, \infty)$ is $1$.
\end{itemize}

\medskip
\textbf{2. Mean, Mode, and Standard Deviation of $X \sim N(\mu, \sigma^2)$:}
\begin{itemize}
\item \textbf{Mean:} Substituting $z = \dfrac{x-\mu}{\sigma}$ ($x = \mu + \sigma z$, $dx = \sigma\,dz$):
\[ E(X) = \frac{1}{\sqrt{2\pi}}\int_{-\infty}^{\infty} (\mu + \sigma z)e^{-z^2/2}\,dz = \mu\underbrace{\frac{1}{\sqrt{2\pi}}\int_{-\infty}^{\infty} e^{-z^2/2}\,dz}_{=1} + \sigma\underbrace{\frac{1}{\sqrt{2\pi}}\int_{-\infty}^{\infty} z e^{-z^2/2}\,dz}_{=0\text{ (odd integrand)}} = \mu. \]
\item \textbf{Mode:} Setting $f'(x) = -\dfrac{x-\mu}{\sigma^2}f(x) = 0$ gives $x = \mu$, with $f''(\mu) = -\dfrac{f(\mu)}{\sigma^2} < 0$. Hence $\text{Mode} = \mu$.
\item \textbf{Variance and Standard Deviation:} Using integration by parts with $u = z$ and $dv = z e^{-z^2/2}\,dz$:
\[ \begin{aligned} \operatorname{Var}(X) &= \frac{\sigma^2}{\sqrt{2\pi}}\int_{-\infty}^{\infty} z^2 e^{-z^2/2}\,dz \\ &= \frac{\sigma^2}{\sqrt{2\pi}}\left(\Big[-z e^{-z^2/2}\Big]_{-\infty}^{\infty} + \int_{-\infty}^{\infty} e^{-z^2/2}\,dz\right) = \sigma^2(0 + 1) = \sigma^2. \end{aligned} \]
Taking the positive square root gives the standard deviation $\text{SD}(X) = \sigma$.
\end{itemize}

\medskip
\textbf{3. Comment on Standard Normal Random Variable $Z = \frac{X - \mu}{\sigma}$:}
For the standard normal variable $Z$, shifting by $-\mu$ centers the distribution at $E(Z) = \frac{E(X) - \mu}{\sigma} = 0$, and scaling by $1/\sigma$ normalizes the spread to $\operatorname{Var}(Z) = \frac{\operatorname{Var}(X)}{\sigma^2} = 1$, so $Z \sim N(0, 1)$.
\ans{For $X \sim N(\mu,\sigma^2)$: $\text{Mean} = \mu$, $\text{Mode} = \mu$, $\text{SD} = \sigma$. For $Z \sim N(0,1)$: $\text{Mean} = 0$ and $\text{Variance} = 1$.}
\end{solution}

%---------------------------------------------------------------- P042
\begin{question}{2021--2022 \textbullet\ Q7(b)}{MATH 245 \textbullet\ 30/04/2023 \textbullet\ P042}{15}
In an industrial process, the buyer sets specifications on the diameter of a ball bearing to be $3.0 \pm 0.01\text{ cm}$, with the implication that no part falling outside these specifications will be accepted. In the process, the diameter of a ball bearing has a normal distribution with mean 3.0 and standard deviation 0.005. On average, how many manufactured ball bearings will be scrapped? (Table 1 is attached.)
\end{question}
\begin{solution}
Let $X \sim N(\mu = 3.0, \sigma = 0.005)$ be the diameter of a ball bearing in cm. The acceptable specification interval is $[3.0 - 0.01, 3.0 + 0.01] = [2.99, 3.01]$. Standardizing the lower and upper limits:
\[ z_1 = \frac{2.99 - 3.0}{0.005} = -2.0, \qquad z_2 = \frac{3.01 - 3.0}{0.005} = +2.0. \]
The probability that a bearing is accepted is:
\[ P(2.99 \le X \le 3.01) = P(-2.0 \le Z \le 2.0) = 2\Phi(2.0) - 1. \]
From the attached Table 1, $\Phi(2.0) = 0.97725$. Therefore, the probability that a bearing falls outside specifications and is scrapped is:
\[ P(\text{scrapped}) = 1 - P(-2.0 \le Z \le 2.0) = 2\big(1 - \Phi(2.0)\big) = 2(1 - 0.97725) = 0.0455. \]
\ans{On average, $4.55\%$ (or $455$ out of every $10{,}000$) of the manufactured ball bearings will be scrapped.}
\end{solution}

%---------------------------------------------------------------- P043
\begin{question}{2020--2021 \textbullet\ Q7(b)}{MATH 245 \textbullet\ 2020--2021 \textbullet\ P043}{18}
A research scientist reports that mice will live an average of 40 months when their diets are sharply restricted and then enriched with vitamins and proteins. Assuming that the lifetimes of such mice are normally distributed with a standard deviation of 6.3 months, find the probability that a given mouse will live:
(i) more than 32 months,
(ii) between 37 and 49 months. (Necessary Table 1 and 2 are attached)
\end{question}
\begin{solution}
Given mouse lifetime $X \sim N(\mu = 40, \sigma = 6.3)$ months:
\begin{enumerate}[(i)]
\item \textbf{Probability that a mouse lives more than 32 months ($P(X > 32)$):}
\[ z = \frac{32 - 40}{6.3} = \frac{-8}{6.3} \approx -1.27. \]
From the attached normal table, $\Phi(-1.27) = 0.1020$:
\[ P(X > 32) = P(Z > -1.27) = 1 - \Phi(-1.27) = 1 - 0.1020 = 0.8980. \]
\item \textbf{Probability that a mouse lives between 37 and 49 months ($P(37 < X < 49)$):}
\[ z_1 = \frac{37 - 40}{6.3} = \frac{-3}{6.3} \approx -0.48, \qquad z_2 = \frac{49 - 40}{6.3} = \frac{9}{6.3} \approx 1.43. \]
From the attached normal tables, $\Phi(1.43) = 0.9236$ and $\Phi(-0.48) = 0.3156$:
\[ P(37 < X < 49) = \Phi(1.43) - \Phi(-0.48) = 0.9236 - 0.3156 = 0.6080. \]
\end{enumerate}
\ans{(i) $P(X > 32) = 0.8980$, \qquad (ii) $P(37 < X < 49) = 0.6080$.}
\end{solution}

%---------------------------------------------------------------- P044
\begin{question}{2018--2019 \textbullet\ Q8(a)}{MATH 245 \textbullet\ 03/11/2019 \textbullet\ P044}{18}
A new factory installed 1000 tube lights with a mean life of 120 days and their length of life is normally distributed with standard deviation 20 days. (i) How many tube lights will expire in less than 90 days? (ii) If it is decided to replace all the tube lights together, what interval should be allowed between replacements if not more than 10 percent should expire before replacement?
\end{question}
\begin{solution}
Let $X \sim N(\mu = 120, \sigma = 20)$ days be the lifetime of a tube light, with $N = 1000$ lights installed.
\begin{enumerate}[(i)]
\item \textbf{Number of tube lights expiring in less than 90 days:}
\[ z = \frac{90 - 120}{20} = -1.50 \implies P(X < 90) = \Phi(-1.50) = 0.0668. \]
Expected number of tube lights expiring before 90 days:
\[ 1000 \times 0.0668 = 66.8 \approx 67\text{ tube lights}. \]
\item \textbf{Replacement interval $t$ such that $P(X \le t) = 0.10$:}
From the standard normal table, $\Phi(z) = 0.10$ corresponds to $z \approx -1.28$ (or $-1.2816$):
\[ \frac{t - 120}{20} = -1.28 \implies t = 120 - 1.28(20) = 120 - 25.6 = 94.4\text{ days} \quad (\text{or } 94.37\text{ days}). \]
\end{enumerate}
\ans{(i) $67\text{ tube lights}$ ($66.8$), \qquad (ii) $94.4\text{ days}$.}
\end{solution}

%---------------------------------------------------------------- P045
\begin{question}{2016--2017 \textbullet\ Q7(b)}{MATH 245 \textbullet\ 05/08/2017 \textbullet\ P045}{10}
A U-shaped component is to be formed from the three parts A, B, and C as shown below. The length of A is normally distributed with a mean of 10 millimeters and a standard deviation of 0.1 millimeter. The thickness of parts B and C is normally distributed with a mean of 2 millimeters and a standard deviation of 0.05 millimeter. Assume all dimensions are independent.
\fig{\ushapefig}
(i) Determine the mean and standard deviation of the length of the gap D.
(ii) What is the probability that the gap D is less than 5.9 millimeters? (Necessary table attached).
\end{question}
\begin{solution}
From the geometry of the U-shaped component, the inner gap length is the linear combination $D = A - B - C$, where $A \sim N(10, 0.1^2)$, $B \sim N(2, 0.05^2)$, and $C \sim N(2, 0.05^2)$ are mutually independent.
\begin{enumerate}[(i)]
\item \textbf{Mean and Standard Deviation of Gap $D$:}
By linearity of expectation and independence of $A, B, C$ (note that variances \textbf{add} even when variables are subtracted):
\begin{align*}
\mu_D &= E(A - B - C) = \mu_A - \mu_B - \mu_C = 10 - 2 - 2 = 6.0\text{ mm}, \\[4pt]
\sigma_D^2 &= \sigma_A^2 + \sigma_B^2 + \sigma_C^2 = (0.1)^2 + (0.05)^2 + (0.05)^2 = 0.015\text{ mm}^2, \\[4pt]
\sigma_D &= \sqrt{0.015} \approx 0.12247\text{ mm}.
\end{align*}
\item \textbf{Probability that $D < 5.9\text{ mm}$:}
Since a linear combination of independent normal variables is normally distributed, $D \sim N(6.0, 0.015)$:
\[ z = \frac{5.9 - 6.0}{\sqrt{0.015}} = \frac{-0.1}{0.122474} \approx -0.82. \]
From the attached normal table, $\Phi(0.82) = 0.7939$:
\[ P(D < 5.9) = \Phi(-0.82) = 1 - 0.7939 = 0.2061. \]
\end{enumerate}
\ans{(i) $\mu_D = 6.0\text{ mm}$, $\sigma_D = \sqrt{0.015} \approx 0.1225\text{ mm}$; \qquad (ii) $P(D < 5.9) \approx 0.2061$.}
\end{solution}

%---------------------------------------------------------------- P046
\begin{question}{2015--2016 \textbullet\ Q8(b)}{MATH 241 \textbullet\ 07/08/2016 \textbullet\ P046}{12}
The mark in a test is assumed to be normally distributed. It is known that 10\% of the students have marks under 65 and 25\% exceed 75. What percentage of students have marks between 60 and 74? (Necessary chart 1 is attached)
\end{question}
\begin{solution}
Let $X \sim N(\mu, \sigma^2)$ denote the test mark. From the two given percentile conditions:
\begin{enumerate}[1.]
\item $P(X < 65) = 0.10 \implies \Phi\left(\dfrac{65 - \mu}{\sigma}\right) = 0.10 \implies \dfrac{65 - \mu}{\sigma} = -1.28 \implies \mu - 1.28\sigma = 65$.
\item $P(X > 75) = 0.25 \implies \Phi\left(\dfrac{75 - \mu}{\sigma}\right) = 0.75 \implies \dfrac{75 - \mu}{\sigma} = 0.67 \implies \mu + 0.67\sigma = 75$.
\end{enumerate}
Subtracting the first equation from the second:
\[ 1.95\sigma = 10 \implies \sigma = \frac{10}{1.95} \approx 5.1282, \qquad \mu = 65 + 1.28(5.1282) \approx 71.5641. \]
Now standardize $x_1 = 60$ and $x_2 = 74$:
\[ z_1 = \frac{60 - 71.5641}{5.1282} \approx -2.255 \ (\approx -2.26), \qquad z_2 = \frac{74 - 71.5641}{5.1282} \approx 0.475 \ (\approx 0.48). \]
From Chart 1, $\Phi(0.48) = 0.6844$ (or $\Phi(0.475) \approx 0.6826$) and $\Phi(-2.26) = 0.0119$ (or $\Phi(-2.255) \approx 0.0121$):
\[ P(60 < X < 74) = \Phi(0.48) - \Phi(-2.26) = 0.6844 - 0.0119 = 0.6725 \ (67.25\%). \]
\ans{Approximately $67.25\%$ (or $67.05\%$ using interpolated $z$-values) of students score between 60 and 74.}
\end{solution}

%---------------------------------------------------------------- P047
\begin{question}{2014--2015 \textbullet\ Q8(c)}{MATH 241 \textbullet\ 06/02/2016 \textbullet\ P047}{11}
A workshop produces 2000 units per day. The average weight of units is 130 kg with a standard deviation of 10 kg. Assuming normal distribution, how many units are expected to weigh less than 142 kg? (Necessary Chart 2 is attached).
\end{question}
\begin{solution}
Let $X \sim N(\mu = 130, \sigma = 10)\text{ kg}$ denote the weight of a unit. Standardizing $x = 142\text{ kg}$:
\[ z = \frac{142 - 130}{10} = 1.20. \]
From Chart 2, $P(X < 142) = \Phi(1.20) = 0.8849$. Out of $N = 2000$ units produced daily, the expected number weighing less than 142 kg is:
\[ 2000 \times 0.8849 = 1769.8 \approx 1770\text{ units}. \]
\ans{Expected number of units weighing less than 142 kg is $1769.8 \approx 1770\text{ units}$.}
\end{solution}

%===================================================================== TOPIC 09
\topic[Chebyshev, Multinomial \& Uniform]{09}{Chebyshev's Theorem, Multinomial \& Uniform Distributions}

\begin{note}
No standalone question on the \emph{Continuous Uniform Distribution} appeared in the 2011--2024 papers (for reference, $X \sim \operatorname{Uniform}(a,b)$ has $f(x) = \frac{1}{b-a}$ on $[a,b]$, mean $\mu = \frac{a+b}{2}$, and variance $\sigma^2 = \frac{(b-a)^2}{12}$).
\end{note}

%---------------------------------------------------------------- P048
\begin{question}{2023--2024 \textbullet\ Q3(a)}{MATH 243 \textbullet\ 21/05/2026 \textbullet\ P048}{11}
State Chebyshev's theorem. For 70 new job positions at an automobile plant, 1000 applicants show up and take a selection test. The mean grade on this test turns out to be 60, and the scores have a standard deviation of 6. Can a person who scores 84 count on getting one of the jobs? Assume that the distribution is symmetric about the mean.
\end{question}
\begin{solution}
\textbf{Part 1: Statement of Chebyshev's Theorem.}
For any random variable $X$ (discrete or continuous) with finite mean $\mu$ and finite non-zero standard deviation $\sigma$, the probability that $X$ assumes a value within $k$ standard deviations of the mean ($k > 0$) is at least $1 - \dfrac{1}{k^2}$:
\[ P(|X - \mu| < k\sigma) \ge 1 - \frac{1}{k^2}, \qquad \text{or equivalently,} \qquad P(|X - \mu| \ge k\sigma) \le \frac{1}{k^2}. \]

\medskip
\textbf{Part 2: Job Applicant Analysis.}
Here $N = 1000$ applicants, $\mu = 60$, and $\sigma = 6$. A score of $x = 84$ lies above the mean by:
\[ x - \mu = 84 - 60 = 24 = 4(6) = 4\sigma \implies k = 4. \]
By Chebyshev's inequality in two-tailed form:
\[ P(|X - 60| \ge 24) = P(X \le 36 \text{ or } X \ge 84) \le \frac{1}{4^2} = \frac{1}{16} = 0.0625. \]
Because the distribution of test scores is given to be \textbf{symmetric about the mean $\mu = 60$}, the two tails have equal probability:
\[ P(X \ge 84) = \frac{1}{2}P(|X - 60| \ge 24) \le \frac{1}{2}\left(\frac{1}{16}\right) = \frac{1}{32} = 0.03125. \]
Therefore, out of $1000$ applicants, the maximum number of applicants who could score $84$ or higher is at most:
\[ 1000 \times \frac{1}{32} = 31.25 \approx 31\text{ applicants}. \]
Since there are $70$ positions available and at most $31$ applicants score $\ge 84$, an applicant who scores $84$ is guaranteed to be among the top $70$ candidates.
\ans{Yes; at most $1000/32 \approx 31$ applicants score $84$ or higher, which is well within the $70$ available job positions.}
\end{solution}

%---------------------------------------------------------------- P049
\begin{question}{2022--2023 \textbullet\ Q2(c)}{MATH 243 \textbullet\ 01/02/2025 \textbullet\ P049}{10}
Compute $p(\mu - 2\sigma < X < \mu + 2\sigma)$, where $X$ has the density function
\[ f(x) = \begin{cases} 6x(1-x), & 0 < x < 1 \\ 0, & \text{elsewhere} \end{cases} \]
and compare with the result given in Chebyshev's theorem.
\end{question}
\begin{solution}
\textbf{Step 1: Find the mean $\mu$ and standard deviation $\sigma$ of $X$.}
\begin{align*}
\mu = E(X) &= \int_{0}^{1} x\cdot 6x(1 - x)\,dx = 6\int_{0}^{1} (x^2 - x^3)\,dx = 6\left[\frac{x^3}{3} - \frac{x^4}{4}\right]_{0}^{1} = 6\left(\frac{1}{12}\right) = \frac{1}{2} = 0.5, \\[4pt]
E(X^2) &= \int_{0}^{1} x^2\cdot 6x(1 - x)\,dx = 6\int_{0}^{1} (x^3 - x^4)\,dx = 6\left[\frac{x^4}{4} - \frac{x^5}{5}\right]_{0}^{1} = 6\left(\frac{1}{20}\right) = \frac{3}{10} = 0.3, \\[4pt]
\sigma^2 &= E(X^2) - \mu^2 = \frac{3}{10} - \frac{1}{4} = \frac{1}{20} = 0.05 \implies \sigma = \frac{1}{\sqrt{20}} = \frac{1}{2\sqrt{5}} \approx 0.223607.
\end{align*}
\textbf{Step 2: Compute the exact probability $P(\mu - 2\sigma < X < \mu + 2\sigma)$.}
The interval endpoints are:
\[ a = \mu - 2\sigma = \frac{1}{2} - \frac{1}{\sqrt{5}} \approx 0.052786, \qquad b = \mu + 2\sigma = \frac{1}{2} + \frac{1}{\sqrt{5}} \approx 0.947214. \]
Since $[a, b] \subset (0, 1)$, integrating $f(x) = 6x - 6x^2$ (whose antiderivative is $F(x) = 3x^2 - 2x^3$) from $a$ to $b$:
\[ \begin{aligned} P(\mu - 2\sigma < X < \mu + 2\sigma) &= \Big[3x^2 - 2x^3\Big]_{0.052786}^{0.947214} = 0.991935 - 0.008065 \\ &= 0.983870 \quad \left(\text{exact: }\frac{11\sqrt{5}}{25}\right). \end{aligned} \]
\textbf{Step 3: Comparison with Chebyshev's Theorem.}
For $k = 2$, Chebyshev's theorem guarantees $P(\mu - 2\sigma < X < \mu + 2\sigma) \ge 1 - \frac{1}{2^2} = \frac{3}{4} = 0.75$. The exact probability ($0.9839$) is substantially higher than the universal distribution-free lower bound ($0.75$).
\ans{Exact $P(\mu - 2\sigma < X < \mu + 2\sigma) = \dfrac{11\sqrt{5}}{25} \approx 0.9839$, which is well above Chebyshev's lower bound of $0.75$.}
\end{solution}

%---------------------------------------------------------------- P050
\begin{question}{2021--2022 \textbullet\ Q6(b)}{MATH 245 \textbullet\ 30/04/2023 \textbullet\ P050}{10}
The probabilities are 0.4, 0.2, 0.3, and 0.1, respectively, that a delegate to a certain convention arrived by air, bus, automobile, or train. What is the probability that among 9 delegates randomly selected at this convention, 3 arrived by air, 3 arrived by bus, 1 arrived by automobile, and 2 arrived by train?
\end{question}
\begin{solution}
This is a multinomial experiment with $n = 9$ independent trials and $k = 4$ mutually exclusive outcomes:
\[ p_1 = 0.4 \ (\text{air}), \quad p_2 = 0.2 \ (\text{bus}), \quad p_3 = 0.3 \ (\text{automobile}), \quad p_4 = 0.1 \ (\text{train}), \]
with counts $x_1 = 3,\ x_2 = 3,\ x_3 = 1,\ x_4 = 2$ ($\sum x_i = 9$). By the multinomial probability formula:
\begin{align*}
P(X_1=3, X_2=3, X_3=1, X_4=2) &= \frac{9!}{3!\,3!\,1!\,2!}\,p_1^3\,p_2^3\,p_3^1\,p_4^2 \\
&= \frac{362{,}880}{6 \times 6 \times 1 \times 2}(0.4)^3(0.2)^3(0.3)^1(0.1)^2 \\
&= 5040 \times (0.064) \times (0.008) \times (0.3) \times (0.01) = 0.00774144.
\end{align*}
\ans{$P = 0.007741$.}
\end{solution}

%===================================================================== TOPIC 10
\topic[Hypothesis Testing]{10}{Hypothesis Testing (z, t, Chi-Square \& Proportion Tests)}

%---------------------------------------------------------------- P051
\begin{question}{2023--2024 \textbullet\ Q6(c)}{MATH 243 \textbullet\ 21/05/2026 \textbullet\ P051}{12}
Explain with an example, stratified random sampling. An online medicine shop claims that the mean delivery time for medicines is less than 120 minutes with a standard deviation of 30 minutes. Is there enough evidence to support this claim at a 0.05 significance level if 49 orders were examined with a mean of 100 minutes? Given that $Z = -1.645$ at 5\% level of significance.
\end{question}
\begin{solution}
\textbf{Part 1: Stratified Random Sampling.}
Stratified random sampling is a probability sampling technique in which a heterogeneous population of size $N$ is first partitioned into non-overlapping, homogeneous subgroups called \emph{strata} (based on a shared characteristic such as age group, department, or geographic region), and a simple random sample is then drawn independently from each stratum (often proportional to stratum size, $n_h = n\frac{N_h}{N}$).
\begin{itemize}
\item \emph{Example:} To survey student satisfaction across BUET's 4 undergraduate levels (L-1 to L-4), the student body is divided into 4 strata by level, and 25 students are randomly selected from each level. This guarantees representation of every level and reduces sampling variance.
\end{itemize}

\medskip
\textbf{Part 2: Hypothesis Test on Medicine Delivery Time.}
Given $\mu_0 = 120\text{ min}$, $\sigma = 30\text{ min}$, $n = 49$, $\bar{x} = 100\text{ min}$, and $\alpha = 0.05$:
\begin{enumerate}[1.]
\item \textbf{Hypotheses:} $H_0: \mu = 120\text{ min}$ vs.\ $H_1: \mu < 120\text{ min}$ (left-tailed test supporting the shop's claim).
\item \textbf{Test Statistic:}
\[ z = \frac{\bar{x} - \mu_0}{\sigma/\sqrt{n}} = \frac{100 - 120}{30/\sqrt{49}} = \frac{-20}{30/7} = -\frac{14}{3} \approx -4.667. \]
\item \textbf{Decision Rule \& Conclusion:} At $\alpha = 0.05$, the critical value is $z_{0.05} = -1.645$. Since $z = -4.667 < -1.645$, we reject $H_0$.
\end{enumerate}
\ans{$z = -4.67 < -1.645$; reject $H_0$. There is strong statistical evidence at the $5\%$ level to support the claim that mean delivery time is less than $120\text{ minutes}$.}
\end{solution}

%---------------------------------------------------------------- P052
\begin{question}{2023--2024 \textbullet\ Q7(a)}{MATH 243 \textbullet\ 21/05/2026 \textbullet\ P052}{10}
Write short notes on the goodness of fit, one-tailed and two-tailed test, Degrees of freedom, and null-hypothesis and alternative hypothesis.
\end{question}
\begin{solution}
\begin{itemize}
\item \textbf{Goodness of Fit:} A statistical test (most commonly Karl Pearson's $\chi^2 = \sum_{i=1}^{k}\frac{(O_i - E_i)^2}{E_i}$) used to determine how well an observed sample frequency distribution ($O_i$) conforms to a hypothesized theoretical distribution ($E_i = Np_i$, such as Binomial, Poisson, or Normal).
\item \textbf{One-Tailed and Two-Tailed Tests:} A \emph{two-tailed (non-directional) test} tests $H_0: \theta = \theta_0$ against $H_1: \theta \neq \theta_0$, splitting the rejection region of area $\alpha$ equally between both tails ($\alpha/2$ each). A \emph{one-tailed (directional) test} tests against $H_1: \theta > \theta_0$ (right-tailed) or $H_1: \theta < \theta_0$ (left-tailed), placing the entire critical area $\alpha$ in a single tail.
\item \textbf{Degrees of Freedom ($\nu$ or df):} The number of independent observations or categories in a sample that are free to vary after imposing $r$ linear constraints or estimating parameters from the data (e.g.\ $\nu = n - 1$ for a single sample variance $s^2$, $\nu = k - 1 - m$ for a $\chi^2$ goodness-of-fit test with $m$ estimated parameters, and $\nu = (r-1)(c-1)$ for an $r\times c$ contingency table).
\item \textbf{Null Hypothesis ($H_0$) and Alternative Hypothesis ($H_1$ or $H_a$):} The \emph{null hypothesis $H_0$} is the baseline assertion of ``no difference'' or ``no effect'' (e.g.\ $\mu = \mu_0$) assumed true until evidence indicates otherwise. The \emph{alternative hypothesis $H_1$} is the research claim accepted if sample evidence falls in the critical region and leads to rejecting $H_0$.
\end{itemize}
\end{solution}

%---------------------------------------------------------------- P053
\begin{question}{2023--2024 \textbullet\ Q7(b)}{MATH 243 \textbullet\ 21/05/2026 \textbullet\ P053}{12}
A machine puts out 16 imperfect articles in a sample of 500. After the machine is overhauled, it puts out 3 imperfect articles in a batch of 100. Has the machine improved at 1\% level of significance? Given that at 1\% level of significance, $Z = \pm 2.58$.
\end{question}
\begin{solution}
Let $p_1$ and $p_2$ denote the population proportions of imperfect articles before and after overhaul, respectively.
Given $n_1 = 500,\ x_1 = 16 \implies \hat{p}_1 = \dfrac{16}{500} = 0.032$, and $n_2 = 100,\ x_2 = 3 \implies \hat{p}_2 = \dfrac{3}{100} = 0.030$.
\begin{enumerate}[1.]
\item \textbf{Hypotheses:} $H_0: p_1 = p_2$ vs.\ $H_1: p_1 > p_2$ (machine improved, so defect rate decreased).
\item \textbf{Pooled Proportion \& Test Statistic:}
Under $H_0$, the pooled sample proportion is:
\[ \hat{p} = \frac{x_1 + x_2}{n_1 + n_2} = \frac{16 + 3}{500 + 100} = \frac{19}{600} \approx 0.03167, \qquad \hat{q} = 1 - \hat{p} = \frac{581}{600} \approx 0.96833. \]
The two-proportion $z$-statistic is:
\[ z = \frac{\hat{p}_1 - \hat{p}_2}{\sqrt{\hat{p}\hat{q}\left(\frac{1}{n_1} + \frac{1}{n_2}\right)}} = \frac{0.032 - 0.030}{\sqrt{(0.03167)(0.96833)\left(\frac{1}{500} + \frac{1}{100}\right)}} = \frac{0.002}{0.01918} \approx 0.1043. \]
\textit{(Note: Even if treated as a one-sample test of $\hat{p}_2 = 0.030$ against historical $p_0 = 0.032$, $|z| = \frac{|0.030-0.032|}{\sqrt{0.032(0.968)/100}} = 0.1136$.)}
\item \textbf{Decision:} At $\alpha = 0.01$, $|z| = 0.1043 \ll 2.58$ (and $z_{0.01} = 2.33$ one-tailed). Fail to reject $H_0$.
\end{enumerate}
\ans{$z = 0.1043 < 2.58$; fail to reject $H_0$. There is no statistically significant evidence at the $1\%$ level that the machine has improved.}
\end{solution}

%---------------------------------------------------------------- P054
\begin{question}{2023--2024 \textbullet\ Q7(c)}{MATH 243 \textbullet\ 21/05/2026 \textbullet\ P054}{13}
Two independent samples of 8 and 7 items, respectively, had the following values of the variable (weight in ounces). Is the difference between the means of the sample significant? Use 5\% level of significance. Note that for $v = 13, t_{0.05} = 2.16$.
\begin{center}\small
\begin{tabular}{@{}lcccccccc@{}}
\toprule
\textbf{Sample I} & 9 & 11 & 13 & 11 & 15 & 9 & 12 & 14 \\
\textbf{Sample II} & 10 & 12 & 10 & 14 & 9 & 8 & 10 & --- \\
\bottomrule
\end{tabular}
\end{center}
\end{question}
\begin{solution}
\begin{enumerate}[1.]
\item \textbf{Hypotheses:} $H_0: \mu_1 = \mu_2$ vs.\ $H_1: \mu_1 \neq \mu_2$ (two-tailed independent two-sample $t$-test with $\nu = n_1 + n_2 - 2 = 8 + 7 - 2 = 13$).
\item \textbf{Sample Means and Sums of Squares:}
\begin{itemize}
\item \textbf{Sample I ($n_1 = 8$):} $\sum x_{1i} = 94 \implies \bar{x}_1 = \dfrac{94}{8} = 11.75$.
\[ \sum (x_{1i} - \bar{x}_1)^2 = \sum x_{1i}^2 - \frac{(\sum x_{1i})^2}{8} = 1138 - \frac{94^2}{8} = 1138 - 1104.5 = 33.50. \]
\item \textbf{Sample II ($n_2 = 7$):} $\sum x_{2j} = 73 \implies \bar{x}_2 = \dfrac{73}{7} \approx 10.4286$.
\[ \sum (x_{2j} - \bar{x}_2)^2 = \sum x_{2j}^2 - \frac{(\sum x_{2j})^2}{7} = 785 - \frac{5329}{7} = 785 - 761.2857 = 23.7143. \]
\end{itemize}
\item \textbf{Pooled Variance and $t$-Statistic:}
\[ s_p^2 = \frac{\sum (x_{1i} - \bar{x}_1)^2 + \sum (x_{2j} - \bar{x}_2)^2}{n_1 + n_2 - 2} = \frac{33.50 + 23.7143}{13} = \frac{57.2143}{13} \approx 4.4011 \implies s_p \approx 2.0979, \]
\[ t = \frac{\bar{x}_1 - \bar{x}_2}{s_p\sqrt{\frac{1}{n_1} + \frac{1}{n_2}}} = \frac{11.75 - 10.4286}{2.0979\sqrt{\frac{1}{8} + \frac{1}{7}}} = \frac{1.3214}{2.0979 \times 0.51755} = \frac{1.3214}{1.0858} \approx 1.217. \]
\item \textbf{Decision:} Since $|t| = 1.217 < t_{0.05, 13} = 2.16$, we fail to reject $H_0$.
\end{enumerate}
\ans{$t = 1.217 < 2.16$; fail to reject $H_0$. The difference between the sample means is not statistically significant at the $5\%$ level.}
\end{solution}

%---------------------------------------------------------------- P055
\begin{question}{2023--2024 \textbullet\ Q8(a)}{MATH 243 \textbullet\ 21/05/2026 \textbullet\ P055}{15}
A random sample of 200 employees was drawn from a business enterprise to see if there is a relationship between educational attainment and job performance. From the following data, do you find any relationship between educational attainment and the job performance of these employees? Test at 5\% level of significance, where $\chi^2_{0.05, 4} = 9.49$.
\begin{center}\small
\begin{tabular}{@{}lcccc@{}}
\toprule
\textbf{Job performance} & \textbf{Below Primary} & \textbf{College} & \textbf{University} & \textbf{Total} \\
\midrule
Excellent & 10 & 40 & 10 & 60 \\
Good & 30 & 30 & 20 & 80 \\
Fair & 10 & 30 & 20 & 60 \\
\midrule
\textbf{Total} & 50 & 100 & 50 & 200 \\
\bottomrule
\end{tabular}
\end{center}
\end{question}
\begin{solution}
\begin{enumerate}[1.]
\item \textbf{Hypotheses:}
\begin{itemize}
\item $H_0$: Educational attainment and job performance are independent (no relationship).
\item $H_1$: Educational attainment and job performance are associated.
\end{itemize}
\item \textbf{Expected Frequencies ($E_{ij} = \dfrac{R_i \times C_j}{N}$ with $N = 200$):}
\begin{center}\small
\begin{tabular}{@{}lccc@{}}
\toprule
\textbf{Job Performance} & \textbf{Below Primary ($C_1=50$)} & \textbf{College ($C_2=100$)} & \textbf{University ($C_3=50$)} \\
\midrule
Excellent ($R_1=60$) & $O_{11}=10,\ E_{11}=15$ & $O_{12}=40,\ E_{12}=30$ & $O_{13}=10,\ E_{13}=15$ \\
Good ($R_2=80$) & $O_{21}=30,\ E_{21}=20$ & $O_{22}=30,\ E_{22}=40$ & $O_{23}=20,\ E_{23}=20$ \\
Fair ($R_3=60$) & $O_{31}=10,\ E_{31}=15$ & $O_{32}=30,\ E_{32}=30$ & $O_{33}=20,\ E_{33}=15$ \\
\bottomrule
\end{tabular}
\end{center}
\item \textbf{Chi-Square Test Statistic ($\nu = (3-1)(3-1) = 4$):}
\begin{align*}
\chi^2 &= \sum_{i=1}^{3}\sum_{j=1}^{3} \frac{(O_{ij} - E_{ij})^2}{E_{ij}} \\
       &= \frac{(-5)^2}{15} + \frac{10^2}{30} + \frac{(-5)^2}{15} + \frac{10^2}{20} + \frac{(-10)^2}{40} + 0 + \frac{(-5)^2}{15} + 0 + \frac{5^2}{15} \\
       &= \frac{25}{15} + \frac{100}{30} + \frac{25}{15} + \frac{100}{20} + \frac{100}{40} + 0 + \frac{25}{15} + 0 + \frac{25}{15} \\
       &= \frac{5}{3} + \frac{10}{3} + \frac{5}{3} + 5 + 2.5 + \frac{5}{3} + \frac{5}{3} = 10 + 7.5 = 17.50.
\end{align*}
\item \textbf{Decision:} Since $\chi^2 = 17.50 > \chi^2_{0.05, 4} = 9.49$, we reject $H_0$.
\end{enumerate}
\ans{$\chi^2 = 17.50 > 9.49$; reject $H_0$. There is a statistically significant relationship between educational attainment and job performance at the $5\%$ level.}
\end{solution}

%---------------------------------------------------------------- P056
\begin{question}{2022--2023 \textbullet\ Q7(a)}{MATH 243 \textbullet\ 01/02/2025 \textbullet\ P056}{10}
Differentiate the following pair of concepts:
(i) Statistic and parameter, (ii) Critical region and acceptance region, (iii) Exact and inexact hypothesis, (iv) Unidirectional and non-directional test, (v) Level of significance and p-value.
\end{question}
\begin{solution}
\begin{enumerate}[(i)]
\item \textbf{Statistic vs.\ Parameter:} A \emph{parameter} is a fixed (usually unknown) numerical characteristic of an entire population (e.g.\ population mean $\mu$, variance $\sigma^2$, proportion $p$). A \emph{statistic} is a random variable computed solely from sample observations (e.g.\ sample mean $\bar{X}$, sample variance $S^2$) used to estimate or test a parameter.
\item \textbf{Critical Region vs.\ Acceptance Region:} The sample space of a test statistic is partitioned into two regions: the \emph{critical (rejection) region} is the set of extreme test statistic values for which the null hypothesis $H_0$ is rejected, whereas the \emph{acceptance (non-rejection) region} is the complementary set of values for which $H_0$ is retained.
\item \textbf{Exact (Simple) vs.\ Inexact (Composite) Hypothesis:} An \emph{exact (simple) hypothesis} completely specifies the population distribution and fixes the parameter to a single value (e.g.\ $H_0: \mu = 50$). An \emph{inexact (composite) hypothesis} specifies a range of parameter values (e.g.\ $H_1: \mu > 50$ or $H_1: \mu \neq 50$).
\item \textbf{Unidirectional vs.\ Non-Directional Test:} A \emph{unidirectional (one-tailed) test} tests for a deviation in one specific direction ($H_1: \theta > \theta_0$ or $H_1: \theta < \theta_0$), whereas a \emph{non-directional (two-tailed) test} tests for deviations in either direction ($H_1: \theta \neq \theta_0$).
\item \textbf{Level of Significance ($\alpha$) vs.\ p-value:} The \emph{level of significance $\alpha$} is the pre-specified maximum allowable probability of committing a Type-I error ($P(\text{Reject } H_0 \mid H_0\text{ true})$). The \emph{p-value} is the data-driven probability, calculated under $H_0$, of observing a test statistic at least as extreme as the value actually obtained from the sample ($H_0$ is rejected iff $\text{p-value} \le \alpha$).
\end{enumerate}
\end{solution}

%---------------------------------------------------------------- P057 + P064 stacked
\repeated{2}
\begin{question}{2022--2023 \textbullet\ Q7(b)}{MATH 243 \textbullet\ 01/02/2025 \textbullet\ P057}{12}
The mean and standard deviation of GPA scores obtained from a random sample of 40 students of a college were 2.8 and 0.35 respectively. Can we conclude that the sample has come from the entire group of students, which has a mean score of 2.4? Compute 99 percent confidence interval for the mean score in the population. (Given that at 1\% level of significance; $z = \pm 2.58$).
\end{question}
\begin{question}{2018--2019 \textbullet\ Q8(b)}{MATH 245 \textbullet\ 03/11/2019 \textbullet\ P064}{17}
The mean and standard deviation of GPA scores obtained from a random sample of 40 students of a college were 2.8 and 0.35 respectively. Can we conclude that the sample has come from the entire group of students, which has a mean score of 2.4? Use a 1\% level of significance. Compute 99\% confidence interval for the mean score in the population. (Given that, at 1\% level of significance; $z = \pm 2.58$.)
\end{question}
\begin{solution}
Given sample size $n = 40$ (large sample, $n \ge 30$), sample mean $\bar{x} = 2.8$, sample standard deviation $s = 0.35$, hypothesized population mean $\mu_0 = 2.4$, and significance level $\alpha = 0.01$ ($z_{\alpha/2} = \pm 2.58$).
\begin{enumerate}[1.]
\item \textbf{Hypotheses:} $H_0: \mu = 2.4$ vs.\ $H_1: \mu \neq 2.4$ (two-tailed test).
\item \textbf{Standard Error and Test Statistic:}
\[ \text{SE} = \frac{s}{\sqrt{n}} = \frac{0.35}{\sqrt{40}} \approx 0.05534, \qquad z = \frac{\bar{x} - \mu_0}{s/\sqrt{n}} = \frac{2.8 - 2.4}{0.05534} \approx 7.228. \]
\item \textbf{Decision:} Since $|z| = 7.228 > 2.58$, we reject $H_0$. We cannot conclude that the sample comes from a population with mean GPA $2.4$.
\item \textbf{99\% Confidence Interval for the Population Mean $\mu$:}
\[ \bar{x} \pm z_{0.005}\frac{s}{\sqrt{n}} = 2.8 \pm 2.58(0.05534) = 2.8 \pm 0.1428 = [2.6572,\ 2.9428]. \]
\end{enumerate}
\ans{$z = 7.23 > 2.58$; reject $H_0$. The $99\%$ confidence interval for $\mu$ is $[2.657,\ 2.943]$.}
\end{solution}

%---------------------------------------------------------------- P058
\begin{question}{2022--2023 \textbullet\ Q7(c)}{MATH 243 \textbullet\ 01/02/2025 \textbullet\ P058}{13}
The average waiting time in a bank counter to cash a cheque for all customers has been 50 minutes. A new service-providing procedure using modern computer facilities is being tried. If a random sample of 12 customers had an average waiting time for service of 42 minutes with a standard deviation of 11.9 minutes under the new system, test the hypothesis that the population mean is now less than 50, using a significance level of 5 percent assuming the waiting time to be normally distributed. (Given that at $v = 11, t_{0.05} = 1.796$).
\end{question}
\begin{solution}
Given $\mu_0 = 50\text{ min}$, $n = 12$, $\bar{x} = 42\text{ min}$, $s = 11.9\text{ min}$, $\nu = n - 1 = 11$, and $\alpha = 0.05$.
\begin{enumerate}[1.]
\item \textbf{Hypotheses:} $H_0: \mu = 50\text{ min}$ vs.\ $H_1: \mu < 50\text{ min}$ (left-tailed $t$-test).
\item \textbf{Test Statistic:}
\[ t = \frac{\bar{x} - \mu_0}{s/\sqrt{n}} = \frac{42 - 50}{11.9/\sqrt{12}} = \frac{-8}{3.4352} \approx -2.329. \]
\textit{(Note: If $11.9$ is treated as unadjusted sample root-mean-square deviation dividing by $\sqrt{n-1}=\sqrt{11}$, $t = \frac{-8}{11.9/\sqrt{11}} \approx -2.230$.)}
\item \textbf{Decision:} For a left-tailed test at $\alpha = 0.05$ with $\nu = 11$, the critical value is $-t_{0.05, 11} = -1.796$. Since $t = -2.329 < -1.796$, we reject $H_0$.
\end{enumerate}
\ans{$t = -2.329 < -1.796$; reject $H_0$. The population mean waiting time under the new system is significantly less than $50\text{ minutes}$.}
\end{solution}

%---------------------------------------------------------------- P059
\begin{question}{2022--2023 \textbullet\ Q8(a)}{MATH 243 \textbullet\ 01/02/2025 \textbullet\ P059}{10}
A certain drug was administered to 500 people out of a total 800 included in the sample to test its efficacy against typhoid. The results are given below:
\begin{center}\small
\begin{tabular}{@{}lccc@{}}
\toprule
& \textbf{Typhoid} & \textbf{No. Typhoid} & \textbf{Total} \\
\midrule
Drug & 200 & 300 & 500 \\
No. Drug & 280 & 20 & 300 \\
\midrule
\textbf{Total} & 480 & 320 & 800 \\
\bottomrule
\end{tabular}
\end{center}
Based on these data, can it be concluded that the drug is effective in preventing typhoid using 5\% level of significance? (Given that at $v = 1, \chi^2_{0.05} = 3.84$).
\end{question}
\begin{solution}
\begin{enumerate}[1.]
\item \textbf{Hypotheses:} $H_0$: The drug and typhoid incidence are independent (drug is not effective) vs.\ $H_1$: The drug is effective in preventing typhoid.
\item \textbf{Expected Frequencies ($E_{ij} = \dfrac{R_i \times C_j}{800}$):}
\begin{itemize}
\item $E_{11} \ (\text{Drug, Typhoid}) = \dfrac{500 \times 480}{800} = 300$, \qquad $E_{12} \ (\text{Drug, No Typhoid}) = \dfrac{500 \times 320}{800} = 200$,
\item $E_{21} \ (\text{No Drug, Typhoid}) = \dfrac{300 \times 480}{800} = 180$, \quad $E_{22} \ (\text{No Drug, No Typhoid}) = \dfrac{300 \times 320}{800} = 120$.
\end{itemize}
\item \textbf{Chi-Square Test Statistic ($\nu = (2-1)(2-1) = 1$):}
\[ \begin{aligned} \chi^2 &= \sum \frac{(O_{ij} - E_{ij})^2}{E_{ij}} = \frac{(-100)^2}{300} + \frac{100^2}{200} + \frac{100^2}{180} + \frac{(-100)^2}{120} \\ &= 33.33 + 50.00 + 55.56 + 83.33 = 222.22. \end{aligned} \]
\item \textbf{Decision:} Since $\chi^2 = 222.22 \gg \chi^2_{0.05, 1} = 3.84$, we reject $H_0$.
\end{enumerate}
\ans{$\chi^2 = 222.22 > 3.84$; reject $H_0$. The drug is highly effective in preventing typhoid.}
\end{solution}

%---------------------------------------------------------------- P060
\begin{question}{2021--2022 \textbullet\ Q8(a)}{MATH 245 \textbullet\ 30/04/2023 \textbullet\ P060}{17}
Ten individuals are chosen at random from a population and their heights are found to be (in inches) 63, 63, 64, 65, 66, 69, 69, 70, 70, 71. Discuss the suggestion that the mean height of universe is 65. Assume that for 9 degrees of freedom $t$ at 5\% level of significance is 2.262.
\end{question}
\begin{solution}
\begin{enumerate}[1.]
\item \textbf{Hypotheses:} $H_0: \mu = 65\text{ inches}$ vs.\ $H_1: \mu \neq 65\text{ inches}$ ($\alpha = 0.05$, $\nu = n - 1 = 9$).
\item \textbf{Sample Mean and Standard Deviation ($n = 10$):}
\[ \bar{x} = \frac{63 + 63 + 64 + 65 + 66 + 69 + 69 + 70 + 70 + 71}{10} = \frac{670}{10} = 67.0\text{ inches}. \]
Deviations $x_i - \bar{x}$ are $-4, -4, -3, -2, -1, +2, +2, +3, +3, +4$, so:
\[ \sum_{i=1}^{10} (x_i - \bar{x})^2 = 16 + 16 + 9 + 4 + 1 + 4 + 4 + 9 + 9 + 16 = 88, \]
\[ s = \sqrt{\frac{\sum (x_i - \bar{x})^2}{n - 1}} = \sqrt{\frac{88}{9}} \approx 3.1269\text{ inches}. \]
\item \textbf{Test Statistic \& Decision:}
\[ t = \frac{\bar{x} - \mu_0}{s/\sqrt{n}} = \frac{67.0 - 65.0}{3.1269/\sqrt{10}} = \frac{2.0}{0.9888} \approx 2.0226. \]
Since $|t| = 2.0226 < t_{0.05, 9} = 2.262$, we fail to reject $H_0$.
\end{enumerate}
\ans{$t = 2.023 < 2.262$; fail to reject $H_0$. The sample is consistent with the suggestion that the population mean height is $65\text{ inches}$.}
\end{solution}

%---------------------------------------------------------------- P061
\begin{question}{2020--2021 \textbullet\ Q8(b)}{MATH 245 \textbullet\ 2020--2021 \textbullet\ P061}{15}
The human requirement for salt is 220 milligrams per day. If a random sample of 20 similar servings of a certain cereal has a mean sodium content of 244 milligrams and a standard deviation of 24.5 milligrams, does this suggest at the 0.05 level of significance that the average sodium content for a single serving of such cereal is greater than 220 milligrams? (Necessary Table 3 is attached)
\end{question}
\begin{solution}
Given $\mu_0 = 220\text{ mg}$, $n = 20$, $\bar{x} = 244\text{ mg}$, $s = 24.5\text{ mg}$, $\nu = n - 1 = 19$, and $\alpha = 0.05$.
\begin{enumerate}[1.]
\item \textbf{Hypotheses:} $H_0: \mu = 220\text{ mg}$ vs.\ $H_1: \mu > 220\text{ mg}$ (right-tailed $t$-test).
\item \textbf{Test Statistic:}
\[ t = \frac{\bar{x} - \mu_0}{s/\sqrt{n}} = \frac{244 - 220}{24.5/\sqrt{20}} = \frac{24}{5.4784} \approx 4.381. \]
\item \textbf{Decision:} From Table 3 (Student's $t$-distribution), for $\nu = 19$ and one-tailed $\alpha = 0.05$, the critical value is $t_{0.05, 19} = 1.729$. Since $t = 4.381 > 1.729$, we reject $H_0$.
\end{enumerate}
\ans{$t = 4.381 > 1.729$; reject $H_0$. The data strongly suggest that the average sodium content per serving exceeds $220\text{ mg}$.}
\end{solution}

%---------------------------------------------------------------- P062
\begin{question}{2020--2021 \textbullet\ Q8(c)}{MATH 245 \textbullet\ 2020--2021 \textbullet\ P062}{10}
In a DNA paternity testing procedure where the child, mother, and putative father are tested, the test is 100\% accurate if the man is not the father and 99.9\% accurate if he is. Consider using the results of DNA paternity testing to decide between the following two hypotheses:
\begin{itemize}
\item $H_0$ : a particular man is the father
\item $H_a$ : a particular man is not the father
\end{itemize}
In the context of this problem, describe Type I and Type II errors.
\end{question}
\begin{solution}
\begin{itemize}
\item \textbf{Type I Error (Rejecting $H_0$ when $H_0$ is true):}
Concluding that the particular man is \textbf{not} the father when in fact he \textbf{is} the true biological father. Since the test is $99.9\%$ accurate when the man is the father, the probability of a Type I error is $\alpha = 1 - 0.999 = 0.001$ ($0.1\%$).
\item \textbf{Type II Error (Failing to reject $H_0$ when $H_a$ is true):}
Concluding that the particular man \textbf{is} the father when in fact he is \textbf{not} the father. Since the test is $100\%$ accurate when the man is not the father, the probability of a Type II error is $\beta = 1 - 1.00 = 0$.
\end{itemize}
\ans{Type I error: declaring the true biological father is not the father ($\alpha = 0.1\%$). Type II error: failing to exclude a man who is not the father ($\beta = 0\%$).}
\end{solution}

%---------------------------------------------------------------- P063
\begin{question}{2019--2020 \textbullet\ Q6(b)}{MATH 245 \textbullet\ 25/09/2021 \textbullet\ P063}{18}
A medical student claims that mice with an average life span of 32 months will live to be about 40 months old when 40\% of the calories in their diet are replaced by vitamins and protein. Is there any reason to believe that $\mu < 40$ if 64 mice that are placed on this diet have an average life of 38 months with a standard deviation of 5.8 months? Use a P-value in your conclusion.
\end{question}
\begin{solution}
Given $\mu_0 = 40\text{ months}$, $n = 64$ (large sample), $\bar{x} = 38\text{ months}$, and $s = 5.8\text{ months}$:
\begin{enumerate}[1.]
\item \textbf{Hypotheses:} $H_0: \mu = 40\text{ months}$ vs.\ $H_1: \mu < 40\text{ months}$ (left-tailed test).
\item \textbf{Test Statistic:}
\[ z = \frac{\bar{x} - \mu_0}{s/\sqrt{n}} = \frac{38 - 40}{5.8/\sqrt{64}} = \frac{-2}{5.8/8} = \frac{-16}{5.8} \approx -2.7586 \ (\approx -2.76). \]
\item \textbf{P-value \& Conclusion:}
The left-tailed P-value is:
\[ \text{P-value} = P(Z \le -2.76) = \Phi(-2.76) \approx 0.0029. \]
Since $\text{P-value} = 0.0029 < 0.01$ (and $< 0.05$), the result is highly statistically significant, so we reject $H_0$.
\end{enumerate}
\ans{$z = -2.76$, $\text{P-value} = 0.0029 < 0.01$; reject $H_0$. There is strong evidence to believe that $\mu < 40\text{ months}$.}
\end{solution}

%---------------------------------------------------------------- P065
\begin{question}{2016--2017 \textbullet\ Q8(b)}{MATH 245 \textbullet\ 05/08/2017 \textbullet\ P065}{15}
A random sample of 100 recorded deaths in Bangladesh during the past year showed an average life span of 61.8 years. Assuming a population standard deviation of 4.9 years, does this seem to indicate that the mean life span today is greater than 60 years? Use a 0.05 level of significance. (Necessary table attached)
\end{question}
\begin{solution}
Given $\mu_0 = 60\text{ years}$, $\sigma = 4.9\text{ years}$, $n = 100$, $\bar{x} = 61.8\text{ years}$, and $\alpha = 0.05$:
\begin{enumerate}[1.]
\item \textbf{Hypotheses:} $H_0: \mu = 60\text{ years}$ vs.\ $H_1: \mu > 60\text{ years}$ (right-tailed $z$-test).
\item \textbf{Test Statistic:}
\[ z = \frac{\bar{x} - \mu_0}{\sigma/\sqrt{n}} = \frac{61.8 - 60}{4.9/\sqrt{100}} = \frac{1.8}{0.49} \approx 3.6735. \]
\item \textbf{Decision:} At $\alpha = 0.05$, the critical value is $z_{0.05} = 1.645$. Since $z = 3.67 > 1.645$ (with $\text{P-value} = 1 - \Phi(3.67) \approx 0.00012$), we reject $H_0$.
\end{enumerate}
\ans{$z = 3.67 > 1.645$; reject $H_0$. The data indicate that the mean life span today is greater than $60\text{ years}$.}
\end{solution}

%---------------------------------------------------------------- P066
\begin{question}{2015--2016 \textbullet\ Q8(c)}{MATH 241 \textbullet\ 07/08/2016 \textbullet\ P066}{10}
The mean weekly sale of a chocolate bar company in candy stores was 153.7 bars per store. After an advertising campaign, the mean weekly sale in 29 stores for a typical week increased to 169.4 and showed a standard deviation of 19.7. Was the advertising successful? Use a 5\% level of significance. (Necessary chart 2 is attached)
\end{question}
\begin{solution}
Given $\mu_0 = 153.7$, $n = 29$, $\bar{x} = 169.4$, $s = 19.7$, $\nu = n - 1 = 28$, and $\alpha = 0.05$:
\begin{enumerate}[1.]
\item \textbf{Hypotheses:} $H_0: \mu = 153.7$ vs.\ $H_1: \mu > 153.7$ (right-tailed $t$-test).
\item \textbf{Test Statistic:}
\[ t = \frac{\bar{x} - \mu_0}{s/\sqrt{n}} = \frac{169.4 - 153.7}{19.7/\sqrt{29}} = \frac{15.7}{3.6582} \approx 4.292. \]
\item \textbf{Decision:} From Chart 2 (for $\nu = 28$), $t_{0.05, 28} = 1.701$ (one-tailed) and $t_{0.025, 28} = 2.048$ (two-tailed). Since $t = 4.292 > 1.701$, we reject $H_0$.
\end{enumerate}
\ans{$t = 4.292 > 1.701$; reject $H_0$. The advertising campaign was successful in increasing sales.}
\end{solution}

%---------------------------------------------------------------- P067
\begin{question}{2014--2015 \textbullet\ Q7(b)}{MATH 241 \textbullet\ 06/02/2016 \textbullet\ P067}{15}
A bulb manufacturing company claims that the average longevity of their bulb is 4 years with a standard deviation 0.16 years. A random sample of 10 bulbs gave mean longevity of 3.45 years. Does the sample mean justify the claim of the manufacturer? Use a 5 percent level of significance. (Necessary Chart 1 is attached)
\end{question}
\begin{solution}
Given $\mu_0 = 4\text{ years}$, $\sigma = 0.16\text{ years}$ (or sample $s = 0.16$), $n = 10$, $\bar{x} = 3.45\text{ years}$, and $\alpha = 0.05$ (Chart 1 is Student's $t$-table with $\nu = 9$, where $t_{0.025, 9} = 2.26$ and $t_{0.05, 9} = 1.83$):
\begin{enumerate}[1.]
\item \textbf{Hypotheses:} $H_0: \mu = 4.0\text{ years}$ vs.\ $H_1: \mu \neq 4.0\text{ years}$ (or $H_1: \mu < 4.0$).
\item \textbf{Test Statistic:}
\[ t = \frac{\bar{x} - \mu_0}{\sigma/\sqrt{n}} = \frac{3.45 - 4.00}{0.16/\sqrt{10}} = \frac{-0.55}{0.05060} \approx -10.870. \]
\item \textbf{Decision:} Since $|t| = 10.870 \gg t_{0.025, 9} = 2.26$ (and also $\gg z_{0.025} = 1.96$ if $0.16$ is treated as the known population $\sigma$), we strongly reject $H_0$.
\end{enumerate}
\ans{$t = -10.87$; reject $H_0$. The sample mean of $3.45\text{ years}$ does not justify the manufacturer's claim.}
\end{solution}

%---------------------------------------------------------------- P068
\begin{question}{2013--2014 \textbullet\ Q8(a)}{MATH 241 \textbullet\ 13/12/2014 \textbullet\ P068}{20}
Talcum powder is packed into tins by a machine. A random sample of 11 tins is drawn and their contents are found to weigh (in lbs) as follows:
0.44, 0.51, 0.49, 0.52, 0.45, 0.48, 0.46, 0.45, 0.47, 0.45 and 0.47.
Test whether the average packing can be taken to be 0.5 lbs at 5\% level of significance. (Given: $t = \pm 2.23$)
\end{question}
\begin{solution}
\begin{enumerate}[1.]
\item \textbf{Hypotheses:} $H_0: \mu = 0.50\text{ lbs}$ vs.\ $H_1: \mu \neq 0.50\text{ lbs}$ ($\alpha = 0.05$, $\nu = n - 1 = 10$, critical value $t_{0.025, 10} = \pm 2.23$).
\item \textbf{Sample Mean and Standard Deviation ($n = 11$):}
\[ \bar{x} = \frac{\sum_{i=1}^{11} x_i}{11} = \frac{5.19}{11} \approx 0.47182\text{ lbs}. \]
Computing $\sum (x_i - \bar{x})^2 = \sum x_i^2 - \frac{(\sum x_i)^2}{11} = 2.4555 - \frac{26.9361}{11} = 2.4555 - 2.448736 = 0.006764$:
\[ s = \sqrt{\frac{0.006764}{10}} \approx 0.02601\text{ lbs}. \]
\item \textbf{Test Statistic \& Decision:}
\[ t = \frac{\bar{x} - \mu_0}{s/\sqrt{n}} = \frac{0.47182 - 0.50}{0.02601/\sqrt{11}} = \frac{-0.02818}{0.007842} \approx -3.594. \]
Since $|t| = 3.594 > 2.23$, we reject $H_0$.
\end{enumerate}
\ans{$t = -3.594$ ($|t| > 2.23$); reject $H_0$. The average packing weight cannot be taken to be $0.5\text{ lbs}$.}
\end{solution}

%---------------------------------------------------------------- P069 + P070 stacked
\repeated{2}
\begin{question}{2013--2014 \textbullet\ Q8(b)}{MATH 241 \textbullet\ 13/12/2014 \textbullet\ P069}{15}
A survey of 320 families with 5 children each revealed the following distribution:
\begin{center}\small
\begin{tabular}{@{}lcccccc@{}}
\toprule
\textbf{No. of boys and girls} & 5B, 0G & 4B, 1G & 3B, 2G & 2B, 3G & 1B, 4G & 0B, 5G \\
\textbf{No. of families} & 18 & 56 & 110 & 88 & 40 & 8 \\
\bottomrule
\end{tabular}
\end{center}
Is the result consistent with the hypothesis that male and female births are equally probable? Use a 5\% level of significance. (Given for $v = 5, \chi^2 = 11.07$)
\end{question}
\begin{question}{2012--2013 \textbullet\ Q7(b)}{MATH 241 \textbullet\ 02/06/2014 \textbullet\ P070}{17}
A survey of 640 families with 5 children each revealed the following distribution:
\begin{center}\small
\begin{tabular}{@{}lcccccc@{}}
\toprule
\textbf{No. of boys / girls} & (5, 0) & (4, 1) & (3, 2) & (2, 3) & (1, 4) & (0, 5) \\
\textbf{No. of families} & 28 & 112 & 220 & 176 & 80 & 24 \\
\bottomrule
\end{tabular}
\end{center}
Is this result consistent with the hypothesis that male and female births are equally probable? (use a 5\% level of significance, given that for $v = 5, \chi^2_{0.05} = 11.07$)
\end{question}
\begin{solution}
Under $H_0: p = q = \dfrac{1}{2}$ (male and female births are equally likely), the probability of having $x$ boys in a family of 5 children is $P(X = x) = \binom{5}{x}\left(\frac{1}{2}\right)^5 = \dfrac{\binom{5}{x}}{32}$ for $x = 5, 4, 3, 2, 1, 0$, giving theoretical proportions $\dfrac{1}{32}, \dfrac{5}{32}, \dfrac{10}{32}, \dfrac{10}{32}, \dfrac{5}{32}, \dfrac{1}{32}$.

\begin{center}\small
\begin{tabular}{@{}ccccccc@{}}
\toprule
\textbf{Boys $(x)$} & \textbf{5} & \textbf{4} & \textbf{3} & \textbf{2} & \textbf{1} & \textbf{0} \\
\midrule
\textbf{2013--14 ($N=320$):} Observed $O_i$ & $18$ & $56$ & $110$ & $88$ & $40$ & $8$ \\
Expected $E_i = 320\binom{5}{x}/32$ & $10$ & $50$ & $100$ & $100$ & $50$ & $10$ \\
$(O_i - E_i)^2 / E_i$ & $6.40$ & $0.72$ & $1.00$ & $1.44$ & $2.00$ & $0.40$ \\
\midrule
\textbf{2012--13 ($N=640$):} Observed $O_i$ & $28$ & $112$ & $220$ & $176$ & $80$ & $24$ \\
Expected $E_i = 640\binom{5}{x}/32$ & $20$ & $100$ & $200$ & $200$ & $100$ & $20$ \\
$(O_i - E_i)^2 / E_i$ & $3.20$ & $1.44$ & $2.00$ & $2.88$ & $4.00$ & $0.80$ \\
\bottomrule
\end{tabular}
\end{center}
\begin{itemize}
\item \textbf{For 2013--2014 ($N = 320$):} $\chi^2 = 6.40 + 0.72 + 1.00 + 1.44 + 2.00 + 0.40 = \mathbf{11.96}$. Since $11.96 > 11.07$, we reject $H_0$.
\item \textbf{For 2012--2013 ($N = 640$):} $\chi^2 = 3.20 + 1.44 + 2.00 + 2.88 + 4.00 + 0.80 = \mathbf{14.32}$. Since $14.32 > 11.07$, we reject $H_0$.
\end{itemize}
\ans{For $N=320$: $\chi^2 = 11.96 > 11.07$ (reject $H_0$). For $N=640$: $\chi^2 = 14.32 > 11.07$ (reject $H_0$). Neither survey is consistent with equal birth probabilities at the $5\%$ level.}
\end{solution}

%---------------------------------------------------------------- P071
\begin{question}{2012--2013 \textbullet\ Q8(b)}{MATH 241 \textbullet\ 02/06/2014 \textbullet\ P071}{17}
In a survey of buying habits, 500 women shoppers are chosen at random in supermarket A; their average weekly food expenditure is Tk. 500 with a standard deviation of Tk. 80. For another group of 500 women shoppers chosen at random in supermarket B located in another area of the same city, the average weekly food expenditure is Tk. 440 with a standard deviation of Tk. 90. Test at 10\% level of significance whether the average weekly food expenditures of the populations of women shoppers are equal.
\end{question}
\begin{solution}
Given $n_1 = 500,\ \bar{x}_1 = 500,\ s_1 = 80$ and $n_2 = 500,\ \bar{x}_2 = 440,\ s_2 = 90$, with $\alpha = 0.10$:
\begin{enumerate}[1.]
\item \textbf{Hypotheses:} $H_0: \mu_1 = \mu_2$ vs.\ $H_1: \mu_1 \neq \mu_2$ (two-tailed large-sample $z$-test).
\item \textbf{Standard Error and Test Statistic:}
\[ \text{SE} = \sqrt{\frac{s_1^2}{n_1} + \frac{s_2^2}{n_2}} = \sqrt{\frac{80^2}{500} + \frac{90^2}{500}} = \sqrt{\frac{6400 + 8100}{500}} = \sqrt{29} \approx 5.3852, \]
\[ z = \frac{\bar{x}_1 - \bar{x}_2}{\text{SE}} = \frac{500 - 440}{5.3852} = \frac{60}{5.3852} \approx 11.142. \]
\item \textbf{Decision:} At $\alpha = 0.10$ (two-tailed), the critical value is $z_{0.05} = \pm 1.645$. Since $z = 11.142 \gg 1.645$, we reject $H_0$.
\end{enumerate}
\ans{$z = 11.14 > 1.645$; reject $H_0$. The population average weekly food expenditures differ significantly.}
\end{solution}

%---------------------------------------------------------------- P072
\begin{question}{2011--2012 \textbullet\ Q8(a)}{MATH 241 \textbullet\ 07/01/2013 \textbullet\ P072}{17}
Explain type-I error and type-II error. 9 items of a sample had the following values:
50, 48, 49, 53, 52, 47, 45, 46, 51.
Can this sample be regarded as taken from the population having 47 as mean at 5\% level of significance? (Given that for $v = 8, t_{0.05} = 3.355$)
\end{question}
\begin{solution}
\begin{note}
The printed critical value $3.355$ is actually $t_{0.005,8}$ (two-tailed $1\%$ level); the two-tailed $5\%$ value for $\nu = 8$ is $t_{0.025,8} = 2.306$. The conclusion is the same with either value.
\end{note}
\textbf{Part 1: Type-I and Type-II Errors.}
\begin{itemize}
\item \textbf{Type-I Error ($\alpha$):} Rejecting the null hypothesis $H_0$ when $H_0$ is actually true (false positive; $\alpha = P(\text{Reject } H_0 \mid H_0\text{ true})$).
\item \textbf{Type-II Error ($\beta$):} Failing to reject the null hypothesis $H_0$ when $H_0$ is actually false (false negative; $\beta = P(\text{Accept } H_0 \mid H_1\text{ true})$).
\end{itemize}

\medskip
\textbf{Part 2: One-Sample $t$-Test ($n = 9,\ \mu_0 = 47,\ \nu = 8$).}
\begin{enumerate}[1.]
\item \textbf{Hypotheses:} $H_0: \mu = 47$ vs.\ $H_1: \mu \neq 47$.
\item \textbf{Sample Mean and Standard Deviation:}
\[ \bar{x} = \frac{50 + 48 + 49 + 53 + 52 + 47 + 45 + 46 + 51}{9} = \frac{441}{9} = 49.0. \]
Deviations $x_i - \bar{x}$ are $1, -1, 0, 4, 3, -2, -4, -3, 2$, so $\sum (x_i - \bar{x})^2 = 1 + 1 + 0 + 16 + 9 + 4 + 16 + 9 + 4 = 60$:
\[ s = \sqrt{\frac{60}{9 - 1}} = \sqrt{7.5} \approx 2.7386. \]
\item \textbf{Test Statistic \& Decision:}
\[ t = \frac{\bar{x} - \mu_0}{s/\sqrt{n}} = \frac{49.0 - 47.0}{2.7386/\sqrt{9}} = \frac{2.0}{0.91287} \approx 2.1909. \]
Since $|t| = 2.191 < 3.355$ (and also less than the standard $5\%$ two-tailed critical value $t_{0.025, 8} = 2.306$), we fail to reject $H_0$.
\end{enumerate}
\ans{$t = 2.191 < 3.355$; fail to reject $H_0$. The sample can be regarded as taken from a population with mean $47$.}
\end{solution}

%---------------------------------------------------------------- P073
\begin{question}{2011--2012 \textbullet\ Q8(b)}{MATH 241 \textbullet\ 07/01/2013 \textbullet\ P073}{18}
5 coins are tossed 3200 times and the number of heads appearing each time are noted. At the end, the following results were obtained:
\begin{center}\small
\begin{tabular}{@{}lcccccc@{}}
\toprule
\textbf{No. of heads} & 0 & 1 & 2 & 3 & 4 & 5 \\
\textbf{Frequency} & 80 & 570 & 1100 & 900 & 500 & 50 \\
\bottomrule
\end{tabular}
\end{center}
Test the hypothesis that the coins are unbiased at 5\% level of significance. (Given that for $v = 5, \chi^2_{0.05} = 11.07$)
\end{question}
\begin{solution}
\begin{enumerate}[1.]
\item \textbf{Hypotheses:} $H_0$: The 5 coins are unbiased ($p = \frac{1}{2}$) vs.\ $H_1$: The coins are biased.
\item \textbf{Expected Frequencies ($E_x = 3200\binom{5}{x}\left(\frac{1}{2}\right)^5 = 100\binom{5}{x}$):}
\begin{center}\small
\begin{tabular}{@{}lcccccc@{}}
\toprule
\textbf{No. of heads ($x$)} & 0 & 1 & 2 & 3 & 4 & 5 \\
\midrule
Observed $O_x$ & $80$ & $570$ & $1100$ & $900$ & $500$ & $50$ \\
Expected $E_x$ & $100$ & $500$ & $1000$ & $1000$ & $500$ & $100$ \\
$(O_x - E_x)^2 / E_x$ & $4.00$ & $9.80$ & $10.00$ & $10.00$ & $0.00$ & $25.00$ \\
\bottomrule
\end{tabular}
\end{center}
\item \textbf{Chi-Square Statistic \& Decision:}
\[ \chi^2 = \sum_{x=0}^{5} \frac{(O_x - E_x)^2}{E_x} = 4.00 + 9.80 + 10.00 + 10.00 + 0.00 + 25.00 = 58.80. \]
Since $\chi^2 = 58.80 > \chi^2_{0.05, 5} = 11.07$, we reject $H_0$.
\end{enumerate}
\ans{$\chi^2 = 58.80 > 11.07$; reject $H_0$. The coins are not unbiased.}
\end{solution}

%===================================================================== TOPIC 11
\topic[Analysis of Variance]{11}{Analysis of Variance (One-Way \& Two-Way ANOVA)}

%---------------------------------------------------------------- P074
\begin{question}{2023--2024 \textbullet\ Q8(b)}{MATH 243 \textbullet\ 21/05/2026 \textbullet\ P074}{20}
The following table gives the retail prices of a commodity in some shops selected in three big cities. Carry out the analysis of variance to determine whether there is a significant difference in the prices charged in these cities at the 5\% level of significance, with $v_1 = 2$ and $v_2 = 9$, and $F_{0.05} = 4.26$.
\begin{center}\small
\begin{tabular}{@{}lcccc@{}}
\toprule
\textbf{City} & \multicolumn{4}{c}{\textbf{Price per unit (TK)}} \\
\midrule
A & 15 & 18 & 23 & 24 \\
B & 23 & 19 & 18 & 24 \\
C & 19 & 16 & 22 & 31 \\
\bottomrule
\end{tabular}
\end{center}
\end{question}
\begin{solution}
We have $k = 3$ cities (treatments), each with $n_i = 4$ observations ($N = 12$ total).
\begin{enumerate}[1.]
\item \textbf{Hypotheses:} $H_0: \mu_A = \mu_B = \mu_C$ vs.\ $H_1$: At least two city mean prices differ.
\item \textbf{Group Means and Grand Mean:}
\begin{itemize}
\item City A: $T_A = 15 + 18 + 23 + 24 = 80 \implies \bar{x}_A = \dfrac{80}{4} = 20.0$.
\item City B: $T_B = 23 + 19 + 18 + 24 = 84 \implies \bar{x}_B = \dfrac{84}{4} = 21.0$.
\item City C: $T_C = 19 + 16 + 22 + 31 = 88 \implies \bar{x}_C = \dfrac{88}{4} = 22.0$.
\item Grand Total $G = 80 + 84 + 88 = 252 \implies \text{Grand Mean } \bar{\bar{x}} = \dfrac{252}{12} = 21.0$.
\end{itemize}
\item \textbf{Sums of Squares ($SST, SSB, SSW$):}
\begin{align*}
SSB &= \sum_{i=1}^{3} n_i(\bar{x}_i - \bar{\bar{x}})^2 = 4\big[(20 - 21)^2 + (21 - 21)^2 + (22 - 21)^2\big] = 4(1 + 0 + 1) = 8.0, \\[4pt]
SSW &= \sum_{i=1}^{3}\sum_{j=1}^{4} (x_{ij} - \bar{x}_i)^2 \\
    &= \big[(-5)^2 + (-2)^2 + 3^2 + 4^2\big] + \big[2^2 + (-2)^2 + (-3)^2 + 3^2\big] + \big[(-3)^2 + (-6)^2 + 0^2 + 9^2\big] \\
    &= 54 + 26 + 126 = 206.0, \\[4pt]
SST &= SSB + SSW = 8.0 + 206.0 = 214.0.
\end{align*}
\item \textbf{One-Way ANOVA Table:}
\begin{center}\small
\begin{tabular}{@{}lcccc@{}}
\toprule
\textbf{Source} & \textbf{SS} & \textbf{df ($\nu$)} & \textbf{MS} & \textbf{$F$-Ratio} \\
\midrule
Between Cities & $SSB = 8.0$ & $k - 1 = 2$ & $MSB = 8.0/2 = 4.000$ & $F = \dfrac{4.000}{22.889} = 0.1748$ \\
Within Cities (Error) & $SSW = 206.0$ & $N - k = 9$ & $MSW = 206.0/9 = 22.889$ & \\
\midrule
\textbf{Total} & $SST = 214.0$ & $N - 1 = 11$ & & \\
\bottomrule
\end{tabular}
\end{center}
\item \textbf{Decision:} Since $F = 0.1748 < F_{0.05(2,9)} = 4.26$, we fail to reject $H_0$.
\end{enumerate}
\ans{$F = 0.1748 < 4.26$; fail to reject $H_0$. There is no significant difference in commodity prices across the three cities.}
\end{solution}

%---------------------------------------------------------------- P075
\begin{question}{2022--2023 \textbullet\ Q8(b)}{MATH 243 \textbullet\ 01/02/2025 \textbullet\ P075}{10}
Suppose you test the lifetime of three electric bulbs of each of four brands (measured in hundreds of hours) as shown below. Can we infer that the mean lifetimes of the four brands of electric bulbs are equal with 5\% and 1\% level of significance? (Necessary charts are attached).
\begin{center}\small
\begin{tabular}{@{}cccc@{}}
\toprule
\textbf{Brand A} & \textbf{Brand B} & \textbf{Brand C} & \textbf{Brand D} \\
\midrule
20 & 25 & 24 & 23 \\
19 & 23 & 20 & 20 \\
21 & 21 & 22 & 20 \\
\bottomrule
\end{tabular}
\end{center}
\end{question}
\begin{solution}
Here $k = 4$ brands, $n_i = 3$ bulbs per brand, and $N = 12$ bulbs total.
\begin{enumerate}[1.]
\item \textbf{Hypotheses:} $H_0: \mu_A = \mu_B = \mu_C = \mu_D$ vs.\ $H_1$: Not all brand means are equal.
\item \textbf{Brand Means and Grand Mean:}
\begin{itemize}
\item $\bar{x}_A = \dfrac{20+19+21}{3} = 20.0$, \quad $\bar{x}_B = \dfrac{25+23+21}{3} = 23.0$, \quad $\bar{x}_C = \dfrac{24+20+22}{3} = 22.0$, \quad $\bar{x}_D = \dfrac{23+20+20}{3} = 21.0$.
\item Grand Mean: $\bar{\bar{x}} = \dfrac{20 + 23 + 22 + 21}{4} = 21.5$.
\end{itemize}
\item \textbf{Sums of Squares:}
\begin{align*}
SSB &= 3\big[(-1.5)^2 + (1.5)^2 + (0.5)^2 + (-0.5)^2\big] = 3(2.25 + 2.25 + 0.25 + 0.25) = 15.0, \\[4pt]
SSW &= (0+1+1) + (4+0+4) + (4+4+0) + (4+1+1) = 2 + 8 + 8 + 6 = 24.0, \\[4pt]
SST &= 15.0 + 24.0 = 39.0.
\end{align*}
\item \textbf{One-Way ANOVA Table ($\nu_1 = 3,\ \nu_2 = 8$):}
\begin{center}\small
\begin{tabular}{@{}lcccc@{}}
\toprule
\textbf{Source} & \textbf{SS} & \textbf{df} & \textbf{MS} & \textbf{$F$-Ratio} \\
\midrule
Between Brands & $15.0$ & $3$ & $MSB = 5.0$ & $F = \dfrac{5.0}{3.0} = 1.6667$ \\
Within Brands (Error) & $24.0$ & $8$ & $MSW = 3.0$ & \\
\midrule
\textbf{Total} & $39.0$ & $11$ & & \\
\bottomrule
\end{tabular}
\end{center}
\item \textbf{Decision:} From the attached $F$-tables for $(\nu_1 = 3, \nu_2 = 8)$, $F_{0.05(3,8)} = 4.07$ and $F_{0.01(3,8)} = 7.59$. Since $F = 1.67 < 4.07 < 7.59$, we fail to reject $H_0$ at both the $5\%$ and $1\%$ levels.
\end{enumerate}
\ans{$F = 1.67$ (critical values: $F_{0.05(3,8)} = 4.07$, $F_{0.01(3,8)} = 7.59$); fail to reject $H_0$ at both $5\%$ and $1\%$ significance levels.}
\end{solution}

%---------------------------------------------------------------- P076
\begin{question}{2022--2023 \textbullet\ Q8(c)}{MATH 243 \textbullet\ 01/02/2025 \textbullet\ P076}{15}
Seeds of four different types of corn are planted in five blocks. Each block is divided into four plots, which are then randomly assigned to the four types of corn. Determine at the 0.05 significance level whether the yields in bushels per acre, as shown in the following Table, vary significantly with differences in (i) the soil (i.e., the five blocks) and (ii) the type of corn. (Necessary charts are attached).
\begin{center}\small
\begin{tabular}{@{}lcccc@{}}
\toprule
& \textbf{Type I} & \textbf{Type II} & \textbf{Type III} & \textbf{Type IV} \\
\midrule
Block A & 12 & 15 & 10 & 14 \\
Block B & 15 & 19 & 12 & 11 \\
Block C & 14 & 18 & 15 & 12 \\
Block D & 11 & 16 & 12 & 16 \\
Block E & 16 & 17 & 11 & 14 \\
\bottomrule
\end{tabular}
\end{center}
\end{question}
\begin{solution}
This is a Randomized Complete Block Design (two-way ANOVA without interaction) with $b = 5$ blocks (rows) and $k = 4$ corn types (columns), giving $N = bk = 20$ observations.
\begin{enumerate}[1.]
\item \textbf{Row (Block) Means, Column (Treatment) Means, and Grand Mean:}
\begin{itemize}
\item \textbf{Block Totals ($R_i$) \& Means ($\bar{x}_{i\cdot} = R_i/4$):}
Block A: $51\ (12.75)$;\quad Block B: $57\ (14.25)$;\quad Block C: $59\ (14.75)$;\quad Block D: $55\ (13.75)$;\quad Block E: $58\ (14.50)$.
\item \textbf{Corn Type Totals ($C_j$) \& Means ($\bar{x}_{\cdot j} = C_j/5$):}
Type I: $68\ (13.60)$;\quad Type II: $85\ (17.00)$;\quad Type III: $60\ (12.00)$;\quad Type IV: $67\ (13.40)$.
\item \textbf{Grand Total \& Grand Mean:} $G = 280 \implies \bar{\bar{x}} = \dfrac{280}{20} = 14.0$.
\end{itemize}
\item \textbf{Sums of Squares:}
\begin{align*}
SST &= \sum_{i=1}^{5}\sum_{j=1}^{4} (x_{ij} - 14.0)^2 = \sum x_{ij}^2 - \frac{G^2}{20} = 4044 - \frac{280^2}{20} = 4044 - 3920 = 124.0, \\[4pt]
SS_{\text{Blocks}} &= 4\sum_{i=1}^{5} (\bar{x}_{i\cdot} - 14.0)^2 = \frac{51^2 + 57^2 + 59^2 + 55^2 + 58^2}{4} - 3920 = 3930 - 3920 = 10.0, \\[4pt]
SS_{\text{Corn}} &= 5\sum_{j=1}^{4} (\bar{x}_{\cdot j} - 14.0)^2 = \frac{68^2 + 85^2 + 60^2 + 67^2}{5} - 3920 = 3987.6 - 3920 = 67.6, \\[4pt]
SSE &= SST - SS_{\text{Blocks}} - SS_{\text{Corn}} = 124.0 - 10.0 - 67.6 = 46.4.
\end{align*}
\item \textbf{Two-Way ANOVA Table:}
\begin{center}\small
\begin{tabular}{@{}lccccc@{}}
\toprule
\textbf{Source of Variation} & \textbf{SS} & \textbf{df} & \textbf{MS} & \textbf{$F$-Ratio} & \textbf{Critical $F_{0.05}$} \\
\midrule
Blocks (Soil) & $10.0$ & $b - 1 = 4$ & $2.5000$ & $F_B = \dfrac{2.5000}{3.8667} = 0.6466$ & $F_{0.05(4,12)} = 3.26$ \\
Corn Types & $67.6$ & $k - 1 = 3$ & $22.5333$ & $F_C = \dfrac{22.5333}{3.8667} = 5.8276$ & $F_{0.05(3,12)} = 3.49$ \\
Error (Residual) & $46.4$ & $(b-1)(k-1) = 12$ & $3.8667$ & & \\
\midrule
\textbf{Total} & $124.0$ & $N - 1 = 19$ & & & \\
\bottomrule
\end{tabular}
\end{center}
\item \textbf{Conclusions at $\alpha = 0.05$:}
\begin{enumerate}[(i)]
\item \textbf{Soil (Blocks):} Since $F_B = 0.65 < F_{0.05(4,12)} = 3.26$, fail to reject $H_{0B}$ (yields do not vary significantly across soil blocks).
\item \textbf{Type of Corn:} Since $F_C = 5.83 > F_{0.05(3,12)} = 3.49$, reject $H_{0C}$ (yields vary significantly among the four types of corn).
\end{enumerate}
\end{enumerate}
\ans{(i) Soil blocks: $F = 0.65 < 3.26$ (not significant); \qquad (ii) Corn types: $F = 5.83 > 3.49$ (significant difference at $\alpha = 0.05$).}
\end{solution}

%---------------------------------------------------------------- P077
\begin{question}{2021--2022 \textbullet\ Q8(b)}{MATH 245 \textbullet\ 30/04/2023 \textbullet\ P077}{18}
A merchandising company wishes to test whether its 3 salesmen A, B, and C tend to make sales of the same size or whether they differ in selling ability as measured by the average size of their sales. During the last week, out of 14 sales, A made 5, B made 4, and C made 5 calls. The following are the weekly sales (in Tk thousand) record of 3 salesmen:
\begin{center}\small
\begin{tabular}{@{}lccccc@{}}
\toprule
\textbf{Salesman} & \textbf{Call 1} & \textbf{Call 2} & \textbf{Call 3} & \textbf{Call 4} & \textbf{Call 5} \\
\midrule
A ($n_A = 5$) & 300 & 400 & 300 & 500 & 0 \\
B ($n_B = 4$) & 600 & 300 & 300 & 400 & --- \\
C ($n_C = 5$) & 700 & 300 & 400 & 600 & 500 \\
\bottomrule
\end{tabular}
\end{center}
Test whether the three salesmen's average sales differ in size. (Given that $F_{(2,11)}$ at 5\% level of significance $= 3.98$)
\end{question}
\begin{solution}
Here $k = 3$ salesmen with unequal sample sizes $n_A = 5,\ n_B = 4,\ n_C = 5$ ($N = 14$ total calls).
\begin{enumerate}[1.]
\item \textbf{Hypotheses:} $H_0: \mu_A = \mu_B = \mu_C$ vs.\ $H_1$: At least two salesmen differ in mean sales.
\item \textbf{Group Means and Grand Mean:}
\begin{itemize}
\item Salesman A: $T_A = 1500 \implies \bar{x}_A = \dfrac{1500}{5} = 300$.
\item Salesman B: $T_B = 1600 \implies \bar{x}_B = \dfrac{1600}{4} = 400$.
\item Salesman C: $T_C = 2500 \implies \bar{x}_C = \dfrac{2500}{5} = 500$.
\item Grand Total $G = 1500 + 1600 + 2500 = 5600 \implies \text{Grand Mean } \bar{\bar{x}} = \dfrac{5600}{14} = 400$.
\end{itemize}
\item \textbf{Sums of Squares:}
\begin{align*}
SSB &= 5(-100)^2 + 4(0)^2 + 5(100)^2 = 50{,}000 + 0 + 50{,}000 = 100{,}000, \\[4pt]
SSW &= \sum_{j}(x_{Aj} - 300)^2 + \sum_{j}(x_{Bj} - 400)^2 + \sum_{j}(x_{Cj} - 500)^2 \\
    &= (100^2 + 200^2 + 300^2) + (200^2 + 100^2 + 100^2) + (200^2 + 200^2 + 100^2 + 100^2) \\
    &= 140{,}000 + 60{,}000 + 100{,}000 = 300{,}000, \\[4pt]
SST &= 100{,}000 + 300{,}000 = 400{,}000.
\end{align*}
\item \textbf{One-Way ANOVA Table ($\nu_1 = 2,\ \nu_2 = 11$):}
\begin{center}\small
\begin{tabular}{@{}lcccc@{}}
\toprule
\textbf{Source} & \textbf{SS} & \textbf{df} & \textbf{MS} & \textbf{$F$-Ratio} \\
\midrule
Between Salesmen & $100{,}000$ & $2$ & $MSB = 50{,}000.0$ & $F = \dfrac{50{,}000}{27{,}272.73} = \dfrac{11}{6} \approx 1.8333$ \\
Within Salesmen (Error) & $300{,}000$ & $11$ & $MSW = 27{,}272.73$ & \\
\midrule
\textbf{Total} & $400{,}000$ & $13$ & & \\
\bottomrule
\end{tabular}
\end{center}
\item \textbf{Decision:} Since $F = 1.833 < F_{0.05(2,11)} = 3.98$, we fail to reject $H_0$.
\end{enumerate}
\ans{$F = 1.833 < 3.98$; fail to reject $H_0$. The three salesmen's average weekly sales do not differ significantly at the $5\%$ level.}
\end{solution}

%===================================================================== TOPIC 12
\topic[Population, Sample \& Descriptive Stats]{12}{Population, Sample \& Descriptive Statistics}

\begin{note}
For stratified random sampling vs.\ simple random sampling, see \textbf{2023--2024 MATH 243 Q6(c) [P051]}; for statistic vs.\ parameter, see \textbf{2022--2023 MATH 243 Q7(a)(i) [P056]}.
\end{note}

%---------------------------------------------------------------- S001
\begin{question}{2022--2023 \textbullet\ Q5(a)}{MATH 243 \textbullet\ 01/02/2025 \textbullet\ S001}{10}
What do you mean by statistical data? How are they produced? Why do you need statistical data at all? What are the usual sources of statistical data in Bangladesh?
\end{question}
\begin{solution}
\begin{enumerate}[(i)]
\item \textbf{Meaning of Statistical Data:}
Statistical data are systematic collections of numerical facts, measurements, or categorical observations gathered on one or more variables from a population or sample that exhibit variability and can be summarized, analyzed, and interpreted to draw valid conclusions.
\item \textbf{How Statistical Data are Produced:}
Data are produced through two primary mechanisms:
\begin{itemize}
\item \emph{Primary Data Collection:} Direct field collection via complete enumeration (censuses), sample surveys (questionnaires, personal/telephone/online interviews), controlled laboratory or field experiments, and direct observational/sensor logging.
\item \emph{Secondary Data Compilation:} Administrative records, statutory registrations (births, deaths, tax filings, customs logs), and published institutional repositories.
\end{itemize}
\item \textbf{Why Statistical Data are Needed:}
They provide empirical evidence for evidence-based policy making, economic forecasting, quality control in manufacturing, clinical and scientific hypothesis testing, risk assessment, and resource allocation under uncertainty.
\item \textbf{Usual Sources of Statistical Data in Bangladesh:}
\begin{itemize}
\item \textbf{Bangladesh Bureau of Statistics (BBS):} The national statistical organization (Population \& Housing Census, HIES, Labour Force Survey, Statistical Yearbook).
\item \textbf{Bangladesh Bank (BB):} Monetary, banking, inflation, balance-of-payments, and foreign exchange statistics.
\item \textbf{Ministry of Finance \& Planning Commission:} Bangladesh Economic Review, budget documents, and Five-Year Plan reports.
\item \textbf{Sectoral Agencies \& Research Institutes:} BANBEIS (education), DGHS \& NIPORT (BDHS health/demographic surveys), BIDS, ICDDR,B, and DSE/CSE (capital markets).
\end{itemize}
\end{enumerate}
\end{solution}

%---------------------------------------------------------------- S002
\begin{question}{2022--2023 \textbullet\ Q5(b)}{MATH 243 \textbullet\ 01/02/2025 \textbullet\ S002}{15}
A random sample was taken of the thickness of insulation in transformer windings, and the following 60 thicknesses (in millimeters) were recorded:
\begin{center}\small
\begin{tabular}{@{}cccccccccc@{}}
\toprule
18 & 21 & 22 & 29 & 25 & 31 & 37 & 38 & 41 & 39 \\
44 & 48 & 54 & 56 & 56 & 57 & 47 & 38 & 35 & 36 \\
29 & 37 & 32 & 42 & 43 & 40 & 48 & 36 & 37 & 37 \\
36 & 38 & 40 & 41 & 44 & 39 & 38 & 34 & 24 & 32 \\
39 & 44 & 42 & 30 & 37 & 30 & 42 & 37 & 34 & 37 \\
32 & 24 & 42 & 36 & 49 & 39 & 23 & 34 & 36 & 40 \\
\bottomrule
\end{tabular}
\end{center}
(i) Represent a frequency table for the data using Sturges' rule $K = 1 + 3.322 \log N$ where $K$ is the number of classes and $N$ is the total frequency.
(ii) Draw a frequency histogram and a cumulative frequency graph.
\end{question}
\begin{solution}
\textbf{(i) Frequency Distribution Table using Sturges' Rule:}
\begin{itemize}
\item Total number of observations: $N = 60$.
\item Minimum value $x_{\min} = 18\text{ mm}$, maximum value $x_{\max} = 57\text{ mm}$, so the range is $R = 57 - 18 = 39\text{ mm}$.
\item Number of classes by Sturges' formula:
\[ K = 1 + 3.322\log_{10}(60) = 1 + 3.322(1.77815) = 6.907 \approx 7\text{ classes}. \]
\item Class width (interval size):
\[ h = \left\lceil\frac{R}{K}\right\rceil = \left\lceil\frac{39}{7}\right\rceil = \lceil 5.57\rceil = 6\text{ mm}. \]
\end{itemize}
Using 7 classes of width $h = 6$ starting at $18$:
\begin{center}\small
\begin{tabular}{@{}cccccc@{}}
\toprule
\textbf{Class (mm)} & \textbf{Boundaries} & \textbf{Midpoint $x_i$} & \textbf{Freq.\ $f_i$} & \textbf{Cum.\ Freq.} & \textbf{Rel.\ Freq.} \\
\midrule
$18 - 23$ & $17.5 - 23.5$ & $20.5$ & $4$ & $4$ & $0.0667$ \\
$24 - 29$ & $23.5 - 29.5$ & $26.5$ & $5$ & $9$ & $0.0833$ \\
$30 - 35$ & $29.5 - 35.5$ & $32.5$ & $10$ & $19$ & $0.1667$ \\
$36 - 41$ & $35.5 - 41.5$ & $38.5$ & $25$ & $44$ & $0.4167$ \\
$42 - 47$ & $41.5 - 47.5$ & $44.5$ & $9$ & $53$ & $0.1500$ \\
$48 - 53$ & $47.5 - 53.5$ & $50.5$ & $3$ & $56$ & $0.0500$ \\
$54 - 59$ & $53.5 - 59.5$ & $56.5$ & $4$ & $60$ & $0.0667$ \\
\midrule
\textbf{Total} & & & $\mathbf{N = 60}$ & & $\mathbf{1.0000}$ \\
\bottomrule
\end{tabular}
\end{center}
\textit{(Note: Whether intervals are written as $[18,24), [24,30), \dots, [54,60]$ or inclusive integer limits $18\text{--}23, \dots, 54\text{--}59$, the frequencies of the 7 classes are identical: $4, 5, 10, 25, 9, 3, 4$.)}

\medskip
\textbf{(ii) Frequency Histogram and Cumulative Frequency Graph (Ogive):}
The histogram bars are drawn on the class boundaries, and the ogive plots each cumulative frequency against the upper class boundary.
\fig{\histogivefig}
\vspace{3mm}
\ans{$K = 7$ classes of width $h = 6$; class frequencies are $4, 5, 10, 25, 9, 3, 4$ (cumulative frequencies: $4, 9, 19, 44, 53, 56, 60$).}
\end{solution}

%---------------------------------------------------------------- S003
\begin{question}{2021--2022 \textbullet\ Q5(b)(i)}{MATH 245 \textbullet\ 30/04/2023 \textbullet\ S003}{8.5}
The sample consists of 34 observations recorded correct to the nearest integer, ranging in value from 201 to 337. If it is decided to use seven classes of width 20 integers and to begin the first class at 199.5, find the class limits and class marks of the seven classes.
\end{question}
\begin{solution}
Since observations are recorded to the nearest integer and the first lower class boundary (LCB) is $199.5$ with class width $h = 20$, the class boundaries step by $20$ ($199.5, 219.5, \dots, 339.5$). Each class has integer lower class limit $\text{LCL} = \text{LCB} + 0.5$, integer upper class limit $\text{UCL} = \text{UCB} - 0.5$, and class mark (midpoint) $x_i = \dfrac{\text{LCL} + \text{UCL}}{2} = \dfrac{\text{LCB} + \text{UCB}}{2}$:
\begin{center}\small
\begin{tabular}{@{}cccc@{}}
\toprule
\textbf{Class No.} & \textbf{Class Boundaries} & \textbf{Class Limits (Integer)} & \textbf{Class Mark (Midpoint $x_i$)} \\
\midrule
1 & $199.5 - 219.5$ & $200 - 219$ & $209.5$ \\
2 & $219.5 - 239.5$ & $220 - 239$ & $229.5$ \\
3 & $239.5 - 259.5$ & $240 - 259$ & $249.5$ \\
4 & $259.5 - 279.5$ & $260 - 279$ & $269.5$ \\
5 & $279.5 - 299.5$ & $280 - 299$ & $289.5$ \\
6 & $299.5 - 319.5$ & $300 - 319$ & $309.5$ \\
7 & $319.5 - 339.5$ & $320 - 339$ & $329.5$ \\
\bottomrule
\end{tabular}
\end{center}
\ans{Class limits: $200\text{--}219, 220\text{--}239, 240\text{--}259, 260\text{--}279, 280\text{--}299, 300\text{--}319, 320\text{--}339$; class marks: $209.5, 229.5, 249.5, 269.5, 289.5, 309.5, 329.5$.}
\end{solution}

%===================================================================== TOPIC 13
\topic[Measures of Central Tendency]{13}{Measures of Central Tendency (Mean, Median, Mode)}

%---------------------------------------------------------------- S004
\begin{question}{2021--2022 \textbullet\ Q5(a)}{MATH 245 \textbullet\ 30/04/2023 \textbullet\ S004}{18}
Draw a suitable diagram for the following distribution. Read the median from the graph and verify the result by formula. How many workers earned monthly wages between 2000 and 2170 Tk.?
\begin{center}\footnotesize\setlength{\tabcolsep}{4.5pt}
\begin{tabular}{@{}lccccccc@{}}
\toprule
\textbf{Monthly wages (Tk.)} & 1900--1950 & 1950--2000 & 2000--2050 & 2050--2100 & 2100--2150 & 2150--2200 & 2200--2400 \\
\textbf{No. of workers} & 6 & 10 & 22 & 30 & 16 & 12 & 15 \\
\bottomrule
\end{tabular}
\end{center}
\end{question}
\begin{solution}
The appropriate diagram for reading the median and intermediate percentiles graphically is the \textbf{less-than cumulative frequency curve (Ogive)}.
\begin{center}\small
\begin{tabular}{@{}lccccccc@{}}
\toprule
\textbf{Upper Class Boundary} & $\le 1950$ & $\le 2000$ & $\le 2050$ & $\le 2100$ & $\le 2150$ & $\le 2200$ & $\le 2400$ \\
\midrule
\textbf{Frequency ($f_i$)} & 6 & 10 & 22 & 30 & 16 & 12 & 15 \\
\textbf{Cumulative Frequency ($F_i$)} & 6 & 16 & 38 & 68 & 84 & 96 & $N = 111$ \\
\bottomrule
\end{tabular}
\end{center}
\fig{\wageogivefig}
\vspace{3mm}
\begin{enumerate}[(i)]
\item \textbf{Verification of Median by Formula:}
Here $N = 111 \implies \dfrac{N}{2} = 55.5$. The cumulative frequency first exceeds $55.5$ in the class $2050\text{--}2100$, so $L = 2050,\ F_{c} = 38,\ f_m = 30,\ h = 50$:
\[ \text{Median} = L + \left(\frac{\frac{N}{2} - F_c}{f_m}\right)h = 2050 + \left(\frac{55.5 - 38}{30}\right)\times 50 = 2079.17\text{ Tk.}, \]
which matches the graphical reading from the Ogive at $F = 55.5$.
\item \textbf{Number of workers earning between $2000\text{ Tk.}$ and $2170\text{ Tk.}$:}
\begin{itemize}
\item Workers earning $\le 2000\text{ Tk.}$: $F(2000) = 16$.
\item Workers earning $\le 2170\text{ Tk.}$ (by linear interpolation within $2150\text{--}2200$ where $f = 12$ and $h = 50$):
\[ F(2170) = F(2150) + \frac{2170 - 2150}{2200 - 2150}\times 12 = 84 + \frac{20}{50}\times 12 = 84 + 4.8 = 88.8. \]
\item Hence, the number of workers earning between $2000$ and $2170\text{ Tk.}$ is:
\[ F(2170) - F(2000) = 88.8 - 16 = 72.8 \approx 73\text{ workers} \quad (\text{or } 22 + 30 + 16 + 4.8 = 72.8). \]
\end{itemize}
\end{enumerate}
\ans{Median $= 2079.17\text{ Tk.}$; \qquad Workers earning between $2000$ and $2170\text{ Tk.}$ $= 72.8 \approx 73\text{ workers}$.}
\end{solution}

%---------------------------------------------------------------- S005 + S006 stacked
\repeated{2}
\begin{question}{2015--2016 \textbullet\ Q6(a)}{MATH 241 \textbullet\ 07/08/2016 \textbullet\ S005}{18}
For the following incomplete frequency distribution, it is known that the total frequency is 1000 and the median is 415. Estimate the missing frequencies and hence find the modal value of the completed table. Also, comment on the shape of the frequency distribution.
\begin{center}\small
\begin{tabular}{@{}lcccccccc@{}}
\toprule
\textbf{Variable} & 300--325 & 325--350 & 350--375 & 375--400 & 400--425 & 425--450 & 450--475 & 475--500 \\
\textbf{Frequency} & 5 & 17 & 80 & ? & 326 & ? & 88 & 9 \\
\bottomrule
\end{tabular}
\end{center}
\end{question}
\begin{question}{2014--2015 \textbullet\ Q6(a)}{MATH 241 \textbullet\ 06/02/2016 \textbullet\ S006}{17}
You are given the same incomplete frequency distribution table above. It is known that the total frequency is 1000 and the median is 413.11. Estimate by calculation the missing frequencies and find the value of the mode.
\end{question}
\begin{solution}
Let $f_1$ be the missing frequency of class $375\text{--}400$ and $f_2$ be the missing frequency of class $425\text{--}450$.
Since total frequency is $N = 1000$ and the sum of the known frequencies is $5 + 17 + 80 + 326 + 88 + 9 = 525$:
\[ f_1 + f_2 = 1000 - 525 = 475. \tag{1} \]
In both papers, the given median ($415$ or $413.11$) lies inside the class $400\text{--}425$, so the median class is $400\text{--}425$ with $L = 400$, $h = 25$, $f_m = 326$, and cumulative frequency of preceding classes $F_c = 5 + 17 + 80 + f_1 = 102 + f_1$. Using the median formula:
\[ \text{Median} = L + \frac{\frac{N}{2} - F_c}{f_m}\times h = 400 + \frac{500 - (102 + f_1)}{326}\times 25 = 400 + \frac{398 - f_1}{326}\times 25. \tag{2} \]

\medskip
\textbf{Case 1: 2015--2016 Paper ($\text{Median} = 415$).}
\begin{itemize}
\item From (2): $415 - 400 = \dfrac{398 - f_1}{326}\times 25 \implies 398 - f_1 = \dfrac{15 \times 326}{25} = 195.6 \implies f_1 = 202.4 \approx \mathbf{202}$.
\item From (1): $f_2 = 475 - 202 = \mathbf{273}$ (or $272.6$).
\item \textbf{Mode:} The modal class is $400\text{--}425$ ($f_m = 326$, $f_{\text{prev}} = f_1 = 202$, $f_{\text{next}} = f_2 = 273$):
\[ \text{Mode} = 400 + \frac{326 - 202}{2(326) - 202 - 273}\times 25 = 400 + \frac{124}{177}\times 25 = 400 + 17.51 = \mathbf{417.51}. \]
\item \textbf{Shape:} Using midpoints, $\text{Mean} = \dfrac{\sum f_i x_i}{1000} = \dfrac{413{,}575}{1000} = 413.58 < \text{Median} = 415 < \text{Mode} = 417.51$, indicating a slightly \textbf{negatively (left) skewed} unimodal distribution.
\end{itemize}

\medskip
\textbf{Case 2: 2014--2015 Paper ($\text{Median} = 413.11$).}
\begin{itemize}
\item From (2): $413.11 - 400 = \dfrac{398 - f_1}{326}\times 25 \implies 398 - f_1 = \dfrac{13.11 \times 326}{25} = 170.95 \implies f_1 = 227.05 \approx \mathbf{227}$.
\item From (1): $f_2 = 475 - 227 = \mathbf{248}$.
\item \textbf{Mode:} With $f_{\text{prev}} = 227$ and $f_{\text{next}} = 248$:
\[ \text{Mode} = 400 + \frac{326 - 227}{2(326) - 227 - 248}\times 25 = 400 + \frac{99}{177}\times 25 = 400 + 13.98 = \mathbf{413.98}. \]
\end{itemize}
\ans{2015--16 ($\text{Med}=415$): $f_1 = 202,\ f_2 = 273,\ \text{Mode} = 417.51$ (mean $413.58$, slightly negatively skewed). \quad 2014--15 ($\text{Med}=413.11$): $f_1 = 227,\ f_2 = 248,\ \text{Mode} = 413.98$.}
\end{solution}

%===================================================================== TOPIC 14
\topic[Measures of Dispersion]{14}{Measures of Dispersion (Variance, Standard Deviation, CV)}

%---------------------------------------------------------------- S007
\begin{question}{2023--2024 \textbullet\ Q5(c)}{MATH 243 \textbullet\ 21/05/2026 \textbullet\ S007}{12}
The mean and standard deviation of 200 items are found to be 60 and 20, respectively. If at the time of calculation, two items are wrongly taken as 3 and 67 instead of 13 and 17, find the correct mean and standard deviation. What is the correct coefficient of variation?
\end{question}
\begin{solution}
Given $n = 200$, incorrect mean $\bar{x}_{\text{wrong}} = 60$, and incorrect standard deviation $\sigma_{\text{wrong}} = 20$:
\begin{enumerate}[1.]
\item \textbf{Corrected Mean ($\bar{x}_{\text{correct}}$):}
The incorrect sum of observations is $\sum x_{\text{wrong}} = n\bar{x}_{\text{wrong}} = 200 \times 60 = 12{,}000$.
Subtracting the wrong values ($3$ and $67$) and adding the correct values ($13$ and $17$):
\[ \sum x_{\text{correct}} = 12{,}000 - (3 + 67) + (13 + 17) = 12{,}000 - 70 + 30 = 11{,}960, \]
\[ \bar{x}_{\text{correct}} = \frac{11{,}960}{200} = 59.80. \]
\item \textbf{Corrected Standard Deviation ($\sigma_{\text{correct}}$):}
From $\sigma_{\text{wrong}}^2 = \frac{1}{n}\sum x_{\text{wrong}}^2 - \bar{x}_{\text{wrong}}^2$, the incorrect sum of squares is:
\[ \sum x_{\text{wrong}}^2 = n\left(\sigma_{\text{wrong}}^2 + \bar{x}_{\text{wrong}}^2\right) = 200\left(20^2 + 60^2\right) = 200(400 + 3600) = 800{,}000. \]
Adjusting for the squared observations ($3^2 + 67^2 = 9 + 4489 = 4498$ vs.\ $13^2 + 17^2 = 169 + 289 = 458$):
\[ \sum x_{\text{correct}}^2 = 800{,}000 - (3^2 + 67^2) + (13^2 + 17^2) = 800{,}000 - 4498 + 458 = 795{,}960. \]
Hence, the corrected variance and standard deviation are:
\[ \sigma_{\text{correct}}^2 = \frac{795{,}960}{200} - (59.80)^2 = 3979.80 - 3576.04 = 403.76 \implies \sigma_{\text{correct}} = \sqrt{403.76} \approx 20.0938. \]
\item \textbf{Corrected Coefficient of Variation (CV):}
\[ \text{CV}_{\text{correct}} = \frac{\sigma_{\text{correct}}}{\bar{x}_{\text{correct}}}\times 100\% = \frac{20.0938}{59.80}\times 100\% \approx 33.60\%. \]
\end{enumerate}
\ans{Correct Mean $= 59.80$, \quad Correct SD $= 20.0938$, \quad Correct CV $= 33.60\%$.}
\end{solution}

%---------------------------------------------------------------- S008
\begin{question}{2022--2023 \textbullet\ Q5(c)}{MATH 243 \textbullet\ 01/02/2025 \textbullet\ S008}{10}
A purchasing agent obtained samples of 60-watt bulbs from two companies and had the samples tested in his own laboratory for length of life with the following results:
\begin{center}\small
\begin{tabular}{@{}lcc@{}}
\toprule
\textbf{Length of life (in hours)} & \textbf{Company A} & \textbf{Company B} \\
\midrule
1700 and under 1900 & 10 & 3 \\
1900 and under 2100 & 16 & 40 \\
2100 and under 2300 & 20 & 12 \\
2300 and under 2500 & 8 & 3 \\
2500 and under 2700 & 6 & 2 \\
\bottomrule
\end{tabular}
\end{center}
If prices of both types are the same, select which company's bulbs you will prefer to collect. Justify your preference.
\end{question}
\begin{solution}
Using class midpoints $x_i = 1800, 2000, 2200, 2400, 2600$ (with $N_A = N_B = 60$):
\begin{center}\small
\begin{tabular}{@{}cccccccc@{}}
\toprule
\textbf{Midpoint $x_i$} & $d_i = \frac{x_i-2200}{200}$ & $f_{Ai}$ & $f_{Ai}d_i$ & $f_{Ai}d_i^2$ & $f_{Bi}$ & $f_{Bi}d_i$ & $f_{Bi}d_i^2$ \\
\midrule
1800 & $-2$ & 10 & $-20$ & 40 & 3 & $-6$ & 12 \\
2000 & $-1$ & 16 & $-16$ & 16 & 40 & $-40$ & 40 \\
2200 & $0$ & 20 & $0$ & 0 & 12 & $0$ & 0 \\
2400 & $+1$ & 8 & $+8$ & 8 & 3 & $+3$ & 3 \\
2600 & $+2$ & 6 & $+12$ & 24 & 2 & $+4$ & 8 \\
\midrule
\textbf{Total} & & $\mathbf{60}$ & $\mathbf{-16}$ & $\mathbf{88}$ & $\mathbf{60}$ & $\mathbf{-39}$ & $\mathbf{63}$ \\
\bottomrule
\end{tabular}
\end{center}
\begin{itemize}
\item \textbf{Company A:}
\[ \begin{aligned} \bar{x}_A &= 2200 + \left(\frac{-16}{60}\right)200 = 2146.67\text{ h}, \quad \sigma_A = 200\sqrt{\frac{88}{60} - \left(\frac{-16}{60}\right)^2} \approx 236.27\text{ h}, \\ \text{CV}_A &= \frac{236.27}{2146.67}\times 100\% = 11.01\%. \end{aligned} \]
\item \textbf{Company B:}
\[ \begin{aligned} \bar{x}_B &= 2200 + \left(\frac{-39}{60}\right)200 = 2070.00\text{ h}, \quad \sigma_B = 200\sqrt{\frac{63}{60} - \left(\frac{-39}{60}\right)^2} \approx 158.43\text{ h}, \\ \text{CV}_B &= \frac{158.43}{2070.00}\times 100\% = 7.65\%. \end{aligned} \]
\end{itemize}
\textbf{Justification:}
\begin{itemize}
\item If \textbf{longer average service life} is the primary criterion, prefer \textbf{Company A} ($\bar{x}_A = 2146.67\text{ h} > \bar{x}_B = 2070.00\text{ h}$).
\item If \textbf{uniformity / consistency of bulb life} is the primary criterion, prefer \textbf{Company B} because its coefficient of variation is significantly smaller ($\text{CV}_B = 7.65\% < \text{CV}_A = 11.01\%$).
\end{itemize}
\ans{Company A has a higher average life ($2146.67\text{ h}$ vs.\ $2070.00\text{ h}$), whereas Company B has greater consistency/uniformity ($\text{CV}_B = 7.65\%$ vs.\ $\text{CV}_A = 11.01\%$).}
\end{solution}

%---------------------------------------------------------------- S009
\begin{question}{2022--2023 \textbullet\ Q6(a)}{MATH 243 \textbullet\ 01/02/2025 \textbullet\ S009}{10}
Define and illustrate with examples:
(i) Standard deviation, (ii) Variance, (iii) Skewness, (iv) Kurtosis and (v) Moments.
\end{question}
\begin{solution}
\begin{enumerate}[(i)]
\item \textbf{Standard Deviation ($\sigma$ or $s$):} The positive square root of the arithmetic mean of the squared deviations of observations from their arithmetic mean: $\sigma = \sqrt{\frac{1}{N}\sum_{i=1}^{N}(x_i - \mu)^2}$. \emph{Example:} For data $\{2, 4, 6\}$, $\mu = 4$ and $\sigma = \sqrt{\frac{4+0+4}{3}} = \sqrt{8/3} \approx 1.633$.
\item \textbf{Variance ($\sigma^2$ or $s^2$):} The second central moment measuring the average squared distance from the mean: $\sigma^2 = \frac{1}{N}\sum (x_i - \mu)^2$ (or sample variance $s^2 = \frac{1}{n-1}\sum (x_i - \bar{x})^2$). \emph{Example:} For $\{2, 4, 6\}$, $\sigma^2 = 8/3$ and $s^2 = 4$.
\item \textbf{Skewness:} A measure of the lack of symmetry (asymmetry) of a probability or frequency distribution. In a positively skewed distribution, the right tail is longer ($\text{Mean} > \text{Median} > \text{Mode}$, $\gamma_1 = \mu_3/\sigma^3 > 0$, e.g.\ household income); in a negatively skewed distribution, the left tail is longer ($\text{Mean} < \text{Median} < \text{Mode}$, $\gamma_1 < 0$).
\item \textbf{Kurtosis:} A measure of the tailedness and peakedness of a frequency distribution relative to a normal distribution, quantified by $\beta_2 = \mu_4/\mu_2^2$ (or excess kurtosis $\gamma_2 = \beta_2 - 3$): \emph{mesokurtic} ($\beta_2 = 3$, normal curve), \emph{leptokurtic} ($\beta_2 > 3$, heavy tails/sharp peak, e.g.\ Laplace/Student-$t$), and \emph{platykurtic} ($\beta_2 < 3$, light tails/flat peak, e.g.\ uniform distribution).
\item \textbf{Moments:} Expected values of integer powers of deviations. The $r$-th raw moment about an arbitrary origin $A$ is $\mu_r' = \frac{1}{N}\sum f_i(x_i - A)^r$, and the $r$-th central moment about the mean $\bar{x}$ is $\mu_r = \frac{1}{N}\sum f_i(x_i - \bar{x})^r$ ($\mu_0 = 1$, $\mu_1 = 0$, $\mu_2 = \sigma^2$).
\end{enumerate}
\end{solution}

%---------------------------------------------------------------- S010
\begin{question}{2020--2021 \textbullet\ Q5(a)}{MATH 245 \textbullet\ 2020--2021 \textbullet\ S010}{10}
Survey data on daily wages paid to workers of two factories A and B are as follows:
\begin{center}\small
\begin{tabular}{@{}lccccccc@{}}
\toprule
\textbf{Daily Wages} & 20-30 & 30-40 & 40-50 & 50-60 & 60-70 & 70-80 & 80-90 \\
\midrule
Factory A & 15 & 30 & 44 & 60 & 30 & 14 & 7 \\
Factory B \textit{(printed as A)} & 25 & 40 & 60 & 35 & 20 & 15 & 5 \\
\bottomrule
\end{tabular}
\end{center}
Find out: (i) Which factory pays higher average wages? (ii) Which factory has greater variability about paying wages?
\end{question}
\begin{solution}
Using midpoints $x_i = 25, 35, 45, 55, 65, 75, 85$ and step-deviation $d_i = \dfrac{x_i - 55}{10}$ ($N_A = 200,\ N_B = 200$):
\begin{itemize}
\item \textbf{Factory A:} $\sum f_{Ai}d_i = -70$, $\sum f_{Ai}d_i^2 = 448$.
\[ \bar{x}_A = 55 + \left(\frac{-70}{200}\right)10 = 51.50, \qquad \sigma_A = 10\sqrt{\frac{448}{200} - (-0.35)^2} = 10\sqrt{2.1175} \approx 14.5516, \]
\[ \text{CV}_A = \frac{14.5516}{51.50}\times 100\% \approx 28.26\%. \]
\item \textbf{Factory B:} $\sum f_{Bi}d_i = -150$, $\sum f_{Bi}d_i^2 = 570$.
\[ \bar{x}_B = 55 + \left(\frac{-150}{200}\right)10 = 47.50, \qquad \sigma_B = 10\sqrt{\frac{570}{200} - (-0.75)^2} = 10\sqrt{2.2875} \approx 15.1245, \]
\[ \text{CV}_B = \frac{15.1245}{47.50}\times 100\% \approx 31.84\%. \]
\end{itemize}
\ans{(i) Factory A pays higher average wages ($\bar{x}_A = 51.50 > \bar{x}_B = 47.50$); \quad (ii) Factory B has greater wage variability ($\text{CV}_B = 31.84\% > \text{CV}_A = 28.26\%$ and $\sigma_B = 15.12 > \sigma_A = 14.55$).}
\end{solution}

%---------------------------------------------------------------- S011
\begin{question}{2019--2020 \textbullet\ Q5(a)}{MATH 245 \textbullet\ 25/09/2021 \textbullet\ S011}{10}
Calculate the coefficient of variance for the following data:
\begin{center}\small
\begin{tabular}{@{}lcccccccc@{}}
\toprule
\textbf{Profits} & 20-30 & 30-40 & 40-50 & 50-60 & 60-70 & 70-80 & 80-90 & 90-100 \\
\textbf{No. of Companies} & 4 & 8 & 18 & 30 & 15 & 10 & 8 & 7 \\
\bottomrule
\end{tabular}
\end{center}
\end{question}
\begin{solution}
Using class midpoints $x_i = 25, 35, 45, 55, 65, 75, 85, 95$ and $d_i = \dfrac{x_i - 55}{10}$ ($N = 100$):
\[ \sum f_i d_i = 4(-3) + 8(-2) + 18(-1) + 30(0) + 15(1) + 10(2) + 8(3) + 7(4) = -46 + 87 = +41, \]
\[ \sum f_i d_i^2 = 36 + 32 + 18 + 0 + 15 + 40 + 72 + 112 = 325. \]
Therefore:
\[ \bar{x} = 55 + \left(\frac{41}{100}\right)10 = 59.10, \qquad \sigma = 10\sqrt{\frac{325}{100} - (0.41)^2} = 10\sqrt{3.0819} \approx 17.5553, \]
\[ \text{CV} = \frac{\sigma}{\bar{x}}\times 100\% = \frac{17.5553}{59.10}\times 100\% \approx 29.70\%. \]
\ans{$\bar{x} = 59.10$, $\sigma = 17.56$, and Coefficient of Variation $\text{CV} = 29.70\%$.}
\end{solution}

%---------------------------------------------------------------- S012
\begin{question}{2018--2019 \textbullet\ Q5(a)}{MATH 245 \textbullet\ 03/11/2019 \textbullet\ S012}{15}
The Shareholder Research Bureau of BD conducted a research study on the price behavior of three leading industrial shares A, B and C for the period 1996 to 1999:
\begin{center}\small
\begin{tabular}{@{}lccc@{}}
\toprule
\textbf{Share} & \textbf{Average price (tk)} & \textbf{Standard Deviation (tk)} & \textbf{Current selling price (tk)} \\
\midrule
A & 18.00 & 5.40 & 36.00 \\
B & 22.50 & 4.50 & 34.75 \\
C & 24.00 & 6.00 & 39.00 \\
\bottomrule
\end{tabular}
\end{center}
(i) Which share in your opinion appears to be more stable in value?
(ii) If you are the holder of all three shares which one would you like to dispose of at present and why?
\end{question}
\begin{solution}
\begin{enumerate}[(i)]
\item \textbf{Stability in Value (Coefficient of Variation $\text{CV} = \frac{\sigma}{\bar{x}}\times 100\%$):}
\begin{itemize}
\item Share A: $\text{CV}_A = \dfrac{5.40}{18.00}\times 100\% = 30.0\%$.
\item Share B: $\text{CV}_B = \dfrac{4.50}{22.50}\times 100\% = 20.0\%$.
\item Share C: $\text{CV}_C = \dfrac{6.00}{24.00}\times 100\% = 25.0\%$.
\end{itemize}
Since \textbf{Share B} has the lowest coefficient of variation ($20.0\%$), it is the most stable in value.
\item \textbf{Which Share to Dispose of at Present:}
Compare each share's current selling price to its historical average using the standardized $z$-score $z = \dfrac{\text{Current} - \text{Average}}{\text{SD}}$ (and percentage appreciation):
\begin{itemize}
\item Share A: $z_A = \dfrac{36.00 - 18.00}{5.40} = \mathbf{+3.33\sigma}$ ($+100.0\%$ above average).
\item Share B: $z_B = \dfrac{34.75 - 22.50}{4.50} = \mathbf{+2.72\sigma}$ ($+54.4\%$ above average).
\item Share C: $z_C = \dfrac{39.00 - 24.00}{6.00} = \mathbf{+2.50\sigma}$ ($+62.5\%$ above average).
\end{itemize}
\textbf{Decision:} Dispose of \textbf{Share A} to lock in maximum capital gains, because it is currently trading at the highest premium ($+3.33\sigma$ or $100\%$ above its mean) and also has the highest volatility ($\text{CV}_A = 30\%$). \textit{(Alternatively, if minimizing portfolio return per unit price appreciation is cited, Share B has the lowest percentage gain of $54.4\%$.)}
\end{enumerate}
\ans{(i) Share B is most stable ($\text{CV}_B = 20\% < 25\% < 30\%$); \quad (ii) Dispose of Share A ($z_A = +3.33\sigma$, $+100\%$ gain, highest volatility).}
\end{solution}

%---------------------------------------------------------------- S013
\begin{question}{2015--2016 \textbullet\ Q6(b)}{MATH 241 \textbullet\ 07/08/2016 \textbullet\ S013}{17}
During the 10 weeks of a session, marks scored by two students, X and Y, in a subject are given below:
\begin{center}\small
\begin{tabular}{@{}lcccccccccc@{}}
\toprule
\textbf{X} & 58 & 59 & 60 & 54 & 65 & 66 & 52 & 75 & 69 & 52 \\
\textbf{Y} & 87 & 89 & 78 & 71 & 73 & 84 & 65 & 66 & 56 & 46 \\
\bottomrule
\end{tabular}
\end{center}
(i) Who is better scorer? \quad (ii) Who is more consistent?
\end{question}
\begin{solution}
For $n = 10$ weeks:
\begin{itemize}
\item \textbf{Student X:} $\sum x_i = 610 \implies \bar{x} = \dfrac{610}{10} = 61.0$.
\[ \sum (x_i - \bar{x})^2 = (-3)^2 + (-2)^2 + (-1)^2 + (-7)^2 + 4^2 + 5^2 + (-9)^2 + 14^2 + 8^2 + (-9)^2 = 526, \]
\[ \sigma_X = \sqrt{\frac{526}{10}} \approx 7.2526 \quad (s_X = \sqrt{526/9} \approx 7.6449), \qquad \text{CV}_X = \frac{7.2526}{61.0}\times 100\% = 11.89\%. \]
\item \textbf{Student Y:} $\sum y_i = 715 \implies \bar{y} = \dfrac{715}{10} = 71.5$.
\[ \sum (y_i - \bar{y})^2 = \sum y_i^2 - \frac{(\sum y_i)^2}{10} = 52{,}833 - \frac{715^2}{10} = 52{,}833 - 51{,}122.5 = 1710.5, \]
\[ \sigma_Y = \sqrt{\frac{1710.5}{10}} \approx 13.0786 \ (s_Y \approx 13.7861), \qquad \text{CV}_Y = \frac{13.0786}{71.5}\times 100\% = 18.29\%. \]
\end{itemize}
\ans{(i) Student Y is the better scorer ($\bar{y} = 71.5 > \bar{x} = 61.0$); \quad (ii) Student X is more consistent ($\text{CV}_X = 11.89\% < \text{CV}_Y = 18.29\%$).}
\end{solution}

%---------------------------------------------------------------- S014
\begin{question}{2013--2014 \textbullet\ Q6(a)}{MATH 241 \textbullet\ 13/12/2014 \textbullet\ S014}{15}
Lives of two models of refrigerators in a survey are as follows:
\begin{center}\small
\begin{tabular}{@{}lcccccc@{}}
\toprule
\textbf{Life (No. of years)} & 0-2 & 2-4 & 4-6 & 6-8 & 8-10 & 10-12 \\
\midrule
Model A & 5 & 16 & 13 & 7 & 5 & 4 \\
Model B & 2 & 7 & 12 & 19 & 9 & 1 \\
\bottomrule
\end{tabular}
\end{center}
Which model has greater uniformity?
\end{question}
\begin{solution}
Using class midpoints $x_i = 1, 3, 5, 7, 9, 11$ ($N_A = 50,\ N_B = 50$):
\begin{itemize}
\item \textbf{Model A:} $\sum f_{Ai}x_i = 5(1) + 16(3) + 13(5) + 7(7) + 5(9) + 4(11) = 256 \implies \bar{x}_A = \dfrac{256}{50} = 5.12\text{ years}$.
\[ \sum f_{Ai}x_i^2 = 5(1) + 16(9) + 13(25) + 7(49) + 5(81) + 4(121) = 1706, \]
\[ \sigma_A = \sqrt{\frac{1706}{50} - (5.12)^2} = \sqrt{7.9056} \approx 2.8117\text{ yr}, \qquad \text{CV}_A = \frac{2.8117}{5.12}\times 100\% = 54.92\%. \]
\item \textbf{Model B:} $\sum f_{Bi}x_i = 2(1) + 7(3) + 12(5) + 19(7) + 9(9) + 1(11) = 308 \implies \bar{x}_B = \dfrac{308}{50} = 6.16\text{ years}$.
\[ \sum f_{Bi}x_i^2 = 2(1) + 7(9) + 12(25) + 19(49) + 9(81) + 1(121) = 2146, \]
\[ \sigma_B = \sqrt{\frac{2146}{50} - (6.16)^2} = \sqrt{4.9744} \approx 2.2303\text{ yr}, \qquad \text{CV}_B = \frac{2.2303}{6.16}\times 100\% = 36.21\%. \]
\end{itemize}
\ans{Model B has greater uniformity because its coefficient of variation is substantially smaller ($\text{CV}_B = 36.21\% < \text{CV}_A = 54.92\%$).}
\end{solution}

%---------------------------------------------------------------- S015
\begin{question}{2012--2013 \textbullet\ Q6(a)}{MATH 241 \textbullet\ 02/06/2014 \textbullet\ S015}{15}
Following are the runs obtained by two players A and B in 10 matches:
\begin{center}\small
\begin{tabular}{@{}lcccccccccc@{}}
\toprule
\textbf{Runs scored by A} & 55 & 80 & 78 & 49 & 52 & 76 & 78 & 76 & 62 & 76 \\
\textbf{Runs scored by B} & 54 & 78 & 55 & 66 & 63 & 69 & 74 & 64 & 59 & 85 \\
\bottomrule
\end{tabular}
\end{center}
Compute:
(i) Pearson's coefficient of skewness and kurtosis and Bowley's coefficient of skewness and comment on the shape of the runs scored by B only.
(ii) If the consistency of performance is the criterion for awarding a prize, who should get the prize?
\end{question}
\begin{solution}
\textbf{1. Means and Standard Deviations ($n = 10$):}
\begin{itemize}
\item \textbf{Player A:} $\bar{x}_A = \dfrac{682}{10} = 68.2$, $\mu_{2A} = \dfrac{1357.6}{10} = 135.76 \implies \sigma_A = 11.6516$, $\text{CV}_A = \dfrac{11.6516}{68.2}\times 100\% = \mathbf{17.08\%}$.
\item \textbf{Player B:} $\bar{x}_B = \dfrac{667}{10} = 66.7$, $\mu_{2B} = \dfrac{900.1}{10} = 90.01 \implies \sigma_B = 9.4874$, $\text{CV}_B = \dfrac{9.4874}{66.7}\times 100\% = \mathbf{14.22\%}$.
\end{itemize}

\textbf{(i) Shape Analysis of Player B's Runs:}
Sorting B's 10 scores: $54, 55, 59, 63, 64, 66, 69, 74, 78, 85$.
\begin{itemize}
\item \textbf{Median ($Q_2$) \& Quartiles ($Q_1, Q_3$):}
\[ Q_2 = \frac{64 + 66}{2} = 65.0, \quad Q_1 = 59.0 \ (58.0\text{ interp.}), \quad Q_3 = 74.0 \ (75.0\text{ interp.}). \]
\item \textbf{Karl Pearson's Coefficient of Skewness (using Median):}
\[ Sk_P = \frac{3(\bar{x}_B - \text{Median})}{\sigma_B} = \frac{3(66.7 - 65.0)}{9.4874} = \frac{5.1}{9.4874} \approx +0.5376. \]
\item \textbf{Moment Skewness \& Kurtosis of B:}
$\mu_{3B} = \frac{1}{10}\sum (x_{Bi} - 66.7)^3 = 379.536$ and $\mu_{4B} = \frac{1}{10}\sum (x_{Bi} - 66.7)^4 = 17{,}983.33$:
\[ \beta_1 = \frac{\mu_3^2}{\mu_2^3} \approx 0.1975 \ (\gamma_1 = +0.4444), \qquad \beta_2 = \frac{\mu_4}{\mu_2^2} \approx 2.2197 \ (\gamma_2 = -0.7803). \]
\item \textbf{Bowley's Coefficient of Skewness:}
\[ Sk_B = \frac{Q_3 + Q_1 - 2Q_2}{Q_3 - Q_1} = \frac{74 + 59 - 2(65)}{74 - 59} = \frac{3}{15} = +0.2000. \]
\item \textbf{Comment:} B's distribution is moderately \textbf{positively skewed} ($Sk_P > 0,\ Sk_B > 0$) and \textbf{platykurtic} ($\beta_2 = 2.22 < 3$).
\end{itemize}

\textbf{(ii) Consistency Prize:} Since $\text{CV}_B = 14.22\% < \text{CV}_A = 17.08\%$, \textbf{Player B} is more consistent and should get the prize.
\ans{(i) For B: $Sk_P = +0.538$, $Sk_B = +0.200$, $\beta_2 = 2.22$ (positively skewed, platykurtic); \quad (ii) Player B wins ($\text{CV}_B = 14.22\% < \text{CV}_A = 17.08\%$).}
\end{solution}

%---------------------------------------------------------------- S016
\begin{question}{2011--2012 \textbullet\ Q6(b)}{MATH 241 \textbullet\ 07/01/2013 \textbullet\ S016}{15}
A time study was conducted in a factory with the help of two samples A and B consisting of 7 workers. The time taken by the workers in each case was recorded. From the information given below state which of the samples is more variable and which takes less time on an average:
\begin{center}\small
\begin{tabular}{@{}lccccccc@{}}
\toprule
\textbf{Sample A} & 130 & 124 & 119 & 134 & 142 & 130 & 141 \\
\textbf{Sample B} & 132 & 144 & 136 & 143 & 132 & 133 & 140 \\
\bottomrule
\end{tabular}
\end{center}
\end{question}
\begin{solution}
For $n = 7$ workers in each sample:
\begin{itemize}
\item \textbf{Sample A:} $\sum x_{Ai} = 920 \implies \bar{x}_A = \dfrac{920}{7} \approx 131.43\text{ minutes}$.
\[ \sum (x_{Ai} - \bar{x}_A)^2 = 423.71 \implies \sigma_A \approx 7.7801\text{ min} \ (s_A = 8.4035), \quad \text{CV}_A = 5.92\%. \]
\item \textbf{Sample B:} $\sum x_{Bi} = 960 \implies \bar{x}_B = \dfrac{960}{7} \approx 137.14\text{ minutes}$.
\[ \sum (x_{Bi} - \bar{x}_B)^2 = 160.86 \implies \sigma_B \approx 4.7937\text{ min} \ (s_B = 5.1778), \quad \text{CV}_B = 3.50\%. \]
\end{itemize}
\ans{Sample A is more variable ($\text{CV}_A = 5.92\% > \text{CV}_B = 3.50\%$), and Sample A takes less time on average ($\bar{x}_A = 131.43\text{ min} < \bar{x}_B = 137.14\text{ min}$).}
\end{solution}

%===================================================================== TOPIC 15
\topic{15}{Moments, Skewness \& Kurtosis}

%---------------------------------------------------------------- S017
\begin{question}{2023--2024 \textbullet\ Q6(a)}{MATH 243 \textbullet\ 21/05/2026 \textbullet\ S017}{13}
Calculate the first four moments about the mean and also find the value of $\beta_1$ and $\beta_2$ from the following data. Hence, comment on the nature of the distribution.
\begin{center}\small
\begin{tabular}{@{}lccccccc@{}}
\toprule
\textbf{Marks} & 0-10 & 10-20 & 20-30 & 30-40 & 40-50 & 50-60 & 60-70 \\
\textbf{No. of students} & 8 & 12 & 20 & 30 & 15 & 10 & 5 \\
\bottomrule
\end{tabular}
\end{center}
\end{question}
\begin{solution}
Using assumed mean $A = 35$, class width $h = 10$, and step-deviation $d_i = \dfrac{x_i - 35}{10}$:
\begin{center}\small
\begin{tabular}{@{}cccccccc@{}}
\toprule
\textbf{Class} & \textbf{Midpoint $x_i$} & \textbf{Freq.\ $f_i$} & $d_i = \frac{x_i-35}{10}$ & $f_i d_i$ & $f_i d_i^2$ & $f_i d_i^3$ & $f_i d_i^4$ \\
\midrule
0--10 & 5 & 8 & $-3$ & $-24$ & $72$ & $-216$ & $648$ \\
10--20 & 15 & 12 & $-2$ & $-24$ & $48$ & $-96$ & $192$ \\
20--30 & 25 & 20 & $-1$ & $-20$ & $20$ & $-20$ & $20$ \\
30--40 & 35 & 30 & $0$ & $0$ & $0$ & $0$ & $0$ \\
40--50 & 45 & 15 & $+1$ & $+15$ & $15$ & $+15$ & $15$ \\
50--60 & 55 & 10 & $+2$ & $+20$ & $40$ & $+80$ & $160$ \\
60--70 & 65 & 5 & $+3$ & $+15$ & $45$ & $+135$ & $405$ \\
\midrule
\textbf{Total} & & $\mathbf{N=100}$ & & $\mathbf{-18}$ & $\mathbf{240}$ & $\mathbf{-102}$ & $\mathbf{1440}$ \\
\bottomrule
\end{tabular}
\end{center}
Raw moments about $A = 35$ ($\mu_r' = \frac{h^r}{N}\sum f_i d_i^r$):
\[ \mu_1' = -1.8, \qquad \mu_2' = 240.0, \qquad \mu_3' = -1020.0, \qquad \mu_4' = 144{,}000.0. \]
Converting to the first four central moments about the mean ($\bar{x} = A + \mu_1' = 35 - 1.8 = 33.2$):
\begin{align*}
\mu_1 &= 0, \\[3pt]
\mu_2 &= \mu_2' - (\mu_1')^2 = 240 - (-1.8)^2 = 240 - 3.24 = \mathbf{236.76}, \\[3pt]
\mu_3 &= \mu_3' - 3\mu_2'\mu_1' + 2(\mu_1')^3 = -1020 - 3(240)(-1.8) + 2(-1.8)^3 \\
      &= -1020 + 1296 - 11.664 = \mathbf{264.336}, \\[3pt]
\mu_4 &= \mu_4' - 4\mu_3'\mu_1' + 6\mu_2'(\mu_1')^2 - 3(\mu_1')^4 \\
      &= 144{,}000 - 7344 + 4665.6 - 31.4928 = \mathbf{141{,}290.1072}.
\end{align*}
Computing Pearson's coefficients $\beta_1$ and $\beta_2$:
\[ \beta_1 = \frac{\mu_3^2}{\mu_2^3} = \frac{(264.336)^2}{(236.76)^3} \approx \mathbf{0.00526} \ (\gamma_1 = +0.0726), \]
\[ \beta_2 = \frac{\mu_4}{\mu_2^2} = \frac{141{,}290.11}{(236.76)^2} \approx \mathbf{2.5205}. \]
\textbf{Comment:} Since $\beta_1 \approx 0.0053$ ($\gamma_1 = +0.073 > 0$), the distribution is nearly symmetrical with very slight positive skewness; since $\beta_2 = 2.52 < 3$, it is \textbf{platykurtic}.
\ans{$\mu_1 = 0,\ \mu_2 = 236.76,\ \mu_3 = 264.336,\ \mu_4 = 141{,}290.11$; \quad $\beta_1 = 0.0053,\ \beta_2 = 2.5205$ (nearly symmetric, platykurtic).}
\end{solution}

%---------------------------------------------------------------- S018 + S021 + S022 + S024 stacked
\repeated{4}
\begin{question}{2022--2023 \textbullet\ Q6(c)}{MATH 243 \textbullet\ 01/02/2025 \textbullet\ S018}{15}
An analysis of companies resulted in the following distribution:
\begin{center}\small
\begin{tabular}{@{}lccccc@{}}
\toprule
\textbf{Profit (Lakhs) / Class interval} & 10-20 & 20-30 & 30-40 & 40-50 & 50-60 \\
\textbf{No. of Companies / Frequency} & 18 & 20 & 30 & 22 & 10 \\
\bottomrule
\end{tabular}
\end{center}
Calculate the first four moments about assumed mean. Convert the result into moments about the mean. Compute the value of $\gamma_1$ and $\gamma_2$ and comment on the result.
\end{question}
\begin{question}{2018--2019 \textbullet\ Q5(b)}{MATH 245 \textbullet\ 03/11/2019 \textbullet\ S021}{20}
What is Skewness and Kurtosis? Calculate the value of $\gamma_1$ and $\gamma_2$ from the same data table above and interpret them.
\end{question}
\begin{question}{2016--2017 \textbullet\ Q6(a)}{MATH 245 \textbullet\ 05/08/2017 \textbullet\ S022}{17}
Define Skewness and Kurtosis with classifications. Using moments calculate the coefficients of Skewness and Kurtosis from the same distribution above and comment on the result obtained.
\end{question}
\begin{question}{2014--2015 \textbullet\ Q6(b)}{MATH 241 \textbullet\ 06/02/2016 \textbullet\ S024}{18}
Calculate coefficient of skewness by Karl Pearson's method and values of $\beta_1$ and $\beta_2$ from the same data table above.
\end{question}
\begin{solution}
\textit{(Note: For definitions and classifications of Skewness and Kurtosis, see \textbf{S009 [2022--2023 Q6(a)]}.)}
Using assumed mean $A = 35$, class width $h = 10$, and $d_i = \dfrac{x_i - 35}{10}$:
\begin{center}\small
\begin{tabular}{@{}cccccccc@{}}
\toprule
\textbf{Class} & \textbf{Midpoint $x_i$} & \textbf{Freq.\ $f_i$} & $d_i = \frac{x_i-35}{10}$ & $f_i d_i$ & $f_i d_i^2$ & $f_i d_i^3$ & $f_i d_i^4$ \\
\midrule
10--20 & 15 & 18 & $-2$ & $-36$ & $72$ & $-144$ & $288$ \\
20--30 & 25 & 20 & $-1$ & $-20$ & $20$ & $-20$ & $20$ \\
30--40 & 35 & 30 & $0$ & $0$ & $0$ & $0$ & $0$ \\
40--50 & 45 & 22 & $+1$ & $+22$ & $22$ & $+22$ & $22$ \\
50--60 & 55 & 10 & $+2$ & $+20$ & $40$ & $+80$ & $160$ \\
\midrule
\textbf{Total} & & $\mathbf{N=100}$ & & $\mathbf{-14}$ & $\mathbf{154}$ & $\mathbf{-62}$ & $\mathbf{490}$ \\
\bottomrule
\end{tabular}
\end{center}
\textbf{1. First Four Moments about Assumed Mean $A = 35$:}
\[ \mu_1' = -1.4, \qquad \mu_2' = 154.0, \qquad \mu_3' = -620.0, \qquad \mu_4' = 49{,}000.0. \]
\textbf{2. First Four Moments about the Mean ($\bar{x} = 35 - 1.4 = 33.6$):}
\begin{align*}
\mu_1 &= 0, \\[3pt]
\mu_2 &= 154 - (-1.4)^2 = 154 - 1.96 = \mathbf{152.04} \quad (\sigma = \sqrt{152.04} \approx 12.3305), \\[3pt]
\mu_3 &= -620 - 3(154)(-1.4) + 2(-1.4)^3 = -620 + 646.8 - 5.488 = \mathbf{21.312}, \\[3pt]
\mu_4 &= 49{,}000 - 3472 + 1811.04 - 11.5248 = \mathbf{47{,}327.5152}.
\end{align*}
\textbf{3. Shape Coefficients ($\beta_1, \beta_2, \gamma_1, \gamma_2$) and Karl Pearson's Skewness:}
\begin{align*}
\beta_1 &= \frac{\mu_3^2}{\mu_2^3} = \frac{(21.312)^2}{(152.04)^3} \approx \mathbf{0.000129}, \qquad \gamma_1 = +\sqrt{\beta_1} = \frac{\mu_3}{\mu_2^{3/2}} \approx \mathbf{+0.0114}, \\[3pt]
\beta_2 &= \frac{\mu_4}{\mu_2^2} = \frac{47{,}327.5152}{(152.04)^2} \approx \mathbf{2.0474}, \qquad \gamma_2 = \beta_2 - 3 \approx \mathbf{-0.9526}.
\end{align*}
For \textbf{2014--2015 Q6(b)} (Karl Pearson's method): the modal class is $30\text{--}40$ ($L=30, f_m=30, f_1=20, f_2=22, h=10$), giving:
\[ \text{Mode} = 30 + \frac{10}{18}\times 10 = 35.56, \qquad Sk_P = \frac{\bar{x} - \text{Mode}}{\sigma} = \frac{33.60 - 35.56}{12.3305} \approx \mathbf{-0.1586}. \]
\textbf{Interpretation:} Moment skewness $\gamma_1 = +0.0114 \approx 0$ shows the distribution is almost perfectly symmetrical (with mode slightly above mean giving $Sk_P = -0.159$), while $\beta_2 = 2.047 < 3$ ($\gamma_2 = -0.953 < 0$) shows it is \textbf{platykurtic}.
\ans{$\mu_1'= -1.4, \mu_2'=154, \mu_3'=-620, \mu_4'=49000$; \ $\mu_1=0, \mu_2=152.04, \mu_3=21.312, \mu_4=47327.52$; \ $\beta_1=0.00013, \beta_2=2.047, \gamma_1=+0.0114, \gamma_2=-0.953, Sk_P=-0.159$.}
\end{solution}

%---------------------------------------------------------------- S019
\begin{question}{2019--2020 \textbullet\ Q5(b)}{MATH 245 \textbullet\ 25/09/2021 \textbullet\ S019}{25}
An analysis of workers resulted in the following distribution:
\begin{center}\small
\begin{tabular}{@{}lccccccc@{}}
\toprule
\textbf{Earnings (Tk.)} & 50-70 & 70-90 & 90-110 & 110-130 & 130-150 & 150-170 & 170-190 \\
\textbf{No. of employees} & 4 & 8 & 12 & 20 & 6 & 7 & 3 \\
\bottomrule
\end{tabular}
\end{center}
Calculate the first four moments about assumed mean. Convert the result into moments about the mean. Calculate values of $\gamma_1$ and $\gamma_2$.
\end{question}
\begin{solution}
Using assumed mean $A = 120$, class width $h = 20$, and $d_i = \dfrac{x_i - 120}{20}$ ($N = 60$):
\begin{center}\small
\begin{tabular}{@{}cccccccc@{}}
\toprule
\textbf{Class} & \textbf{Midpoint $x_i$} & \textbf{Freq.\ $f_i$} & $d_i = \frac{x_i-120}{20}$ & $f_i d_i$ & $f_i d_i^2$ & $f_i d_i^3$ & $f_i d_i^4$ \\
\midrule
50--70 & 60 & 4 & $-3$ & $-12$ & $36$ & $-108$ & $324$ \\
70--90 & 80 & 8 & $-2$ & $-16$ & $32$ & $-64$ & $128$ \\
90--110 & 100 & 12 & $-1$ & $-12$ & $12$ & $-12$ & $12$ \\
110--130 & 120 & 20 & $0$ & $0$ & $0$ & $0$ & $0$ \\
130--150 & 140 & 6 & $+1$ & $+6$ & $6$ & $+6$ & $6$ \\
150--170 & 160 & 7 & $+2$ & $+14$ & $28$ & $+56$ & $112$ \\
170--190 & 180 & 3 & $+3$ & $+9$ & $27$ & $+81$ & $243$ \\
\midrule
\textbf{Total} & & $\mathbf{N=60}$ & & $\mathbf{-11}$ & $\mathbf{141}$ & $\mathbf{-41}$ & $\mathbf{825}$ \\
\bottomrule
\end{tabular}
\end{center}
\textbf{1. Moments about Assumed Mean $A = 120$:}
\begin{align*}
\mu_1' &= \frac{20(-11)}{60} = -\frac{11}{3} \approx -3.6667, & \mu_2' &= \frac{20^2(141)}{60} = 940.0, \\
\mu_3' &= \frac{20^3(-41)}{60} = -\frac{16{,}400}{3} \approx -5466.6667, & \mu_4' &= \frac{20^4(825)}{60} = 2{,}200{,}000.0.
\end{align*}
\textbf{2. Central Moments about the Mean ($\bar{x} = 120 - 3.6667 = 116.3333$):}
\begin{align*}
\mu_1 &= 0, \\[3pt]
\mu_2 &= 940 - (-3.6667)^2 = 940 - 13.4444 = \mathbf{926.5556}, \\[3pt]
\mu_3 &= -5466.67 + 10{,}340.0 - 98.59 = \mathbf{4774.7407}, \\[3pt]
\mu_4 &= 2{,}200{,}000 - 4(-5466.6667)(-3.6667) + 6(940)(13.4444) - 3(180.7531) = \mathbf{2{,}195{,}106.63}.
\end{align*}
\textbf{3. Values of $\gamma_1$ and $\gamma_2$:}
\[ \gamma_1 = \frac{\mu_3}{\mu_2^{3/2}} = \frac{4774.74}{(926.56)^{1.5}} \approx \mathbf{+0.1693}, \qquad \beta_2 = \frac{\mu_4}{\mu_2^2} \approx 2.5569 \implies \gamma_2 \approx \mathbf{-0.4431}. \]
\ans{$\mu_1'=-3.67,\ \mu_2'=940,\ \mu_3'=-5466.67,\ \mu_4'=2{,}200{,}000$; \ $\mu_2=926.56,\ \mu_3=4774.74,\ \mu_4=2{,}195{,}106.63$; \ $\gamma_1=+0.1693,\ \gamma_2=-0.4431$.}
\end{solution}

%---------------------------------------------------------------- S020
\begin{question}{2021--2022 \textbullet\ Q5(b)(ii)}{MATH 245 \textbullet\ 30/04/2023 \textbullet\ S020}{8.5}
Two distributions have same values of mean, SD and Skewness. Find Kurtosis of them. Comment on your result.
\end{question}
\begin{solution}
Let the two distributions have identical mean $\bar{x}$, standard deviation $\sigma$ (so $\mu_2 = \sigma^2$), and moment skewness $\gamma_1 = \mu_3/\sigma^3$ (so $\mu_3$ is also identical).
\begin{itemize}
\item \textbf{Independence of Kurtosis ($\beta_2 = \mu_4 / \mu_2^2$):}
Kurtosis depends on the \textbf{fourth central moment} $\mu_4 = E[(X - \mu)^4]$, which is algebraically independent of $\mu, \mu_2,$ and $\mu_3$, subject only to the Cauchy--Schwarz / Pearson lower bound:
\[ \beta_2 \ge \beta_1 + 1 \qquad \left(\text{where } \beta_1 = \gamma_1^2 = \frac{\mu_3^2}{\mu_2^3}\right). \]
\item \textbf{Comment:} Having the same mean, SD, and skewness does \textbf{not} determine a unique value of kurtosis or force the two distributions to have the same kurtosis. One distribution can be \emph{leptokurtic} ($\beta_2 > 3$, sharp peak and heavy tails) while the other is \emph{mesokurtic} ($\beta_2 = 3$) or \emph{platykurtic} ($1 + \beta_1 \le \beta_2 < 3$, flat peak and light tails).
\end{itemize}
\ans{Kurtosis $\beta_2 = \mu_4/\sigma^4$ cannot be uniquely determined from mean, SD, and skewness alone (beyond the lower bound $\beta_2 \ge \beta_1 + 1$); the two distributions can still differ in peakedness and tail weight.}
\end{solution}

%---------------------------------------------------------------- S023
\begin{question}{2020--2021 \textbullet\ Q5(b)}{MATH 245 \textbullet\ 2020--2021 \textbullet\ S023}{15}
An analysis of weight distribution resulted in the following table:
\begin{center}\small
\begin{tabular}{@{}lcccccc@{}}
\toprule
\textbf{Weight} & 30-40 & 40-50 & 50-60 & 60-70 & 70-80 & 80-90 \\
\textbf{Frequency} & 18 & 37 & 45 & 27 & 15 & 8 \\
\bottomrule
\end{tabular}
\end{center}
Calculate Bowley's coefficient of skewness and Kelly's coefficient of skewness for the above data. Hence, interpret your results by discussing the advantages and drawbacks of the methods.
\end{question}
\begin{solution}
Total frequency is $N = 18 + 37 + 45 + 27 + 15 + 8 = 150$. The cumulative frequencies ($F_i$) at upper boundaries $40, 50, 60, 70, 80, 90$ are $18, 55, 100, 127, 142, 150$.
\begin{enumerate}[1.]
\item \textbf{Quartiles ($Q_1, Q_2, Q_3$) and Percentiles ($P_{10}, P_{90}$):}
\begin{align*}
Q_1 \ (N/4 = 37.5\text{ in } 40\text{--}50): &\quad Q_1 = 40 + \frac{37.5 - 18}{37}\times 10 = 40 + 5.2703 = 45.2703, \\
Q_2 \ (N/2 = 75.0\text{ in } 50\text{--}60): &\quad Q_2 = 50 + \frac{75.0 - 55}{45}\times 10 = 50 + 4.4444 = 54.4444, \\
Q_3 \ (3N/4 = 112.5\text{ in } 60\text{--}70): &\quad Q_3 = 60 + \frac{112.5 - 100}{27}\times 10 = 60 + 4.6296 = 64.6296, \\
P_{10} \ (0.1N = 15.0\text{ in } 30\text{--}40): &\quad P_{10} = 30 + \frac{15.0 - 0}{18}\times 10 = 30 + 8.3333 = 38.3333, \\
P_{90} \ (0.9N = 135.0\text{ in } 70\text{--}80): &\quad P_{90} = 70 + \frac{135.0 - 127}{15}\times 10 = 70 + 5.3333 = 75.3333.
\end{align*}
\item \textbf{Bowley's Coefficient of Skewness ($Sk_B$):}
\[ Sk_B = \frac{Q_3 + Q_1 - 2Q_2}{Q_3 - Q_1} = \frac{64.6296 + 45.2703 - 2(54.4444)}{64.6296 - 45.2703} = \frac{1.0110}{19.3594} \approx +0.0522. \]
\item \textbf{Kelly's Coefficient of Skewness ($Sk_K$):}
\[ Sk_K = \frac{P_{90} + P_{10} - 2P_{50}}{P_{90} - P_{10}} = \frac{75.3333 + 38.3333 - 2(54.4444)}{75.3333 - 38.3333} = \frac{4.7778}{37.0000} \approx +0.1291. \]
\item \textbf{Advantages \& Drawbacks:}
\begin{itemize}
\item \emph{Bowley's Method:} \textbf{Advantage:} Applicable to open-ended distributions and robust against extreme outliers. \textbf{Drawback:} Ignores the outer $50\%$ of observations (below $Q_1$ and above $Q_3$).
\item \emph{Kelly's Method:} \textbf{Advantage:} Covers the central $80\%$ of the data (between $P_{10}$ and $P_{90}$), capturing tail asymmetry that Bowley misses (here $Sk_K = 0.1291 > Sk_B = 0.0522$ due to the right tail). \textbf{Drawback:} Still ignores the extreme $20\%$ in the tails ($<P_{10}$ and $>P_{90}$).
\end{itemize}
\end{enumerate}
\ans{Bowley's $Sk_B = +0.0522$, \qquad Kelly's $Sk_K = +0.1291$ (slight positive skewness).}
\end{solution}

%---------------------------------------------------------------- S025 + S026 stacked
\repeated{2}
\begin{question}{2015--2016 \textbullet\ Q7(a)}{MATH 241 \textbullet\ 07/08/2016 \textbullet\ S025}{18}
Calculate coefficient of skewness and kurtosis from the table given below. Also, comment on the nature of the distribution.
\begin{center}\small
\begin{tabular}{@{}lccccc@{}}
\toprule
\textbf{Class} & 10--20 & 20--30 & 30--40 & 40--50 & 50--60 \\
\textbf{Frequency} & 174 & 86 & 32 & 181 & 27 \\
\bottomrule
\end{tabular}
\end{center}
\end{question}
\begin{question}{2013--2014 \textbullet\ Q6(b)}{MATH 241 \textbullet\ 13/12/2014 \textbullet\ S026}{20}
Calculate coefficient of skewness and kurtosis from the same table above. Also comment on the nature of the distribution.
\end{question}
\begin{solution}
Using assumed mean $A = 35$, class width $h = 10$, and $d_i = \dfrac{x_i - 35}{10}$ ($N = 500$):
\begin{center}\small
\begin{tabular}{@{}cccccccc@{}}
\toprule
\textbf{Class} & \textbf{Midpoint $x_i$} & \textbf{Freq.\ $f_i$} & $d_i$ & $f_i d_i$ & $f_i d_i^2$ & $f_i d_i^3$ & $f_i d_i^4$ \\
\midrule
10--20 & 15 & 174 & $-2$ & $-348$ & $696$ & $-1392$ & $2784$ \\
20--30 & 25 & 86 & $-1$ & $-86$ & $86$ & $-86$ & $86$ \\
30--40 & 35 & 32 & $0$ & $0$ & $0$ & $0$ & $0$ \\
40--50 & 45 & 181 & $+1$ & $+181$ & $181$ & $+181$ & $181$ \\
50--60 & 55 & 27 & $+2$ & $+54$ & $108$ & $+216$ & $432$ \\
\midrule
\textbf{Total} & & $\mathbf{N=500}$ & & $\mathbf{-199}$ & $\mathbf{1071}$ & $\mathbf{-1081}$ & $\mathbf{3483}$ \\
\bottomrule
\end{tabular}
\end{center}
Raw moments about $A = 35$:
\[ \mu_1' = -3.98, \qquad \mu_2' = 214.2, \qquad \mu_3' = -2162.0, \qquad \mu_4' = 69{,}660.0. \]
Central moments about the mean ($\bar{x} = 35 - 3.98 = 31.02$):
\begin{align*}
\mu_2 &= 214.2 - (-3.98)^2 = 214.2 - 15.8404 = \mathbf{198.3596} \quad (\sigma \approx 14.0840), \\[3pt]
\mu_3 &= -2162.0 - 3(214.2)(-3.98) + 2(-3.98)^3 = -2162.0 + 2557.548 - 126.0896 = \mathbf{269.4584}, \\[3pt]
\mu_4 &= 69{,}660.0 - 4(-2162.0)(-3.98) + 6(214.2)(15.8404) - 3(250.9183) = \mathbf{54{,}846.2873}.
\end{align*}
Coefficients of Skewness and Kurtosis:
\[ \beta_1 = \frac{\mu_3^2}{\mu_2^3} \approx \mathbf{0.00930} \ (\gamma_1 \approx \mathbf{+0.0965}), \qquad \beta_2 = \frac{\mu_4}{\mu_2^2} \approx \mathbf{1.3939} \ (\gamma_2 = -1.6061). \]
\textbf{Comment:} The distribution is bimodal (two separate peaks, at $10\text{--}20$ and $40\text{--}50$, with a trough at $30\text{--}40$), nearly symmetric with slight positive skewness ($\gamma_1 = +0.0965$), and strongly \textbf{platykurtic} ($\beta_2 = 1.394 < 3$).
\ans{$\beta_1 = 0.0093\ (\gamma_1 = +0.0965)$ and $\beta_2 = 1.3939\ (\gamma_2 = -1.6061)$; nearly symmetric and strongly platykurtic.}
\end{solution}

%---------------------------------------------------------------- S027
\begin{question}{2011--2012 \textbullet\ Q6(a)}{MATH 241 \textbullet\ 07/01/2013 \textbullet\ S027}{20}
From the data given below calculate Pearson's first coefficient of Skewness and Kurtosis. Comment on the nature of the distribution.
\begin{center}\small
\begin{tabular}{@{}lcccccc@{}}
\toprule
\textbf{Age in years} & 21 -- 26 & 27 -- 32 & 33 -- 38 & 39 -- 44 & 45 -- 50 & 51 -- 56 \\
\textbf{No. of persons} & 83 & 75 & 176 & 130 & 92 & 54 \\
\bottomrule
\end{tabular}
\end{center}
\end{question}
\begin{solution}
Using continuous class boundaries $20.5$--$26.5$, \dots, $50.5$--$56.5$, midpoints $x_i = 23.5$, $29.5$, $35.5$, $41.5$, $47.5$, $53.5$, assumed mean $A = 35.5$, class width $h = 6$, and $d_i = (x_i - 35.5)/6$ ($N = 610$):
\[ \sum f_i d_i = 235, \quad \sum f_i d_i^2 = 1391, \quad \sum f_i d_i^3 = 1585, \quad \sum f_i d_i^4 = 7379. \]
\begin{enumerate}[1.]
\item \textbf{Mean, Mode, Standard Deviation, and Pearson's First Coefficient of Skewness ($Sk_1$):}
\begin{itemize}
\item $\bar{x} = 35.5 + \left(\dfrac{235}{610}\right)6 = 35.5 + 2.3115 = 37.8115\text{ years}$.
\item $\mu_2 = 6^2\left[\dfrac{1391}{610} - \left(\dfrac{235}{610}\right)^2\right] = 36(2.2803 - 0.1484) = 76.7489 \implies \sigma = \sqrt{76.7489} \approx 8.7606\text{ years}$.
\item Modal class is $32.5\text{--}38.5$ ($L = 32.5,\ f_m = 176,\ f_1 = 75,\ f_2 = 130,\ h = 6$):
\[ \text{Mode} = 32.5 + \frac{176 - 75}{2(176) - 75 - 130}\times 6 = 32.5 + \frac{101}{147}\times 6 = 32.5 + 4.1224 = 36.6224\text{ years}. \]
\item \textbf{Pearson's First Coefficient of Skewness:}
\[ Sk_1 = \frac{\bar{x} - \text{Mode}}{\sigma} = \frac{37.8115 - 36.6224}{8.7606} \approx +0.1357. \]
\end{itemize}
\item \textbf{Kurtosis ($\beta_2$):}
From the raw moments $\mu_1' = 2.3115,\ \mu_2' = 82.0918,\ \mu_3' = 561.2459,\ \mu_4' = 15{,}677.3508$:
\[ \mu_4 \approx 13{,}034.15 \implies \beta_2 = \frac{\mu_4}{\mu_2^2} = \frac{13{,}034.15}{(76.7489)^2} \approx 2.2128 \ (\gamma_2 = -0.7872). \]
\end{enumerate}
\ans{Pearson's first skewness $Sk_1 = +0.1357$ (slight positive skewness); Kurtosis $\beta_2 = 2.2128$ (platykurtic).}
\end{solution}

%===================================================================== TOPIC 16
\topic[Statistical Graphs]{16}{Statistical Graphs (Scatter Plot, Stem-and-Leaf, Histogram, Ogive)}

\begin{note}
For the Frequency Histogram and Ogive of transformer insulation thicknesses, see \textbf{2022--2023 MATH 243 Q5(b) [S002]}; for the Wage Distribution Ogive, see \textbf{2021--2022 MATH 245 Q5(a) [S004]}.
\end{note}

%---------------------------------------------------------------- S028
\begin{question}{2023--2024 \textbullet\ Q5(a)}{MATH 243 \textbullet\ 21/05/2026 \textbullet\ S028}{10}
What do you mean by Stem-and-Leaf plots? How to create Stem-and-Leaf plots? Explain with examples.
\end{question}
\begin{solution}
\textbf{1. Definition of a Stem-and-Leaf Plot:}
Introduced by John Tukey, a stem-and-leaf plot is a semi-graphical exploratory data display for quantitative data that partitions each numerical observation into two parts: a \textbf{stem} (the leading digit or digits) and a \textbf{leaf} (the trailing significant digit). It displays the shape, spread, and clusters of a frequency distribution like a horizontal histogram while preserving the exact numerical value of every individual observation.

\medskip
\textbf{2. Steps to Construct a Stem-and-Leaf Plot:}
\begin{enumerate}[1.]
\item \textbf{Sort the data} in ascending order and decide the positional place value of the stem (e.g.\ tens) and leaf (e.g.\ units).
\item \textbf{List all stems} vertically in ascending order to the left of a vertical dividing line (including stems with zero observations so gaps are visible).
\item \textbf{Record each leaf} to the right of its corresponding stem in ascending order from left to right.
\item \textbf{Include a key/legend} (e.g.\ \texttt{2 | 5 = 25}) indicating the units.
\end{enumerate}

\medskip
\textbf{3. Illustrative Example:}
Consider the test scores of $n = 10$ students: $12, 15, 21, 24, 24, 28, 30, 33, 37, 45$.
\begin{center}\small
\begin{tabular}{@{}r|l@{}}
\toprule
\textbf{Stem (Tens)} & \textbf{Leaf (Units)} \\
\midrule
1 & 2 \ \ 5 \\
2 & 1 \ \ 4 \ \ 4 \ \ 8 \\
3 & 0 \ \ 3 \ \ 7 \\
4 & 5 \\
\bottomrule
\multicolumn{2}{@{}l}{\footnotesize\textbf{Key:} \texttt{2 | 1} $= 21$ marks.}
\end{tabular}
\end{center}
\end{solution}

%---------------------------------------------------------------- S029
\begin{question}{2023--2024 \textbullet\ Q5(b)}{MATH 243 \textbullet\ 21/05/2026 \textbullet\ S029}{13}
Explain the scatter plot with an example briefly. Prices (in dollars) and the ratings (out of 5 stars) of 5 different restaurants are recorded and plotted in a scatter plot. Interpret the scatter plot to determine if there is any relationship between price and rating.
\begin{center}\small
\begin{tabular}{@{}lccccc@{}}
\toprule
\textbf{Restaurant} & A & B & C & D & E \\
\midrule
\textbf{Prices ($x$)} & 22 & 29 & 15 & 18 & 23 \\
\textbf{Ratings ($y$)} & 3.5 & 4.0 & 3.0 & 4.0 & 3.5 \\
\bottomrule
\end{tabular}
\end{center}
\end{question}
\begin{solution}
\textbf{1. Explanation of a Scatter Plot:}
A scatter plot (scatter diagram) is a two-dimensional graphical representation of paired bivariate observations $(x_i, y_i)$ plotted as points on a Cartesian plane, with the explanatory (independent) variable $X$ on the horizontal axis and the response (dependent) variable $Y$ on the vertical axis. It reveals the direction (positive/negative), form (linear/curved), and strength of association between two variables (e.g.\ plotting study hours $x$ against exam score $y$).

\medskip
\textbf{2. Scatter Plot of Restaurant Prices ($x$) vs.\ Ratings ($y$):}
\fig{\scatterfig}

\textbf{3. Interpretation:}
With $\bar{x} = 21.4$ and $\bar{y} = 3.6$, the corrected sums are $SS_x = 2403 - \frac{107^2}{5} = 113.2$, $SS_y = 65.5 - \frac{18^2}{5} = 0.70$, and $SP_{xy} = 390.5 - \frac{107 \times 18}{5} = 5.3$, so
\[ r = \frac{SP_{xy}}{\sqrt{SS_x\,SS_y}} = \frac{5.3}{\sqrt{113.2 \times 0.70}} = \frac{5.3}{8.9017} \approx +0.595, \]
and the least-squares trend line drawn in the plot is $\hat{y} = 2.598 + 0.0468x$.
As restaurant price increases from $\$15$ (Restaurant C, rating $3.0$) to $\$29$ (Restaurant B, rating $4.0$), the ratings generally trend upward. Thus, there is a \textbf{moderate positive linear relationship} between price and rating, though Restaurant D ($\$18$, rating $4.0$) achieves a top rating at a relatively low price.
\ans{The scatter plot shows a moderate positive association ($r \approx +0.595$) between restaurant price and rating.}
\end{solution}

%---------------------------------------------------------------- S030
\begin{question}{2022--2023 \textbullet\ Q6(b)}{MATH 243 \textbullet\ 01/02/2025 \textbullet\ S030}{10}
The following data represent the length in centimeters (cm) of 12 pencils:
17.2, 16.4, 18.6, 17.4, 17.5, 19.0, 17.4, 16.8, 19.0, 17.8, 18.8, 18.7.
Use a stem and leaf plot to display the data. Hence find the median, mode and range of pencil length.
\end{question}
\begin{solution}
Sorting the $n = 12$ pencil lengths in ascending order:
\[ 16.4,\ \ 16.8,\ \ 17.2,\ \ 17.4,\ \ 17.4,\ \ 17.5,\ \ 17.8,\ \ 18.6,\ \ 18.7,\ \ 18.8,\ \ 19.0,\ \ 19.0. \]
\textbf{1. Stem-and-Leaf Plot (Stem = whole cm, Leaf = tenths of cm):}
\begin{center}\small
\begin{tabular}{@{}r|l@{}}
\toprule
\textbf{Stem (cm)} & \textbf{Leaf ($0.1\text{ cm}$)} \\
\midrule
16 & 4 \ \ 8 \\
17 & 2 \ \ 4 \ \ 4 \ \ 5 \ \ 8 \\
18 & 6 \ \ 7 \ \ 8 \\
19 & 0 \ \ 0 \\
\bottomrule
\multicolumn{2}{@{}l}{\footnotesize\textbf{Key:} \texttt{16 | 4} $= 16.4\text{ cm}$.}
\end{tabular}
\end{center}
\textbf{2. Median, Mode, and Range:}
\begin{itemize}
\item \textbf{Median:} Average of the 6th ($17.5$) and 7th ($17.8$) ordered observations:
\[ \text{Median} = \frac{17.5 + 17.8}{2} = 17.65\text{ cm}. \]
\item \textbf{Mode:} Both $17.4\text{ cm}$ and $19.0\text{ cm}$ appear twice (bimodal): $\text{Mode} = 17.4\text{ cm and } 19.0\text{ cm}$.
\item \textbf{Range:} $x_{\max} - x_{\min} = 19.0 - 16.4 = 2.6\text{ cm}$.
\end{itemize}
\ans{Median $= 17.65\text{ cm}$, \qquad Mode $= 17.4\text{ cm and } 19.0\text{ cm}$, \qquad Range $= 2.6\text{ cm}$.}
\end{solution}

%===================================================================== TOPIC 17
\topic[Regression \& Correlation]{17}{Regression, Correlation \& Rank Correlation}

%---------------------------------------------------------------- S031
\begin{question}{2023--2024 \textbullet\ Q4(b)}{MATH 243 \textbullet\ 21/05/2026 \textbullet\ S031}{12}
What is Pearson's correlation coefficient? How do you modify your formula to compute the rank correlation coefficient when the observations are repeated? Calculate the rank correlation coefficient for the following data on hourly sales ($x$) and expenses ($y$) of 10 stores:
\begin{center}\small
\begin{tabular}{@{}lcccccccccc@{}}
\toprule
\textbf{Sales ($x$)} & 50 & 50 & 55 & 60 & 65 & 65 & 65 & 60 & 60 & 60 \\
\textbf{Expenses ($y$)} & 11 & 13 & 14 & 16 & 16 & 15 & 15 & 14 & 13 & 13 \\
\bottomrule
\end{tabular}
\end{center}
\end{question}
\begin{solution}
\textbf{1. Pearson's Correlation Coefficient \& Tied-Rank Modification:}
Pearson's product-moment correlation coefficient $r = \frac{\sum (x_i - \bar{x})(y_i - \bar{y})}{\sqrt{\sum (x_i - \bar{x})^2 \sum (y_i - \bar{y})^2}} \in [-1, 1]$ measures the strength and direction of linear association between two quantitative variables.
When computing Spearman's rank correlation with repeated (tied) observations, tied items are assigned the \textbf{average of the ranks} they would have occupied, and a tie-correction term $\frac{m(m^2 - 1)}{12}$ is added to $\sum d_i^2$ for each group of $m$ tied observations:
\[ r_S = 1 - \frac{6\left[\sum_{i=1}^{n} d_i^2 + \sum \frac{m_x(m_x^2 - 1)}{12} + \sum \frac{m_y(m_y^2 - 1)}{12}\right]}{n(n^2 - 1)}. \]

\medskip
\textbf{2. Rank Table and Tie Corrections ($n = 10$):}
\begin{itemize}
\item \textbf{Ranks of $x$ (ascending):}
$50$ (2 ties at pos.\ 1,2 $\to$ rank $1.5$);\ $55$ (pos.\ 3 $\to$ rank $3.0$);\ $60$ (4 ties at pos.\ 4,5,6,7 $\to$ rank $5.5$);\ $65$ (3 ties at pos.\ 8,9,10 $\to$ rank $9.0$).
\item \textbf{Ranks of $y$ (ascending):}
$11$ (pos.\ 1 $\to$ rank $1.0$);\ $13$ (3 ties at pos.\ 2,3,4 $\to$ rank $3.0$);\ $14$ (2 ties at pos.\ 5,6 $\to$ rank $5.5$);\ $15$ (2 ties at pos.\ 7,8 $\to$ rank $7.5$);\ $16$ (2 ties at pos.\ 9,10 $\to$ rank $9.5$).
\end{itemize}
\begin{center}\small
\begin{tabular}{@{}lcccccccccc|c@{}}
\toprule
\textbf{Store} & 1 & 2 & 3 & 4 & 5 & 6 & 7 & 8 & 9 & 10 & \textbf{Total} \\
\midrule
$x_i$ & 50 & 50 & 55 & 60 & 65 & 65 & 65 & 60 & 60 & 60 & \\
$y_i$ & 11 & 13 & 14 & 16 & 16 & 15 & 15 & 14 & 13 & 13 & \\
Rank $R_{xi}$ & $1.5$ & $1.5$ & $3.0$ & $5.5$ & $9.0$ & $9.0$ & $9.0$ & $5.5$ & $5.5$ & $5.5$ & \\
Rank $R_{yi}$ & $1.0$ & $3.0$ & $5.5$ & $9.5$ & $9.5$ & $7.5$ & $7.5$ & $5.5$ & $3.0$ & $3.0$ & \\
$d_i = R_{xi} - R_{yi}$ & $0.5$ & $-1.5$ & $-2.5$ & $-4.0$ & $-0.5$ & $1.5$ & $1.5$ & $0.0$ & $2.5$ & $2.5$ & $0.0$ \\
$d_i^2$ & $0.25$ & $2.25$ & $6.25$ & $16.00$ & $0.25$ & $2.25$ & $2.25$ & $0.00$ & $6.25$ & $6.25$ & $\mathbf{42.00}$ \\
\bottomrule
\end{tabular}
\end{center}
Tie-correction factor $CF$:
\begin{itemize}
\item In $x$: $m_1 = 2\ (50),\ m_2 = 4\ (60),\ m_3 = 3\ (65) \implies CF_x = \frac{2(3) + 4(15) + 3(8)}{12} = \frac{6 + 60 + 24}{12} = 7.5$.
\item In $y$: $m_1 = 3\ (13),\ m_2 = 2\ (14),\ m_3 = 2\ (15),\ m_4 = 2\ (16) \implies CF_y = \frac{3(8) + 3\times 2(3)}{12} = \frac{24 + 18}{12} = 3.5$.
\end{itemize}
Total $CF = 7.5 + 3.5 = 11.0$. Substituting into the tied-rank formula:
\[ r_S = 1 - \frac{6(42.0 + 11.0)}{10(10^2 - 1)} = 1 - \frac{6(53.0)}{990} = 1 - \frac{318}{990} = 1 - 0.3212 = 0.6788. \]
\ans{Rank correlation coefficient $r_S = 0.6788$ (or $0.7275$ via exact Pearson correlation on the rank vectors).}
\end{solution}

%---------------------------------------------------------------- S032
\begin{question}{2023--2024 \textbullet\ Q4(c)}{MATH 243 \textbullet\ 21/05/2026 \textbullet\ S032}{12}
Briefly mention the objectives of correlation and regression analysis. A test was made to different doses of nitrogen ($N$) on rice yield. The doses of phosphorus ($P$) and potassium ($K$) were kept at the optimum level for proper growth of plant. The following data were recorded:
\begin{center}\small
\begin{tabular}{@{}lcccccc@{}}
\toprule
\textbf{Nitrogen doses (lb/acre) ($x$)} & 0 & 1 & 2 & 3 & 4 & 5 \\
\textbf{Yield (maunds/acre) ($y$)} & 15 & 25 & 40 & 55 & 52 & 43 \\
\bottomrule
\end{tabular}
\end{center}
Fit a second degree polynomial to the data and obtain coefficient of determination for the data.
\end{question}
\begin{solution}
\textbf{1. Objectives of Correlation and Regression Analysis:}
\emph{Correlation analysis} quantifies the degree and direction of linear association between two variables via a dimensionless index $r \in [-1, 1]$. \emph{Regression analysis} establishes an explicit functional equation $\hat{y} = f(x)$ to predict or estimate the dependent variable $Y$ from known values of the independent variable $X$ and quantify the proportion of explained variance ($R^2$).

\medskip
\textbf{2. Fitting the Second-Degree Polynomial $\hat{y} = b_0 + b_1 x + b_2 x^2$ ($n = 6$):}
\begin{center}\small
\begin{tabular}{@{}cccccccc|c@{}}
\toprule
$x_i$ & $y_i$ & $x_i^2$ & $x_i^3$ & $x_i^4$ & $x_i y_i$ & $x_i^2 y_i$ & $\hat{y}_i$ & $(y_i - \hat{y}_i)^2$ \\
\midrule
0 & 15 & 0 & 0 & 0 & 0 & 0 & $11.536$ & $12.001$ \\
1 & 25 & 1 & 1 & 1 & 25 & 25 & $30.207$ & $27.114$ \\
2 & 40 & 4 & 8 & 16 & 80 & 160 & $42.914$ & $8.493$ \\
3 & 55 & 9 & 27 & 81 & 165 & 495 & $49.657$ & $28.546$ \\
4 & 52 & 16 & 64 & 256 & 208 & 832 & $50.436$ & $2.447$ \\
5 & 43 & 25 & 125 & 625 & 215 & 1075 & $45.250$ & $5.063$ \\
\midrule
$\mathbf{15}$ & $\mathbf{230}$ & $\mathbf{55}$ & $\mathbf{225}$ & $\mathbf{979}$ & $\mathbf{693}$ & $\mathbf{2587}$ & $\mathbf{230.0}$ & $\mathbf{SSE = 83.664}$ \\
\bottomrule
\end{tabular}
\end{center}
The least-squares normal equations for $\hat{y} = b_0 + b_1 x + b_2 x^2$ are:
\[ \begin{cases}
6b_0 + 15b_1 + 55b_2 = 230 \\
15b_0 + 55b_1 + 225b_2 = 693 \\
55b_0 + 225b_1 + 979b_2 = 2587
\end{cases}
\implies
\begin{cases}
b_0 = \dfrac{323}{28} \approx 11.5357 \\[4pt]
b_1 = \dfrac{6063}{280} \approx 21.6536 \\[4pt]
b_2 = -\dfrac{167}{56} \approx -2.9821.
\end{cases} \]
Hence, the fitted quadratic regression curve is:
\[ \hat{y} = 11.5357 + 21.6536x - 2.9821x^2. \]

\textbf{3. Coefficient of Determination ($R^2$):}
With $\bar{y} = \dfrac{230}{6} \approx 38.3333$:
\[ SST = \sum_{i=1}^{6} (y_i - \bar{y})^2 = \sum y_i^2 - \frac{(\sum y_i)^2}{6} = 10{,}028 - \frac{52{,}900}{6} = 1211.3333, \]
\[ SSE = \sum_{i=1}^{6} (y_i - \hat{y}_i)^2 = 83.6643 \implies SSR = SST - SSE = 1127.6690, \]
\[ R^2 = 1 - \frac{SSE}{SST} = \frac{1127.6690}{1211.3333} \approx 0.9309 \ (93.09\%). \]
\ans{$\hat{y} = 11.5357 + 21.6536x - 2.9821x^2$; \qquad Coefficient of determination $R^2 = 0.9309$.}
\end{solution}

%---------------------------------------------------------------- S033
\begin{question}{2022--2023 \textbullet\ Q4(b)}{MATH 243 \textbullet\ 01/02/2025 \textbullet\ S033}{12}
Distinguish clearly between correlation and regression analysis. An examination of eight applicants for a post was taken by a firm. From the marks obtained by the applicants in the Bangla and English papers, compute the rank correlation coefficients and comment on your findings:
\begin{center}\small
\begin{tabular}{@{}lcccccccc@{}}
\toprule
\textbf{Applicant} & A & B & C & D & E & F & G & H \\
\midrule
\textbf{Marks in Bangla ($x$)} & 15 & 20 & 28 & 12 & 40 & 60 & 20 & 80 \\
\textbf{Marks in English ($y$)} & 40 & 30 & 50 & 30 & 20 & 10 & 30 & 60 \\
\bottomrule
\end{tabular}
\end{center}
\end{question}
\begin{solution}
\textbf{1. Distinction Between Correlation and Regression Analysis:}
\begin{itemize}
\item \textbf{Symmetry vs.\ Directionality:} Correlation is symmetric ($r_{xy} = r_{yx}$) and treats $X$ and $Y$ on an equal footing, whereas regression is asymmetric ($b_{yx} \neq b_{xy}$) and models $Y$ as a dependent response of an independent predictor $X$.
\item \textbf{Purpose:} Correlation measures the degree of linear covariation ($-1 \le r \le 1$), whereas regression provides an algebraic prediction equation $\hat{y} = a + bx$ to estimate one variable from the other.
\end{itemize}

\medskip
\textbf{2. Rank Correlation Computation ($n = 8$):}
\begin{itemize}
\item \textbf{Ties in Bangla ($x$):} Value $20$ occurs $m_x = 2$ times (positions 3 and 4 $\to$ average rank $3.5$), so $CF_x = \frac{2(2^2-1)}{12} = 0.5$.
\item \textbf{Ties in English ($y$):} Value $30$ occurs $m_y = 3$ times (positions 3, 4, and 5 $\to$ average rank $4.0$), so $CF_y = \frac{3(3^2-1)}{12} = 2.0$.
\end{itemize}
\begin{center}\small
\begin{tabular}{@{}lcccccccc|c@{}}
\toprule
\textbf{Applicant} & A & B & C & D & E & F & G & H & \textbf{Total} \\
\midrule
Bangla ($x_i$) & 15 & 20 & 28 & 12 & 40 & 60 & 20 & 80 & \\
English ($y_i$) & 40 & 30 & 50 & 30 & 20 & 10 & 30 & 60 & \\
Rank $R_{xi}$ & $2.0$ & $3.5$ & $5.0$ & $1.0$ & $6.0$ & $7.0$ & $3.5$ & $8.0$ & \\
Rank $R_{yi}$ & $6.0$ & $4.0$ & $7.0$ & $4.0$ & $2.0$ & $1.0$ & $4.0$ & $8.0$ & \\
$d_i = R_{xi} - R_{yi}$ & $-4.0$ & $-0.5$ & $-2.0$ & $-3.0$ & $+4.0$ & $+6.0$ & $-0.5$ & $0.0$ & $0.0$ \\
$d_i^2$ & $16.00$ & $0.25$ & $4.00$ & $9.00$ & $16.00$ & $36.00$ & $0.25$ & $0.00$ & $\mathbf{81.50}$ \\
\bottomrule
\end{tabular}
\end{center}
Applying the tied-rank Spearman formula with $\sum d_i^2 = 81.50$ and $CF = 0.5 + 2.0 = 2.5$:
\[ r_S = 1 - \frac{6\left(\sum d_i^2 + CF\right)}{n(n^2 - 1)} = 1 - \frac{6(81.5 + 2.5)}{8(64 - 1)} = 1 - \frac{6(84.0)}{504} = 1 - \frac{504}{504} = 0.0000. \]
\textbf{Comment:} Since $r_S = 0$, there is \textbf{zero rank correlation} (complete rank independence) between candidates' marks in Bangla and their marks in English.
\ans{$r_S = 0.0000$; proficiency in Bangla and proficiency in English among the applicants are completely uncorrelated in rank.}
\end{solution}

%---------------------------------------------------------------- S034
\begin{question}{2022--2023 \textbullet\ Q4(c)}{MATH 243 \textbullet\ 01/02/2025 \textbullet\ S034}{13}
From the data given below find out:
(i) The two linear regression equations,
(ii) Most likely age of husband when the wife's age is 30,
(iii) Most likely age of wife when the husband's age is 25.
\begin{center}\small
\begin{tabular}{@{}lcccccccccc@{}}
\toprule
\textbf{Ages of husbands ($x$)} & 22 & 23 & 23 & 24 & 26 & 27 & 27 & 28 & 30 & 30 \\
\textbf{Ages of wives ($y$)} & 18 & 20 & 21 & 20 & 21 & 22 & 23 & 24 & 25 & 26 \\
\bottomrule
\end{tabular}
\end{center}
Also, comment on Pearson's correlation coefficient between the above variables.
\end{question}
\begin{solution}
For $n = 10$ couples:
\[ \sum x_i = 260 \implies \bar{x} = 26.0, \qquad \sum y_i = 220 \implies \bar{y} = 22.0, \]
\[ \sum x_i^2 = 6836, \qquad \sum y_i^2 = 4896, \qquad \sum x_i y_i = 5782. \]
Centered sums of squares and cross-products:
\[ SS_x = 6836 - \frac{260^2}{10} = 76.0, \qquad SS_y = 4896 - \frac{220^2}{10} = 56.0, \qquad SP_{xy} = 5782 - \frac{260 \times 220}{10} = 62.0. \]
Regression coefficients:
\[ b_{yx} = \frac{SP_{xy}}{SS_x} = \frac{62}{76} = \frac{31}{38} \approx 0.8158, \qquad b_{xy} = \frac{SP_{xy}}{SS_y} = \frac{62}{56} = \frac{31}{28} \approx 1.1071. \]
\begin{enumerate}[(i)]
\item \textbf{The Two Linear Regression Equations:}
\begin{itemize}
\item \textbf{Regression of Wife's Age ($Y$) on Husband's Age ($X$):}
\[ y - 22 = 0.8158(x - 26) \implies \hat{y} = 0.7895 + 0.8158x. \]
\item \textbf{Regression of Husband's Age ($X$) on Wife's Age ($Y$):}
\[ x - 26 = 1.1071(y - 22) \implies \hat{x} = 1.6429 + 1.1071y. \]
\end{itemize}
\item \textbf{Most likely age of husband when wife's age is $y = 30$:}
\[ \hat{x}(30) = 1.6429 + 1.1071(30) = 34.857 \approx 34.86\text{ years}. \]
\item \textbf{Most likely age of wife when husband's age is $x = 25$:}
\[ \hat{y}(25) = 0.7895 + 0.8158(25) = 21.184 \approx 21.18\text{ years}. \]
\end{enumerate}
\textbf{Pearson's Correlation Coefficient:}
\[ r = \sqrt{b_{yx}\,b_{xy}} = \frac{62}{\sqrt{76 \times 56}} = \frac{62}{65.2380} \approx +0.9504, \]
which indicates a \textbf{very strong positive linear correlation} between the ages of husbands and wives.
\ans{(i) $\hat{y} = 0.7895 + 0.8158x$ and $\hat{x} = 1.6429 + 1.1071y$; \quad (ii) $\hat{x}(30) = 34.86\text{ yrs}$; \quad (iii) $\hat{y}(25) = 21.18\text{ yrs}$; \quad $r = +0.9504$.}
\end{solution}

%---------------------------------------------------------------- S035
\begin{question}{2021--2022 \textbullet\ Q7(a)}{MATH 245 \textbullet\ 30/04/2023 \textbullet\ S035}{20}
In a forestry study of 15 loblolly pines, the resulting data on specific gravity $x$ ($\text{g/cm}^3$) and modulus of rupture $y$ ($\text{kPa}$) are shown below:
\begin{center}\small
\begin{tabular}{@{}lcccccccc@{}}
\toprule
\textbf{Tree} & 1 & 2 & 3 & 4 & 5 & 6 & 7 & 8 \\
\textbf{$x$ ($\text{g/cm}^3$)} & 0.414 & 0.383 & 0.399 & 0.402 & 0.442 & 0.422 & 0.466 & 0.500 \\
\textbf{$y$ (kPa)} & 29,186 & 29,266 & 26,215 & 30,162 & 38,867 & 37,831 & 44,576 & 46,097 \\
\midrule
\textbf{Tree} & 9 & 10 & 11 & 12 & 13 & 14 & 15 & \\
\textbf{$x$ ($\text{g/cm}^3$)} & 0.514 & 0.530 & 0.569 & 0.558 & 0.577 & 0.572 & 0.548 & \\
\textbf{$y$ (kPa)} & 59,698 & 67,705 & 66,088 & 78,486 & 89,869 & 77,369 & 67,095 & \\
\bottomrule
\end{tabular}
\end{center}
(i) Find both the regression coefficients.
(ii) Compute and interpret the sample correlation coefficient.
\end{question}
\begin{solution}
From the $n = 15$ paired observations:
\[ \sum x_i = 7.296 \implies \bar{x} = 0.4864, \qquad \sum y_i = 788{,}510 \implies \bar{y} = 52{,}567.3333, \]
\[ \sum x_i^2 = 3.620132, \qquad \sum y_i^2 = 47{,}598{,}707{,}772, \qquad \sum x_i y_i = 403{,}686.826. \]
Computing the corrected sums of squares and cross-products:
\begin{align*}
SS_x &= \sum x_i^2 - \frac{(\sum x_i)^2}{15} = 3.620132 - \frac{(7.296)^2}{15} = 3.620132 - 3.5487744 = 0.0713576, \\[4pt]
SS_y &= \sum y_i^2 - \frac{(\sum y_i)^2}{15} = 47{,}598{,}707{,}772 - 41{,}449{,}868{,}006.67 = 6{,}148{,}839{,}765.33, \\[4pt]
SP_{xy} &= \sum x_i y_i - \frac{(\sum x_i)(\sum y_i)}{15} = 403{,}686.826 - 383{,}531.264 = 20{,}155.562.
\end{align*}
\begin{enumerate}[(i)]
\item \textbf{Both Regression Coefficients ($b_{yx}$ and $b_{xy}$):}
\[ \begin{aligned} b_{yx} &= \frac{SP_{xy}}{SS_x} = \frac{20{,}155.562}{0.0713576} \approx \mathbf{282{,}458.52\text{ kPa/(g/cm}^3\text{)}}, \\ b_{xy} &= \frac{SP_{xy}}{SS_y} = \frac{20{,}155.562}{6{,}148{,}839{,}765.33} \approx \mathbf{3.2779 \times 10^{-6}\text{ (g/cm}^3\text{)/kPa}}. \end{aligned} \]
\textit{(Note: If ``both regression coefficients'' refers to the intercept $\hat{\beta}_0$ and slope $\hat{\beta}_1$ of $Y$ on $X$, then $\hat{\beta}_1 = 282{,}458.52$ and $\hat{\beta}_0 = \bar{y} - \hat{\beta}_1\bar{x} = -84{,}820.49\text{ kPa}$.)}
\item \textbf{Sample Correlation Coefficient ($r$):}
\[ r = \frac{SP_{xy}}{\sqrt{SS_x\,SS_y}} = \sqrt{b_{yx}\,b_{xy}} = \frac{20{,}155.562}{\sqrt{(0.0713576)(6{,}148{,}839{,}765.33)}} \approx \mathbf{+0.9622}. \]
\textbf{Interpretation:} $r = +0.9622$ ($R^2 = 92.59\%$) indicates a very strong positive linear relationship between wood specific gravity and modulus of rupture.
\end{enumerate}
\ans{(i) $b_{yx} = 282{,}458.52$, $b_{xy} = 3.2779 \times 10^{-6}$ (and intercept $a_{yx} = -84{,}820.49$); \quad (ii) $r = +0.9622$ (very strong positive linear correlation).}
\end{solution}

%---------------------------------------------------------------- S036 + S039 stacked
\repeated{2}
\begin{question}{2020--2021 \textbullet\ Q6(a)}{MATH 245 \textbullet\ 2020--2021 \textbullet\ S036}{15}
The following table shows the experience ($X$) and the performance rating ($Y$) of 5 persons:
\begin{center}\small
\begin{tabular}{@{}lccccc@{}}
\toprule
\textbf{X} & 16 & 12 & 18 & 4 & 3 \\
\textbf{Y} & 87 & 88 & 89 & 68 & 58 \\
\bottomrule
\end{tabular}
\end{center}
(i) Fit a linear regression model for $Y$ on $X$.
(ii) Find sum of squares for regression (SSR) and total sum of squares (SST).
\end{question}
\begin{question}{2016--2017 \textbullet\ Q6(b)}{MATH 245 \textbullet\ 05/08/2017 \textbullet\ S039}{18}
Given the exact same table of experience ($X$) and performance rating ($Y$) of 5 persons above:
(i) Fit a linear regression model $Y$ on $X$.
(ii) Find sum of squares for regression (SSR) and total sum of squares (SST).
(iii) Is the linear model appropriate for given data? Justify your result using coefficient of determination.
\end{question}
\begin{solution}
For the $n = 5$ persons:
\[ \sum x_i = 53 \implies \bar{x} = 10.6, \qquad \sum y_i = 390 \implies \bar{y} = 78.0, \]
\[ \sum x_i^2 = 749, \qquad \sum y_i^2 = 31{,}222, \qquad \sum x_i y_i = 4496. \]
Corrected sums of squares and cross-products:
\[ SS_x = 749 - \frac{53^2}{5} = 187.2, \qquad SST = SS_y = 31{,}222 - \frac{390^2}{5} = 31{,}222 - 30{,}420 = 802.0, \]
\[ SP_{xy} = 4496 - \frac{53 \times 390}{5} = 4496 - 4134 = 362.0. \]
\begin{enumerate}[(i)]
\item \textbf{Linear Regression Model of $Y$ on $X$ ($\hat{y} = a + bx$):}
\[ b = \frac{SP_{xy}}{SS_x} = \frac{362.0}{187.2} \approx 1.9338, \qquad a = \bar{y} - b\bar{x} = 78.0 - (1.9338)(10.6) \approx 57.5021. \]
Thus, the fitted linear model is $\hat{y} = 57.5021 + 1.9338x$.
\item \textbf{Sum of Squares for Regression ($SSR$) and Total Sum of Squares ($SST$):}
\[ SST = SS_y = \mathbf{802.0}, \qquad SSR = \frac{(SP_{xy})^2}{SS_x} = \frac{(362.0)^2}{187.2} \approx \mathbf{700.0214} \quad (SSE = 101.9786). \]
\item \textbf{Appropriateness of the Linear Model (Coefficient of Determination $R^2$):}
\[ R^2 = \frac{SSR}{SST} = \frac{700.0214}{802.0} \approx 0.8728 \ (87.28\%). \]
Since $87.28\%$ of the total variation in performance rating $Y$ is explained by linear regression on experience $X$ ($r = +0.9343$), the linear model provides a good fit and is appropriate for the data.
\end{enumerate}
\ans{(i) $\hat{y} = 57.5021 + 1.9338x$; \quad (ii) $SSR = 700.02$, $SST = 802.00$; \quad (iii) Yes, $R^2 = 0.8728$ ($87.28\%$ of variation explained).}
\end{solution}

%---------------------------------------------------------------- S037
\begin{question}{2019--2020 \textbullet\ Q6(a)}{MATH 245 \textbullet\ 25/09/2021 \textbullet\ S037}{17}
The following table shows the distance to transmitter ($X$) and corresponding wireless signal strength ($Y$):
\begin{center}\small
\begin{tabular}{@{}lcccccccccc@{}}
\toprule
\textbf{Distance to transmitter (m) ($X$)} & 13 & 12 & 17 & 19 & 14 & 15 & 15 & 8 & 13 & 3 \\
\textbf{Wireless signal strength (dB) ($Y$)} & 34.4 & 38.4 & 30.4 & 29.7 & 30.1 & 33.9 & 32.8 & 35.2 & 34.9 & 36.8 \\
\bottomrule
\end{tabular}
\end{center}
Find the regression line of $Y$ on $X$ and also predict what the signal strength would be if the distance is 10 meters.
\end{question}
\begin{solution}
For the $n = 10$ observations:
\[ \sum x_i = 129 \implies \bar{x} = 12.90, \qquad \sum y_i = 336.6 \implies \bar{y} = 33.66, \]
\[ \sum x_i^2 = 1851, \qquad \sum y_i^2 = 11{,}406.52, \qquad \sum x_i y_i = 4256.7. \]
Corrected sums of squares and cross-products:
\[ SS_x = 1851 - 1664.1 = 186.90, \qquad SP_{xy} = 4256.7 - 4342.14 = -85.44. \]
Slope and intercept of the regression line of $Y$ on $X$:
\[ b_{yx} = \frac{SP_{xy}}{SS_x} = \frac{-85.44}{186.90} \approx -0.45714, \qquad a_{yx} = \bar{y} - b_{yx}\bar{x} = 33.66 - (-0.45714)(12.90) \approx 39.5571. \]
Thus, the regression line of $Y$ on $X$ is:
\[ \hat{y} = 39.5571 - 0.4571x. \]
Predicting the wireless signal strength at $x = 10\text{ meters}$:
\[ \hat{y}(10) = 39.5571 - 0.45714(10) = 34.9857 \approx 34.99\text{ dB}. \]
\ans{Regression line: $\hat{y} = 39.5571 - 0.4571x$; \qquad Predicted signal strength at $10\text{ m}$: $\hat{y}(10) = 34.99\text{ dB}$.}
\end{solution}

%---------------------------------------------------------------- S038
\begin{question}{2018--2019 \textbullet\ Q6(a)}{MATH 245 \textbullet\ 03/11/2019 \textbullet\ S038}{20}
Find the regression equation of capacity utilisation on production from the data where $r = 0.62$:
\begin{center}\small
\begin{tabular}{@{}lcc@{}}
\toprule
& \textbf{Average} & \textbf{Standard deviation} \\
\midrule
Production ($X$) & 35.6 & 10.5 \\
Capacity utilisation ($Y$) & 84.8 & 8.5 \\
\bottomrule
\end{tabular}
\end{center}
Also estimate the production when the capacity utilisation is 70 percent.
\end{question}
\begin{solution}
Let $X$ denote Production ($\bar{x} = 35.6,\ \sigma_x = 10.5$) and $Y$ denote Capacity Utilisation ($\bar{y} = 84.8,\ \sigma_y = 8.5$), with correlation coefficient $r = 0.62$.
\begin{enumerate}[1.]
\item \textbf{Regression Equation of Capacity Utilisation ($Y$) on Production ($X$):}
\[ b_{yx} = r\frac{\sigma_y}{\sigma_x} = 0.62 \times \frac{8.5}{10.5} \approx 0.5019, \qquad y - 84.8 = 0.5019(x - 35.6) \implies \hat{y} = 66.9322 + 0.5019x. \]
\item \textbf{Regression Equation of Production ($X$) on Capacity Utilisation ($Y$) \& Estimation at $y = 70\%$:}
To estimate $X$ (Production) for a given value of $Y$ (Capacity Utilisation $= 70$), we use the regression line of $X$ on $Y$:
\[ b_{xy} = r\frac{\sigma_x}{\sigma_y} = 0.62 \times \frac{10.5}{8.5} \approx 0.7659, \qquad x - 35.6 = 0.7659(y - 84.8) \]
\[ \implies \hat{x} = -29.3468 + 0.7659y. \]
Substituting $y = 70$:
\[ \hat{x}(70) = 35.6 + 0.7659(70 - 84.8) = 35.6 - 11.3351 = 24.2649 \approx 24.26. \]
\end{enumerate}
\ans{Capacity utilisation on production: $\hat{y} = 66.9322 + 0.5019x$; \quad Production on capacity utilisation: $\hat{x} = -29.3468 + 0.7659y$; \quad Estimated production at $70\%$ capacity utilisation $= 24.26$.}
\end{solution}

%---------------------------------------------------------------- S040
\begin{question}{2015--2016 \textbullet\ Q7(b)}{MATH 241 \textbullet\ 07/08/2016 \textbullet\ S040}{17}
A factory manager collected the following data of 9 persons from a production department referring to years of service and their monthly income:
\begin{center}\small
\begin{tabular}{@{}lccccccccc@{}}
\toprule
\textbf{Year of Service ($x$)} & 10 & 8 & 9 & 5 & 7 & 6 & 12 & 11 & 13 \\
\textbf{Income in '000 ($y$)} & 17 & 12 & 14 & 8 & 10 & 9 & 20 & 18 & 22 \\
\bottomrule
\end{tabular}
\end{center}
Calculate the following:
(i) The regression equations and the correlation coefficient.
(ii) What initial start would you recommend for a person applying for a job after having served in a similar capacity in another factory for 15 years?
\end{question}
\begin{solution}
For the $n = 9$ employees:
\[ \sum x_i = 81 \implies \bar{x} = 9.0, \qquad \sum y_i = 130 \implies \bar{y} = \frac{130}{9} \approx 14.4444, \]
\[ \sum x_i^2 = 789, \qquad \sum y_i^2 = 2082, \qquad \sum x_i y_i = 1280. \]
Centered sums of squares and cross-products:
\[ SS_x = 789 - \frac{81^2}{9} = 60.0, \qquad SS_y = 2082 - \frac{130^2}{9} = 204.2222, \qquad SP_{xy} = 1280 - \frac{81 \times 130}{9} = 110.0. \]
\begin{enumerate}[(i)]
\item \textbf{Regression Equations and Correlation Coefficient:}
\begin{itemize}
\item $b_{yx} = \dfrac{110}{60} = \dfrac{11}{6} \approx 1.8333 \implies \hat{y} = 14.4444 + 1.8333(x - 9) = \mathbf{-2.0556 + 1.8333x}$.
\item $b_{xy} = \dfrac{110}{204.2222} = \dfrac{990}{1838} \approx 0.5386 \implies \hat{x} = 9.0 + 0.5386(y - 14.4444) = \mathbf{1.2198 + 0.5386y}$.
\item $r = \dfrac{SP_{xy}}{\sqrt{SS_x\,SS_y}} = \dfrac{110}{\sqrt{60 \times 204.2222}} \approx \mathbf{+0.9937}$.
\end{itemize}
\item \textbf{Recommended Starting Income for $x = 15\text{ years}$ of service:}
\[ \hat{y}(15) = -2.0556 + 1.8333(15) = 25.4444\text{ ('000)} = \mathbf{25{,}444.44\text{ Tk.}} \]
\end{enumerate}
\ans{(i) $\hat{y} = -2.0556 + 1.8333x$, $\hat{x} = 1.2198 + 0.5386y$, $r = +0.9937$; \quad (ii) Recommended monthly income $= 25.44\text{ thousand}$ ($25{,}444.44\text{ Tk.}$).}
\end{solution}

%---------------------------------------------------------------- S041
\begin{question}{2014--2015 \textbullet\ Q7(a)}{MATH 241 \textbullet\ 06/02/2016 \textbullet\ S041}{20}
The following table gives the relative values of two variables:
\begin{center}\small
\begin{tabular}{@{}lccccccc@{}}
\toprule
\textbf{X} & 42 & 44 & 58 & 55 & 89 & 98 & 66 \\
\textbf{Y} & 56 & 49 & 53 & 58 & 65 & 76 & 51 \\
\bottomrule
\end{tabular}
\end{center}
Find the regression line of $Y$ on $X$ and $X$ on $Y$. Find also the correlation coefficient between $X$ and $Y$.
\end{question}
\begin{solution}
For $n = 7$ observations:
\[ \sum x_i = 452 \ (\bar{x} = 64.57), \ \sum y_i = 408 \ (\bar{y} = 58.29), \ \sum x_i^2 = 31{,}970, \ \sum y_i^2 = 24{,}312, \ \sum x_i y_i = 27{,}371. \]
Centered sums of squares and cross-products:
\[ SS_x = 2783.7143, \qquad SS_y = 531.4286, \qquad SP_{xy} = 1025.8571. \]
\begin{itemize}
\item \textbf{Regression line of $Y$ on $X$:} $b_{yx} = \dfrac{1025.8571}{2783.7143} \approx 0.3685 \implies \hat{y} = \mathbf{34.4898 + 0.3685x}$.
\item \textbf{Regression line of $X$ on $Y$:} $b_{xy} = \dfrac{1025.8571}{531.4286} \approx 1.9304 \implies \hat{x} = \mathbf{-47.9419 + 1.9304y}$.
\item \textbf{Correlation Coefficient:} $r = \sqrt{b_{yx}\,b_{xy}} = \dfrac{1025.8571}{\sqrt{(2783.7143)(531.4286)}} \approx \mathbf{+0.8434}$.
\end{itemize}
\ans{$Y$ on $X$: $\hat{y} = 34.4898 + 0.3685x$; \quad $X$ on $Y$: $\hat{x} = -47.9419 + 1.9304y$; \quad $r = +0.8434$.}
\end{solution}

%---------------------------------------------------------------- S042
\begin{question}{2013--2014 \textbullet\ Q7(a)}{MATH 241 \textbullet\ 13/12/2014 \textbullet\ S042}{20}
Calculate the two regression equations and the coefficient of correlation from the data given below:
\begin{center}\small
\begin{tabular}{@{}lcccccccc@{}}
\toprule
\textbf{Age of husband (in years) ($x$)} & 22 & 25 & 28 & 31 & 35 & 32 & 37 & 39 \\
\textbf{Age of wife (in years) ($y$)} & 19 & 18 & 22 & 29 & 31 & 23 & 30 & 33 \\
\bottomrule
\end{tabular}
\end{center}
Also estimate the most likely age of husband when wife's age is 24.
\end{question}
\begin{solution}
For $n = 8$ couples:
\[ \sum x_i = 249 \ (\bar{x} = 31.125), \ \sum y_i = 205 \ (\bar{y} = 25.625), \ \sum x_i^2 = 7993, \ \sum y_i^2 = 5489, \ \sum x_i y_i = 6601. \]
Centered sums of squares and cross-products:
\[ SS_x = 242.875, \qquad SS_y = 235.875, \qquad SP_{xy} = 220.375. \]
\begin{itemize}
\item \textbf{Regression of $Y$ on $X$:} $b_{yx} = \dfrac{220.375}{242.875} \approx 0.9074 \implies \hat{y} = \mathbf{-2.6166 + 0.9074x}$.
\item \textbf{Regression of $X$ on $Y$:} $b_{xy} = \dfrac{220.375}{235.875} \approx 0.9343 \implies \hat{x} = \mathbf{7.1839 + 0.9343y}$.
\item \textbf{Correlation Coefficient:} $r = \dfrac{220.375}{\sqrt{242.875 \times 235.875}} \approx \mathbf{+0.9207}$.
\item \textbf{Estimated Age of Husband when Wife's Age is $y = 24$:}
\[ \hat{x}(24) = 7.1839 + 0.9343(24) = \mathbf{29.61\text{ years}}. \]
\end{itemize}
\ans{$\hat{y} = -2.6166 + 0.9074x$, \ $\hat{x} = 7.1839 + 0.9343y$, \ $r = +0.9207$, \ and $\hat{x}(24) = 29.61\text{ years}$.}
\end{solution}

%---------------------------------------------------------------- S043
\begin{question}{2012--2013 \textbullet\ Q7(a)}{MATH 241 \textbullet\ 02/06/2014 \textbullet\ S043}{18}
Calculate two regression equations and regression coefficients from the data given below:
\begin{center}\small
\begin{tabular}{@{}lcccccccccc@{}}
\toprule
\textbf{Marks in Economics ($x$)} & 47 & 38 & 34 & 44 & 32 & 41 & 35 & 32 & 42 & 30 \\
\textbf{Marks in Business ($y$)} & 38 & 34 & 42 & 39 & 34 & 33 & 30 & 41 & 38 & 28 \\
\bottomrule
\end{tabular}
\end{center}
Estimate the marks in economics if a student has scored 36 in business and the mark in business if another student has scored 44 in economics. Also comment on the correlation of the above mentioned marks obtained in economics and business.
\end{question}
\begin{solution}
For $n = 10$ students:
\[ \sum x_i = 375 \ (\bar{x} = 37.5), \ \sum y_i = 357 \ (\bar{y} = 35.7), \ \sum x_i^2 = 14{,}363, \ \sum y_i^2 = 12{,}939, \ \sum x_i y_i = 13{,}461. \]
Centered sums of squares and cross-products:
\[ SS_x = 300.50, \qquad SS_y = 194.10, \qquad SP_{xy} = 73.50. \]
\begin{itemize}
\item \textbf{Regression Coefficients:} $b_{yx} = \dfrac{73.50}{300.50} \approx \mathbf{0.2446}$, \qquad $b_{xy} = \dfrac{73.50}{194.10} \approx \mathbf{0.3787}$.
\item \textbf{Regression Equations:}
\[ Y\text{ on }X:\ \hat{y} = \mathbf{26.5278 + 0.2446x}, \qquad X\text{ on }Y:\ \hat{x} = \mathbf{23.9815 + 0.3787y}. \]
\item \textbf{Predictions:}
\begin{itemize}
\item Estimated Economics mark when Business mark is $y = 36$: $\hat{x}(36) = 23.9815 + 0.3787(36) = \mathbf{37.61}$.
\item Estimated Business mark when Economics mark is $x = 44$: $\hat{y}(44) = 26.5278 + 0.2446(44) = \mathbf{37.29}$.
\end{itemize}
\item \textbf{Correlation Comment:} $r = \sqrt{b_{yx}\,b_{xy}} = \dfrac{73.50}{\sqrt{300.50 \times 194.10}} \approx \mathbf{+0.3043}$, indicating a \textbf{weak positive linear correlation} between marks in Economics and Business.
\end{itemize}
\ans{$b_{yx} = 0.2446,\ b_{xy} = 0.3787$; \ $\hat{y} = 26.53 + 0.2446x,\ \hat{x} = 23.98 + 0.3787y$; \ $\hat{x}(36) = 37.61,\ \hat{y}(44) = 37.29$; \ $r = +0.3043$ (weak positive correlation).}
\end{solution}

%---------------------------------------------------------------- S044
\begin{question}{2011--2012 \textbullet\ Q7(a)}{MATH 241 \textbullet\ 07/01/2013 \textbullet\ S044}{20}
Find the two regression lines and the coefficient of correlation from the table given below:
\begin{center}\small
\begin{tabular}{@{}lccccccccc@{}}
\toprule
\textbf{X} & 105 & 112 & 104 & 114 & 117 & 108 & 106 & 117 & 111 \\
\textbf{Y} & 58 & 56 & 47 & 46 & 51 & 64 & 66 & 50 & 55 \\
\bottomrule
\end{tabular}
\end{center}
\end{question}
\begin{solution}
For $n = 9$ observations:
\[ \sum x_i = 994 \ (\bar{x} = 110.44), \quad \sum y_i = 493 \ (\bar{y} = 54.78), \]
\[ \sum x_i^2 = 109{,}980, \quad \sum y_i^2 = 27{,}403, \quad \sum x_i y_i = 54{,}324. \]
Centered sums of squares and cross-products:
\[ SS_x = 198.2222, \qquad SS_y = 397.5556, \qquad SP_{xy} = -125.1111. \]
\begin{itemize}
\item \textbf{Regression line of $Y$ on $X$:} $b_{yx} = \dfrac{-125.1111}{198.2222} \approx -0.6312 \implies \hat{y} = \mathbf{124.4865 - 0.6312x}$.
\item \textbf{Regression line of $X$ on $Y$:} $b_{xy} = \dfrac{-125.1111}{397.5556} \approx -0.3147 \implies \hat{x} = \mathbf{127.6831 - 0.3147y}$.
\item \textbf{Coefficient of Correlation:} $r = -\sqrt{b_{yx}\,b_{xy}} = \dfrac{-125.1111}{\sqrt{(198.2222)(397.5556)}} \approx \mathbf{-0.4457}$.
\end{itemize}
\ans{$Y$ on $X$: $\hat{y} = 124.4865 - 0.6312x$; \quad $X$ on $Y$: $\hat{x} = 127.6831 - 0.3147y$; \quad $r = -0.4457$.}
\end{solution}

%===================================================================== TOPIC 18
\topic[Out of Syllabus]{18}{Out of Syllabus (CLT, Negative Binomial, Exponential Process, Sampling)}

%---------------------------------------------------------------- X001
\begin{question}{2023--2024 \textbullet\ Q6(b)}{MATH 243 \textbullet\ 21/05/2026 \textbullet\ X001}{10}
What do you mean by the central limit theorem in statistics? Write an algorithm to solve problems involving the central limit theorem. Discuss its applications in Computer Science.
\end{question}
\begin{solution}
\textbf{1. Statement of the Central Limit Theorem (CLT):}
Let $X_1, X_2, \dots, X_n$ be a random sample of size $n$ drawn from any population distribution with finite mean $\mu$ and finite non-zero variance $\sigma^2$. As the sample size $n \to \infty$ (in practice, $n \ge 30$), the sampling distribution of the standardized sample mean $\bar{X} = \frac{1}{n}\sum_{i=1}^{n} X_i$ (and sample sum $S_n = \sum_{i=1}^{n} X_i$) converges in distribution to the standard normal distribution $N(0, 1)$:
\[ Z = \frac{\bar{X} - \mu}{\sigma/\sqrt{n}} = \frac{S_n - n\mu}{\sigma\sqrt{n}} \xrightarrow{d} N(0, 1) \quad \text{as } n \to \infty. \]

\medskip
\textbf{2. Algorithm to Solve CLT Problems:}
\begin{enumerate}[Step 1:]
\item \textbf{Input:} Population mean $\mu$, standard deviation $\sigma$, sample size $n$, target query type (\texttt{MEAN} $\bar{X}$ or \texttt{SUM} $S_n$), and interval bounds $[a, b]$.
\item \textbf{Compute Sampling Parameters:}
\begin{itemize}
\item If query is on sample mean $\bar{X}$: set $\mu_T \leftarrow \mu$ and $\sigma_T \leftarrow \sigma / \sqrt{n}$ (apply finite population correction $\sqrt{(N-n)/(N-1)}$ if sampling without replacement from finite $N$ with $n/N > 0.05$).
\item If query is on sample sum $S_n$: set $\mu_T \leftarrow n\mu$ and $\sigma_T \leftarrow \sigma\sqrt{n}$.
\end{itemize}
\item \textbf{Standardize Bounds:} Compute $z_a \leftarrow (a - \mu_T)/\sigma_T$ and $z_b \leftarrow (b - \mu_T)/\sigma_T$.
\item \textbf{Evaluate Normal CDF:} Return $P(a \le T \le b) = \Phi(z_b) - \Phi(z_a)$ (or $\Phi(z_b)$ for $\le b$, $1 - \Phi(z_a)$ for $\ge a$).
\end{enumerate}

\medskip
\textbf{3. Applications in Computer Science:}
\begin{itemize}
\item \textbf{Network Traffic \& Queueing Systems:} Aggregating packet arrivals or server request latencies across thousands of independent client sessions to provision cloud bandwidth and SLA tail-latency bounds.
\item \textbf{A/B Testing \& Machine Learning Evaluation:} Constructing confidence intervals and $z$-tests for click-through rates, conversion rates, and cross-validation loss averages over large test sets.
\item \textbf{Monte Carlo Simulation \& Randomized Algorithms:} Bounding estimation error ($\mathcal{O}(1/\sqrt{n})$) in stochastic ray tracing, reinforcement learning policy evaluation, and probabilistic data structures.
\end{itemize}
\end{solution}

%---------------------------------------------------------------- X002
\begin{question}{2020--2021 \textbullet\ Q6(c)}{MATH 245 \textbullet\ 2020--2021 \textbullet\ X002}{10}
In a sociological study, about two-thirds of the 20 million persons who take Valium are women. Assuming this figure to be a valid estimate, find the probability that on a given day the fifth prescription written by a doctor for Valium is:
(i) the first prescribing Valium for a woman,
(ii) the third prescribing Valium for a woman.
\end{question}
\begin{solution}
Each prescription is an independent Bernoulli trial where ``success'' (prescribing to a woman) has probability $p = \dfrac{2}{3}$ and ``failure'' (prescribing to a man) has probability $q = 1 - p = \dfrac{1}{3}$.
The probability that the $x$-th trial is the $k$-th success follows the \textbf{Negative Binomial (Pascal) distribution}:
\[ b^*(x; k, p) = \binom{x - 1}{k - 1}p^k q^{x - k}, \qquad x = k, k+1, \dots \]
\begin{enumerate}[(i)]
\item \textbf{The 5th prescription ($x = 5$) is the 1st for a woman ($k = 1$, Geometric distribution):}
The first 4 prescriptions must be for men ($q^4$) and the 5th for a woman ($p$):
\[ P(X = 5 \text{ for } k = 1) = \binom{5 - 1}{1 - 1}\left(\frac{2}{3}\right)^1\left(\frac{1}{3}\right)^4 = 1 \times \frac{2}{3} \times \frac{1}{81} = \frac{2}{243} \approx 0.00823. \]
\item \textbf{The 5th prescription ($x = 5$) is the 3rd for a woman ($k = 3$):}
Exactly 2 of the first 4 prescriptions must be for women, and the 5th must be for a woman:
\[ P(X = 5) = \binom{4}{2}\left(\frac{2}{3}\right)^3\left(\frac{1}{3}\right)^{2} = 6 \times \frac{8}{243} = \frac{16}{81} \approx 0.19753. \]
\end{enumerate}
\ans{(i) $\dfrac{2}{243} \approx 0.00823$, \qquad (ii) $\dfrac{16}{81} \approx 0.19753$.}
\end{solution}

%---------------------------------------------------------------- X003
\begin{question}{2019--2020 \textbullet\ Q8(b)}{MATH 245 \textbullet\ 25/09/2021 \textbullet\ X003}{25}
Suppose that the time between emergency calls to a small suburban fire station follows an exponential distribution with an average rate of 1.8 calls per day.
(i) A fireman has just clocked on. What is the chance of a call in the next 15 minutes?
(ii) The fireman has nearly finished his shift: 15 minutes to go, and there has been no call during his shift so far. What is the chance of a call in the next 15 minutes?
(iii) Another firefighter works a 10-hour shift, Mondays to Thursdays. What is the probability that she has no calls in a shift?
(iv) What is the probability that she has no calls in four successive days?
(v) In 10\% of shifts, there is a call in the first $x$ hours of the shift. What is $x$, to one decimal place?
\end{question}
\begin{solution}
The arrival of calls forms a Poisson process with rate $\lambda_d = 1.8\text{ calls/day}$, which in hourly units ($1\text{ day} = 24\text{ hours}$) is:
\[ \lambda_h = \frac{1.8}{24} = 0.075\text{ calls/hour}. \]
The waiting time $T$ (in hours) until the next call has CDF $F(t) = P(T \le t) = 1 - e^{-\lambda_h t} = 1 - e^{-0.075t}$ for $t \ge 0$.
\begin{enumerate}[(i)]
\item \textbf{Chance of a call in the next 15 minutes ($t = 15\text{ min} = 0.25\text{ hours}$):}
\[ P(T \le 0.25) = 1 - e^{-0.075(0.25)} = 1 - e^{-0.01875} \approx 1 - 0.981425 = 0.018575. \]
\item \textbf{Chance of a call in the last 15 minutes given no call so far:}
By the \textbf{memoryless property} of the exponential distribution, $P(T \le s + t \mid T > s) = P(T \le t)$:
\[ P(\text{call in next 15 min} \mid \text{no call so far}) = 1 - e^{-0.01875} \approx 0.018575. \]
\item \textbf{Probability of no calls in a 10-hour shift ($t = 10\text{ hours}$):}
\[ P(T > 10) = e^{-0.075(10)} = e^{-0.75} \approx 0.472367. \]
\item \textbf{Probability of no calls in four successive days:}
\begin{itemize}
\item For \textbf{four successive 10-hour shifts} (Monday to Thursday):
\[ \big[P(\text{no call in one 10-hour shift})\big]^4 = \left(e^{-0.75}\right)^4 = e^{-3.0} \approx 0.049787. \]
\item \textit{(If interpreted as four full 24-hour calendar days: $e^{-4 \times 1.8} = e^{-7.2} \approx 0.000747$.)}
\end{itemize}
\item \textbf{Finding $x$ (in hours) such that $P(T \le x) = 0.10$:}
\[ 1 - e^{-0.075x} = 0.10 \implies e^{-0.075x} = 0.90 \implies x = \frac{-\ln(0.90)}{0.075} = \frac{0.105361}{0.075} \approx 1.4\text{ hours}. \]
\end{enumerate}
\ans{(i) $0.0186$; \quad (ii) $0.0186$ (memoryless); \quad (iii) $e^{-0.75} \approx 0.4724$; \quad (iv) $e^{-3.0} \approx 0.0498$ (4 shifts) or $e^{-7.2} \approx 0.00075$ (4 full days); \quad (v) $x = 1.4\text{ hours}$.}
\end{solution}

%---------------------------------------------------------------- X004
\begin{question}{2016--2017 \textbullet\ Q7(c)}{MATH 245 \textbullet\ 05/08/2017 \textbullet\ X004}{15}
The masses of 1500 ball bearings are normally distributed, with a mean of 22.40 g and a standard deviation of 0.048 g. If 300 random samples of size 36 are drawn from this population, determine the expected mean and standard deviation of the sampling distribution of means if the sampling is done with replacement and without replacement. How many of the random samples would have their means (i) between 22.39 and 22.41 g, (ii) greater than 22.42 g, and (iii) less than 22.38 g or more than 22.41 g? (Necessary table attached).
\end{question}
\begin{solution}
Given population size $N = 1500$, population mean $\mu = 22.40\text{ g}$, population SD $\sigma = 0.048\text{ g}$, sample size $n = 36$, and $M = 300$ samples.

\medskip
\textbf{Part 1: Mean and Standard Deviation (Standard Error) of $\bar{X}$:}
\begin{itemize}
\item \textbf{Sampling With Replacement:}
\[ \mu_{\bar{x}} = \mu = 22.40\text{ g}, \qquad \sigma_{\bar{x}} = \frac{\sigma}{\sqrt{n}} = \frac{0.048}{\sqrt{36}} = \frac{0.048}{6} = 0.0080\text{ g}. \]
\item \textbf{Sampling Without Replacement (using Finite Population Correction):}
\[ \mu_{\bar{x}} = \mu = 22.40\text{ g}, \qquad \sigma_{\bar{x}} = \frac{\sigma}{\sqrt{n}}\sqrt{\frac{N - n}{N - 1}} = 0.0080\sqrt{\frac{1464}{1499}} \approx 0.007906\text{ g}. \]
\end{itemize}

\medskip
\textbf{Part 2: Expected Number of Samples out of $M = 300$ (using $\sigma_{\bar{x}} = 0.008\text{ g}$):}
\begin{enumerate}[(i)]
\item \textbf{Sample mean between $22.39\text{ g}$ and $22.41\text{ g}$:}
\[ z_1 = \frac{22.39 - 22.40}{0.008} = -1.25, \qquad z_2 = \frac{22.41 - 22.40}{0.008} = +1.25. \]
From the normal table, $\Phi(1.25) = 0.8944$ and $\Phi(-1.25) = 0.1056$:
\[ P(22.39 < \bar{X} < 22.41) = 0.8944 - 0.1056 = 0.7888 \implies 300 \times 0.7888 = 236.64 \approx \mathbf{237}. \]
\textit{(Without replacement, $z = \pm 0.01/0.007906 \approx \pm 1.26 \implies P = 0.8962 - 0.1038 = 0.7924 \implies 300 \times 0.7924 = 237.7 \approx \mathbf{238\text{ samples}}$.)}
\item \textbf{Sample mean greater than $22.42\text{ g}$:}
\[ z = \frac{22.42 - 22.40}{0.008} = +2.50 \implies P(\bar{X} > 22.42) = 1 - \Phi(2.50) = 1 - 0.9938 = 0.0062, \]
\[ 300 \times 0.0062 = 1.86 \approx \mathbf{2\text{ samples}} \quad (\text{without replacement: } 300 \times 0.0057 = 1.71 \approx \mathbf{2}). \]
\item \textbf{Sample mean less than $22.38\text{ g}$ or more than $22.41\text{ g}$:}
For $\bar{X} < 22.38$, $z = -2.50 \implies P(\bar{X} < 22.38) = 0.0062$. For $\bar{X} > 22.41$, $z = +1.25 \implies P(\bar{X} > 22.41) = 1 - 0.8944 = 0.1056$.
Adding the two mutually exclusive tail probabilities:
\[ P(\bar{X} < 22.38\text{ or }\bar{X} > 22.41) = 0.0062 + 0.1056 = 0.1118 \implies 300(0.1118) = 33.54 \approx \mathbf{34}. \]
\textit{(Without replacement: $0.0057 + 0.1038 = 0.1095 \implies 300 \times 0.1095 = 32.85 \approx \mathbf{33\text{ samples}}$.)}
\end{enumerate}
\ans{With replacement: $\mu_{\bar{x}} = 22.40\text{ g},\ \sigma_{\bar{x}} = 0.0080\text{ g}$; without replacement: $\mu_{\bar{x}} = 22.40\text{ g},\ \sigma_{\bar{x}} = 0.0079\text{ g}$. Counts: (i) $237$ (or $238$), (ii) $2$, (iii) $34$ (or $33$) samples.}
\end{solution}

\end{document}